% Options for packages loaded elsewhere
\PassOptionsToPackage{unicode}{hyperref}
\PassOptionsToPackage{hyphens}{url}
%
\documentclass[
  11pt,
]{article}
\usepackage{amsmath,amssymb}
\usepackage{iftex}
\ifPDFTeX
  \usepackage[T1]{fontenc}
  \usepackage[utf8]{inputenc}
  \usepackage{textcomp} % provide euro and other symbols
\else % if luatex or xetex
  \usepackage{unicode-math} % this also loads fontspec
  \defaultfontfeatures{Scale=MatchLowercase}
  \defaultfontfeatures[\rmfamily]{Ligatures=TeX,Scale=1}
\fi
\usepackage{lmodern}
\ifPDFTeX\else
  % xetex/luatex font selection
\fi
% Use upquote if available, for straight quotes in verbatim environments
\IfFileExists{upquote.sty}{\usepackage{upquote}}{}
\IfFileExists{microtype.sty}{% use microtype if available
  \usepackage[]{microtype}
  \UseMicrotypeSet[protrusion]{basicmath} % disable protrusion for tt fonts
}{}
\makeatletter
\@ifundefined{KOMAClassName}{% if non-KOMA class
  \IfFileExists{parskip.sty}{%
    \usepackage{parskip}
  }{% else
    \setlength{\parindent}{0pt}
    \setlength{\parskip}{6pt plus 2pt minus 1pt}}
}{% if KOMA class
  \KOMAoptions{parskip=half}}
\makeatother
\usepackage{xcolor}
\usepackage[margin=1in]{geometry}
\usepackage{listings}
\newcommand{\passthrough}[1]{#1}
\lstset{defaultdialect=[5.3]Lua}
\lstset{defaultdialect=[x86masm]Assembler}
\setlength{\emergencystretch}{3em} % prevent overfull lines
\providecommand{\tightlist}{%
  \setlength{\itemsep}{0pt}\setlength{\parskip}{0pt}}
\setcounter{secnumdepth}{-\maxdimen} % remove section numbering
\lstset{basicstyle=\ttfamily\scriptsize,breaklines=true,columns=fullflexible,keepspaces=true}
\emergencystretch=2em
\ifLuaTeX
  \usepackage{selnolig}  % disable illegal ligatures
\fi
\usepackage{bookmark}
\IfFileExists{xurl.sty}{\usepackage{xurl}}{} % add URL line breaks if available
\urlstyle{same}
\hypersetup{
  pdftitle={Completed algebraic zero-Hessian ascent on pure analytic curves},
  pdfauthor={Unverified research draft},
  hidelinks,
  pdfcreator={LaTeX via pandoc}}

\title{Completed algebraic zero-Hessian ascent on pure analytic curves}
\author{Unverified research draft}
\date{3 October 2026}

\begin{document}
\maketitle

\subsection{1. Exact result and logical
scope}\label{exact-result-and-logical-scope}

Put \(D=d'd''\) and \(t_f=-\log|f|\). All smooth forms are
\textbf{ordinary CLD forms}, represented on ordinary open neighbourhoods
by analytic unit moments. All measures use the canonical calibration of
the actual analytic algebra, including its Artinian lengths. No space is
replaced by its reduction for integration.

\textbf{Theorem (completed algebraic curve ascent).} Let \(K\) be
complete and nontrivially valued, let \(M/K\) be an algebraic extension
with an extending absolute value, and let \(L=\widehat M\). Let \(X\) be
a good, separated, strict, boundaryless, pure one-dimensional
\(K\)-analytic space. Singular and nonreduced spaces are allowed. For
every continuous finite real function \(h\) and the scalar projection
\(\pi:X_L\to X\), \[
             D[h]=0\quad\Longleftrightarrow\quad
             D[h\circ\pi]=0.                         \tag{1}
\] The statement holds on every ordinary open subspace of \(X\). In
particular it holds for the completion of an algebraic closure, for
completed separable algebraic extensions, and in positive characteristic
for arbitrary mixtures of inseparability and separability. \(L/K\)
itself need not be algebraic. No upstairs current equality is imposed.

\textbf{Exact compact scalar test descent.} For a good separated
\(K\)-analytic \(X\), with the same fields, every
\(\beta\in A_c^{0,0}(X_L)\) is the exact smooth pullback of an
\(\alpha\in A_c^{0,0}(X_E)\) for some finite \(K\subset E\subset M\).
Here \(q_E:X_L\to X_E\) is the actual scalar projection and
\(\beta=q_E^*\alpha\) as ordinary smooth scalar sections, so also
\(D\beta=q_E^*D\alpha\). This subsidiary result needs no pure dimension
or boundarylessness. Its proof, including domains, ordinary open loci,
units and compact supports, is in Section 12. It is local descent to a
finite stage; it does not assert descent of every analytic function or
an arbitrary affinoid presentation.

The only approximation needed at the completed stage is of finitely many
analytic functions in Banach norm. Strict ultrametric inequalities turn
it into \textbf{exact} domain and unit-norm identities. No
differentiated limit or infinite group average is taken at that stage.
For finite separable ascent, the separate approximation assertion proved
in Sections 3--10 is: for invariant \(\beta\in A_c^{0,0}(X_E)\) there
are \(\alpha_j\in A_c^{0,0}(X)\) and a fixed compact \(C\subset X\) with
\[
 \operatorname{supp}\alpha_j\subset C,\qquad
 \big\|(D(\beta-\pi_E^*\alpha_j))_{\mathrm{Rad},E}\big\|_{\mathrm{TV}}
                 \longrightarrow0,                  \tag{2}
\] and every error supported in \(\pi_E^{-1}C\). Both approximation
limits, including collision-contained calibrated segments, are
reconstructed. No earlier draft theorem, review, hash or native note is
a premise.

On a pure curve the current equation is precisely \[
 B_h(\psi):=\int_X h\,(D\psi)_{\mathrm{Rad},K}=0
                  \quad(\psi\in A_c^{0,0}(X)).         \tag{3}
\] Only the differentiated smooth top form is represented by a measure;
\(D[h]\) is not presumed to be a Radon measure. Current-duality signs do
not affect vanishing.

This is a scalar transport theorem, a partial result for original P06.2.
It gives no recognition of \(h\), no two-sign CLD-psh regularization, no
universal curve-probe transport, and no higher-dimensional
separable-current ascent. P05.2 is not addressed. The original status is
\passthrough{\lstinline!partial!} and
\passthrough{\lstinline!candidate=false!}.

Sections 2--10 reconstruct the finite separable dependency. Section 11
reconstructs finite Frobenius descent and the actual finite mixed-field
decomposition. Section 12 proves the completion and support gate.
Section 13 concludes the current theorem. Section 15 includes the used
cached primary passages and proof bodies in this document itself.

\subsection{2. Exact analytic dependencies and source
checks}\label{exact-analytic-dependencies-and-source-checks}

CLD refers to Chambert-Loir--Ducros, \emph{Formes différentielles
réelles et courants sur les espaces de Berkovich}, arXiv:1204.6277v2, 25
July 2025. GR refers to Gubler--Rabinoff, \emph{Classically psh and
pluriharmonic functions on Berkovich spaces}, arXiv:2506.13548v2, 17
September 2025. Numbering is that of these cached primary snapshots,
whose statements and proof bodies, not just the former drafts'
descriptions, were inspected.

\begin{enumerate}
\def\labelenumi{\arabic{enumi}.}
\item
  \textbf{Affinoid foundations.} CLD 11.1.1--2: an affinoid spectrum is
  compact Hausdorff; its algebra is noetherian with closed ideals; the
  map to its prime spectrum is surjective. CLD 11.1.5: every point of an
  affinoid has rational affinoid neighbourhoods inside any prescribed
  ordinary open, with the point satisfying all associated strict
  inequalities. Rational systems have positive real radii and
  denominator/numerators with no common zero (Section 12 uses this
  stronger neighbourhood clause of 11.1.5). CLD 11.1.7--8: a good space
  has an affinoid neighbourhood basis; ordinary opens are analytic
  domains; analytic sections form a sheaf also for ordinary opens, and
  \(\mathcal O(\mathcal M(A))=A\). CLD 11.1.9: affinoid fibre products
  and scalar extensions use completed tensor products, and scalar
  projection is surjective. CLD 11.4.4: a finite algebra with its
  finite-module Banach topology is affinoid. CLD 11.7.7: over a
  nontrivially valued field, every point of a good strict space has a
  strict affinoid neighbourhood. Applying this to ordinary opens gives
  such a neighbourhood inside the open. Particular tensor, orbit,
  compactness, shrinking and unit claims are proved in Section 3.
\item
  \textbf{Dimension and boundary.} CLD 11.2.2: separated analytic spaces
  are topologically separated. CLD 11.2.3: strict affinoid dimension is
  Krull dimension, dimension is preserved by complete scalar extension,
  and domains preserve local analytic dimension. If the minimal-prime
  components of \(\mathcal M(A)\) are \(X_i\), then
  \(\dim_xX=\sup_{x\in X_i}\dim X_i\). Pure dimension means constant
  local dimension. GR 2.11(5), equivalently (5'), states for a morphism
  \(Y\to Z\) and any analytic extension \(L/K\) that
  \(\partial(Y_L/Z_L)\subseteq\pi^{-1}(\partial(Y/Z))\). Hence a
  boundaryless space remains boundaryless. GR 2.11(3),(4) gives
  boundarylessness for an ordinary open in a boundaryless space: its
  relative interior is its ordinary topological interior and interiors
  compose. Section 3 proves purity component by component; equality of
  total dimensions alone is not the proof.
\item
  \textbf{Ordinary forms.} CLD Corollary 15.1.2 states that a function
  is smooth near \(x\) if it has a representation
  \(\varphi(t_{f_1},\ldots,t_{f_m})\) on a neighbourhood of \(x\), where
  the \(f_i\) are analytic units and \(\varphi\) is smooth on a
  Euclidean neighbourhood of the coordinate value. CLD Proposition
  15.1.1(b) gives ordinary open moment descriptions for forms in every
  bidegree. CLD 6.2.1--2 defines ordinary forms and says they form an
  ordinary sheaf stable under wedge, \(J,d',d''\). Moment pullback
  commutes with the exterior calculus (CLD 3.3.12). In a moment a
  logarithm is an affine coordinate, so \(Dt_f=0\) as a \textbf{form}.
  CLD 3.4.4: for a G-tropical chart \(f\) and a coefficient form \(a\),
  its pullback is zero in the G-form sheaf exactly when \(a\) restricts
  to zero on \(f(V)\) for every compact domain \(V\). This identifies
  the ordinary form sheaves when the analytic units, domains and moment
  images are the same; it is not descent through scalar projection.
\item
  \textbf{Compact cutoffs only.} CLD 15.2.2: a good topologically
  separated space is locally holomorphically separated. CLD Corollary
  15.2.5: on a locally holomorphically separated analytic space, for any
  compact \(F\) and ordinary open neighbourhood \(U\) of \(F\), there is
  an ordinary smooth function equal to one near \(F\) with compact
  support contained in \(U\). Its proof uses local differential
  separation (15.2.4) and the smooth bump result 6.2.4. Section 4
  constructs the only partition statement needed, for compact \(F\) and
  finite covers. The overbroad fixed-index countable-at-infinity reading
  of CLD 15.2.9 is not used.
\item
  \textbf{Canonical top measures.} CLD Theorem 15.7.1: a pure
  \(n\)-dimensional topologically separated analytic space with its
  canonical calibration and analytic boundary is a tropical space of
  level \(n\). CLD 15.7.2: the induced structure on a domain agrees with
  that domain's analytic structure. CLD 5.4.4: every compact tropical
  chart is pleasant, and its tropical polyhedron is the top-dimensional
  part of its compact moment image. CLD 14.3.2--5, 14.5.1--2: canonical
  calibrations are standard and positive, and on a top moment cell \(Q\)
  the calibration is \(d_Q\tau_{Q,*}|e_1\wedge\cdots\wedge e_n|\), where
  \(\tau_Q\) is the affine inverse of an injective torus projection and
  \(d_Q\) is the sum of the actual finite-flat degrees of the skeletal
  cells above \(Q\). Those degrees retain fibre-algebra dimensions,
  including lengths. CLD 5.8.3--5 constructs a top form's measure by
  contraction on its supporting skeleton. CLD Proposition 5.10.1(c,d):
  for a pleasant chart \(f\) and a Borel \((n,n)\) coefficient form
  \(a\), the induced measure is locally finite exactly when \(a\) is
  locally integrable on the calibrated tropical polyhedron, and then
  \(f_*((f^*a)_{\mathrm{Rad}})=\langle a,\mu_{\Pi_f}\rangle_{\mathrm{Rad}}\);
  the corresponding integrals agree. In particular a smooth top form is
  finite on a compact moment domain. Lower-dimensional cells have no top
  measure. Compact moment images are finite polyhedra by the
  polyhedrality assertion of CLD Theorem 12.2.4. The dimension bound is
  specifically CLD 12.2.10(1)--(2):
  \(\dim_{\mathrm{trop},x}f\le d_K(x)\le\dim_xX\) and, for compact
  \(X\), \(\dim f_{\mathrm{trop}}(X)=\sup_x\dim_{\mathrm{trop},x}f\).
  Thus compact curve images have dimension at most one. We do not use a
  claim that a moment image always has dimension equal to its source;
  constant moments rule out that reading. CLD 7.3.5 defines integration
  currents for locally integrable multipliers; a continuous finite
  multiplier is locally integrable. Section 8 proves the carrier
  total-variation estimate needed here, before any possibly cancelling
  signed pushforward.
\item
  \textbf{Scalar projection has no field-degree factor.} CLD Proposition
  14.3.6: on a level-\(n\) analytic space, if \(\Sigma\) has canonical
  calibration \(\mu\) and \(\Sigma_L\) its scalar extension has
  canonical calibration \(\mu_L\), then \(\pi_*\mu_L=\mu\). CLD
  Proposition 15.3.1: for a tropically Borel \((n,n)\) G-form
  \(\omega\), \(\pi_*\omega_{L,\mathrm{Bor}}=\omega_{\mathrm{Bor}}\);
  for a locally integrable such form, \[
             \pi_*\omega_{L,\mathrm{Rad}}
                       =\omega_{\mathrm{Rad}}.          \tag{4}
  \] Here is the primary proof mechanism, retaining its lengths. On a
  convenient skeletal cell choose a torus projection injective on the
  cell. Its finite-flat degree downstairs equals the sum of its degrees
  at all upstairs skeletal lifts. This is CLD Lemma 13.3.3(a): degrees
  at \(y\) equal the sum of degrees over \(y\) after complete base
  change and choice of a target lift. The lemma reduces to a fibre; a
  zero-dimensional local fibre algebra is Artinian, and its tensor
  decomposes into its local Artinian factors. The sum of their
  dimensions is the original dimension after scalar extension.
  Positive-dimensional fibres give infinite degrees on both sides. Thus
  the cell volume vectors sum to the downstairs vector, proving 14.3.6.
  Contracting the form and taking Radon measures proves 15.3.1. We apply
  (4) only to smooth top forms, then pair with \(h\); this never assumes
  pullback of a continuous current.
\item
  \textbf{Same-carrier field restriction has a degree factor.} CLD
  Remark 14.4.3: for a finite extension \(F/K\), an \(F\)-analytic space
  \(Z\) of level \(n\), and its restricted \(K\) structure \(Z_0\),
  dimensions and skeletal subsets agree and \(\mu_{Z_0}=[F:K]\mu_Z\).
  The proof uses the finite-flat degree \([F:K]\) of the projection of
  the split torus on this restricted carrier. CLD 15.4.1 and Proposition
  15.4.3: for a morphism \(f:Y\to Z\) of level-\(n\) spaces finite off a
  Zariski closed set of dimension \(<n\), a locally integrable \((p,n)\)
  G-form \(\omega\), \(p\in\{n-1,n\}\), has locally integrable pullback
  and
  \(f_*((f^*\omega)_{\mathrm{Rad}})=\mathrm{dc}(f)\omega_{\mathrm{Rad}}\).
  Its proof first proves the calibrated volume-vector identity on the
  inverse skeletal set (14.4.2), then contracts and integrates. We use
  only an actual isomorphism on \(Z_0\), whose calibrated degree is one:
  both local Artinian rings are identical. Cancelling the same-carrier
  factor \([F:K]\) proves invariance over \(F\).
\end{enumerate}

All these source statements apply to the actual, potentially singular
and nonreduced structure. They require no normalization and no
quasi-smoothness. No weakly smooth GJR descent, positive-current
recognition, or catalogue theorem is used as an ordinary form theorem.

\subsubsection{2.1. Direct proofs of the neighbourhood and
compact-cutoff facts used
here}\label{direct-proofs-of-the-neighbourhood-and-compact-cutoff-facts-used-here}

The finite proof uses the primary foundations above. For the two
topological localization facts decisive at completion, the following
direct constructions supply their proof bodies as well.

On \(\mathcal M(A)\) the ordinary topology is initial for all
evaluations \(|a|\) (CLD 11.1.1). A prescribed neighbourhood of a point
therefore contains a finite intersection of conditions
\(|a_j|>\ell_j>0\) and \(|b_k|<u_k\), where thresholds can be chosen
strictly between the point's value and the excluded endpoint. Missing
lower or upper bounds impose no condition. Put \(g=\prod_j a_j\), with
\(g=1\) when there are no lower bounds. Choose \(C>0\) so that
\(C^{-1}<\prod_j\ell_j\). The rational conjunction \[
 |1|\le C|g|,\qquad
 |\prod_{j'\ne j}a_{j'}|\le\ell_j^{-1}|g|\ (\forall j),\qquad
 |b_kg|\le u_k|g|\ (\forall k)                         \tag{R}
\] has no common zero among its denominator and numerators, because one
numerator is \(1\). It is a rational analytic domain by the universal
rational-domain construction of CLD 11.1.5. Its validity locus lies in
the originally prescribed neighbourhood if the thresholds were chosen
inside that neighbourhood's intervals: its first inequality excludes
\(g=0\), so division of norms gives all chosen lower and upper bounds.
The original point satisfies every strict inequality, including the
first. Non-strict thresholds in (R) were chosen inside the original open
intervals, so validity of the weak system suffices for containment. This
proves the strong rational-neighbourhood clause needed in Section 12
using the initial topology and the rational-domain construction. Real
positive thresholds suffice even when the value group is discrete.

Here is also a direct compact-cutoff construction in every good
Hausdorff analytic space occurring in the proof. For a point \(x\) in an
ordinary open \(U\), take an affinoid neighbourhood \(V_0\) and within
its ordinary interior a rational neighbourhood \(R\) of \(x\) contained
in \(U\), with strict inequalities \(|f_i(x)|<r_i|g(x)|\). Its
denominator satisfies \(g(x)\ne0\). Choose \(b_i\) with
\(|f_i(x)|/|g(x)|<b_i<r_i\) and choose \(0<a<|g(x)|\). On the ordinary
nonzero locus of \(g\) in \(V_0\), a smooth function of \(\log|g|\)
equals one near \(x\) and vanishes where \(|g|\le a\). For each \(i\),
choose a smooth real function of \(\log|f_i/g|\) that is one when
\(|f_i/g|\le b_i\) and zero when \(|f_i/g|\ge r_i'\) for some
\(b_i<r_i'<r_i\). Extend that factor as the constant one at zeros of
\(f_i\): near such a zero its ratio is below \(b_i\), so the extension
is locally constant and ordinary smooth. On its nonzero locus the ratio
is an analytic unit, giving the required ordinary moment description.
Multiply the factors with the denominator cutoff. At points where
\(g=0\), extend by zero; the denominator cutoff already vanishes on a
neighbourhood of each such point.

The support of this product lies in the compact set in \(V_0\) defined
by \(|g|\ge a\) and \(|f_i|\le r_i'|g|\) for all \(i\). That compact set
lies in the strict locus of \(R\), hence in
\(U\cap\operatorname{int}_X(V_0)\). Extend the product by zero outside
\(V_0\): every point outside the compact support has a zero
neighbourhood, so the extension is ordinary smooth. It equals one near
\(x\). Taking squares gives nonnegative cutoffs. For compact
\(F\subset U\), finitely many positive loci of these point cutoffs cover
\(F\). Their sum \(s\) is positive near \(F\). Compose with a real
smooth function that vanishes near zero and equals one on a
neighbourhood of the compact positive set \(s(F)\). This gives a compact
cutoff equal to one near \(F\) and supported in \(U\). This reconstructs
the compact CLD 15.2.5 conclusion in our actual good Hausdorff scope.
Section 4 then proves its finite-cover partition specialization. No
countable-at-infinity fixed-index compact-support assertion is used.

In particular all cutoffs are obtained from ordinary analytic-unit
moments, flat constant extensions near analytic zeros and compact
extension by zero. No weakly smooth tropicalization is substituted.

\subsection{3. Finite fields, fibres, analytic structures, and
purity}\label{finite-fields-fibres-analytic-structures-and-purity}

We first treat a finite Galois extension \(E/K\), with group \(G\) and
\(d=[E:K]=|G|\). The elementary finite field facts used below can be
justified as follows. A finite separable extension has as many distinct
embeddings in an algebraic closure as its degree: extend embeddings
successively by the distinct roots of each separable minimal polynomial.
The field \(K\) is infinite. For two separable generators \(u,v\), all
but finitely many \(c\in K\) make the values of \(u+cv\) distinct for
every pair of distinct embeddings; if their \(v\) values agree their
\(u\) values already differ. Its minimal polynomial then has at least as
many roots as the extension degree, so \(u+cv\) generates the extension.
Induction gives a primitive element. The splitting field of its
separable minimal polynomial is finite separable. Every embedding of
this splitting field permutes the roots and is onto, so it is Galois and
its automorphism group has its degree as order. Its fixed field is
\(K\): the embeddings of any intermediate separable simple field extend
successively to the splitting field, so a fixed element has a minimal
polynomial of degree one. This also provides the normal closure used in
Section 10.

Finite-dimensional norms over a complete non-Archimedean field are
equivalent to the coefficient maximum norm in a fixed basis. Here is a
proof. Induct on dimension. The span of all but the last vector is
complete by induction, hence closed. The last vector has positive
distance from this span, bounding its coefficient in any combination.
The ultrametric inequality and induction then bound all other
coefficients. The reverse bound is the triangle inequality. Consequently
finite-dimensional normed spaces are complete and their coordinate maps
bounded. Existence of an absolute value on a finite field extension can
also be seen by giving the finite algebra its finite-module Banach
topology, using CLD 11.4.4, and taking a point of its nonempty affinoid
spectrum. Its kernel is zero since the algebra is a field. Uniqueness
follows from norm equivalence: compare two extending norms on \(z^m\)
and take \(m\)th roots. Hence all finite field automorphisms preserve
the absolute value, and every finite field here is complete.

For a \(K\)-affinoid algebra \(A\), tensor the bounded coordinate maps
of \(E\) with \(A\). The tensor norm upper bound and these coordinate
maps identify its topology with a finite sum of copies of \(A\). Already
complete, this proves \[
             B=A\widehat\otimes_K E=A\otimes_K E.        \tag{5}
\] This is the actual algebra, including nilpotents. Taking the kernel
of \(E\to\bigoplus_{g\in G}E\), \(z\mapsto(gz-z)_g\), and tensoring over
the field gives \(B^G=A\). Invariant units descend as units since their
inverses are invariant too. There is no field averaging or division by
\(d\). An affinoid element nonzero at every analytic point is a unit:
otherwise a maximal ideal containing it and spectral surjectivity give a
zero.

The fibre over \(x\) is \(\mathcal M(\mathcal H(x)\otimes_K E)\). A
primitive separable polynomial factors without repeated factors over
\(\mathcal H(x)\), so the tensor is a finite product of fields. Each has
one analytic point by the existence and uniqueness argument above. The
\(G\) action on these finitely many points is transitive: two orbits
would supply a nontrivial invariant idempotent, but the invariant
algebra is \(\mathcal H(x)\). Stabilizers can be nontrivial; no local
covering-space description is claimed or used.

Inverse images of affinoid ordinary interiors give goodness upstairs.
Strict presentations base change to strict presentations. Base change of
the closed diagonal gives separatedness. Boundarylessness follows from
Section 2(2). In particular all carriers are Hausdorff. The inverse
image of any compact \(F\subset X\) is compact: cover \(F\) by finitely
many interiors of affinoid neighbourhoods \(V_i\); its inverse image is
closed in the finite union of the compact \(V_{i,E}\).

\textbf{Full-fibre shrinking.} In a downstairs strict affinoid
neighbourhood \(V_0\) of \(x\), let \(N\subset V_{0,E}\) be a relative
ordinary open containing the whole fibre. The compact complement has
closed image excluding \(x\). Choose a smaller strict affinoid
neighbourhood \(V\) of \(x\) in the complement of that image and inside
\(\operatorname{int}_X(V_0)\). Then \(V_E\subset N\). Thus finitely many
unit or strict norm conditions holding on the whole fibre persist on one
neighbourhood. This includes conditions on elements of zero seminorm at
\(x\) and never replaces \(A\) by its reduction.

For completeness, purity is checked at minimal components. In a strict
affinoid \(\mathcal M(A)\) in \(X\), each minimal-prime component has
dimension one: choose a point with kernel that prime by spectral
surjectivity and apply the local dimension formula. Let \(J\) be the
nilradical. It is nilpotent since \(A\) is noetherian. The injection
\(A/J\hookrightarrow\prod_P\operatorname{Frac}(A/P)\) stays injective
after tensoring with a separable field. Each field tensor is a product
of fields, hence \((A/J)\otimes E\) is reduced and the nilradical of
\(B\) is exactly \(J\otimes E\). The finitely many kernels of its maps
to those field factors have zero intersection. Their product is zero, so
a minimal prime contains one of these kernels, and by minimality equals
it. Its contraction is the corresponding minimal \(P\) of \(A\). The
finite integral injection \(A/P\hookrightarrow B/Q\) preserves Krull
dimension. One elementary justification is lying over and going up from
the determinant argument for a finite module over a local ring:
\(\mathfrak pD_{\mathfrak p}\ne D_{\mathfrak p}\), since otherwise the
determinant of \(I-M\), \(M\) with entries in \(\mathfrak p\), is a unit
annihilating that nonzero module. Choose a maximal ideal above it and
then repeat after quotienting by a chosen prime. For incomparability
quotient by the smaller prime and invert the nonzero elements of the
base domain. The resulting finite integral domain over a field is a
field, since a least-degree monic equation for a nonzero element
expresses its inverse. A strictly larger prime of the same contraction
could not survive, nor could a nonzero element of a domain vanish under
that localization. Going up lifts chains and incomparability makes
contractions strict. The dimensions agree. Every minimal component
upstairs therefore has dimension one, proving purity. The same proof
works for any finite separable stage. Reduction here is a dimension
check only; all measures below use \(B\) itself.

Finally consider the same carrier \(Z=X_E\) with its restricted
\(K\)-analytic structure \(Z_0\). A finite \(E/K\) is \(K\)-affinoid; an
\(E\)-affinoid presentation is therefore \(K\)-affinoid under finite
field restriction. The finite-module topology is unchanged. Rational
domains are defined by the same elements and inequalities and have the
same analytic rings. They give a common local basis and the same unit
germs and compact moment images. The moment nullity criterion of Section
2(3) identifies the ordinary form sheaves, their exterior calculus and
supports. Each \(g\in G\) is a bounded \(K\)-algebra automorphism of
(5), gluing to a \(K\)-analytic automorphism of \(Z_0\). Its isomorphism
pushforward formula, followed by cancellation of the same-carrier factor
\(d\), gives for a smooth compact top form \(\omega\) \[
             g_*((g^*\omega)_{\mathrm{Rad},E})
                            =\omega_{\mathrm{Rad},E}.    \tag{6}
\] This is different from (4). Both are valid with actual generic
lengths.

\subsection{4. Compact partitions and full-fibre scalar
descriptions}\label{compact-partitions-and-full-fibre-scalar-descriptions}

Let compact \(F\) have an ordinary open cover. At every point of \(F\),
apply Section 2(4)'s bump theorem to the singleton in a chosen cover
member. Squaring the bump makes it nonnegative, still equal to one near
that point and compactly supported there. Select finitely many whose
positive loci cover \(F\), and group them by cover index as \(\psi_i\).
Then \(s=\sum_i\psi_i>0\) near \(F\). Take a bump \(\chi=1\) near \(F\)
compactly supported in \(\{s>0\}\) and put \[
                \theta_i=\chi\psi_i/s                  \tag{7}
\] there, extending by zero. Products, smooth composition, and
reciprocal on \(s>0\) are ordinary smooth operations in moment
coordinates. The compact support of \(\chi\) lies inside \(\{s>0\}\), so
extension by zero is smooth also at its complement. Each support is
compact and contained in its cover member. The sum is one near \(F\).
This proves the needed finite compact partition without any noncompact
fixed-cover-index assertion.

Let \(\beta\) be an invariant ordinary smooth scalar function near a
whole fibre over \(x\). Its fibre has finitely many points. Choose
pairwise disjoint ordinary neighbourhoods \(N_i\) of them with local
representations \(\beta=\varphi_i(t_{f_{i,1}},\ldots)\), using Corollary
15.1.2. Shrink \(N_i\) so the coordinate values are in a relatively
compact Euclidean region on which \(\varphi_i\) is smooth. Full-fibre
shrinking gives a strict affinoid \(V\) with \(V_E\subset\bigcup_iN_i\).
The pieces \(C_i=V_E\cap N_i\) are disjoint ordinary opens and closed in
\(V_E\). They are compact. Analytic gluing gives idempotents \(e_i\)
which are one on \(C_i\) and zero on the other pieces, and gives a
global unit \(\widetilde f_{i,k}\) equal to \(f_{i,k}\) on \(C_i\) and
one elsewhere. Its inverse glues in the same way. These are elements of
the actual affinoid algebra \(B\), not formal topological selectors.

Choose \(a\in K^\times\) with \(|a|\ne1\). The analytic unit
\(s_i=e_i+a(1-e_i)\) has \(t_{s_i}=0\) on \(C_i\) and
\(t_{s_i}=-\log|a|\) elsewhere. Choose a real smooth \(q\) equal to one
near zero and zero near \(-\log|a|\). Extend \(\varphi_i\) by a
Euclidean cutoff equal to one near its compact coordinate image. Thus \[
 \beta=\sum_i q(t_{s_i})\varphi_i(t_{\widetilde f_{i,1}},\ldots)
                          =\Phi(T)\quad\text{on }V_E,   \tag{8}
\] for one global unit tuple \(T\) and a Euclidean smooth \(\Phi\). All
group translates of these units exist on \(V_E\); adjoining them gives
one fixed tuple. The real average \(|G|^{-1}\sum_g\Phi(T\circ g)\)
equals \(\beta\) there. This construction works for every orbit and
stabilizer, including singleton fibres. It does not assume exact
invariant ordinary scalar descent.

\subsection{5. Polarization removes the orbit-correlation
difficulty}\label{polarization-removes-the-orbit-correlation-difficulty}

For a smooth Euclidean function on a neighbourhood of a compact box,
there are real polynomials converging through order two on that box.
Here is a proof of precisely the version used. Extend it by a compact
Euclidean cutoff, convolve with a Gaussian, and differentiate under the
integral. Uniform continuity of the derivatives through order two and
the Gaussian tail give uniform derivative convergence on the box. For
fixed convolution width, the Gaussian and its first two derivatives are
uniform limits of their Taylor polynomials on the compact difference of
the box and the original coefficient support. Integrate these
polynomials to get polynomial approximants through order two, then
choose widths and truncation degrees diagonally.

For \(k\ge1\) and real symbols \(z_1,\ldots,z_k\), including repeated
symbols, inclusion-exclusion and the multinomial expansion give \[
 z_1\cdots z_k=\frac1{k!}\sum_{J\subseteq\{1,\ldots,k\}}
            (-1)^{k-|J|}\left(\sum_{j\in J}z_j\right)^k. \tag{9}
\] Indeed each term missing one symbol cancels, and the only surviving
degree-\(k\) term containing every symbol is their product with
coefficient \(k!\). For unit logarithms, a sum is the logarithm of the
corresponding product unit. Consequently the real group average of
\textbf{every} scalar polynomial in an arbitrary tuple of unit
logarithms is a finite real linear combination of constants and terms \[
        T_{f,k}=\frac1d\sum_{g\in G}t_{gf}^{,k},          \tag{10}
\] with actual product units \(f\). This recovers mixed orbit
correlations; no claim that separate symmetric sums for the original
individual units suffice is used. The number of product units may grow
with polynomial degree. Section 9 handles this explicitly; no common
neighbourhood for infinitely many product units is assumed.

\subsection{6. Actual root valuations and all coefficient-zero
cases}\label{actual-root-valuations-and-all-coefficient-zero-cases}

For a unit \(f\in B^\times\) set \[
 P_f(Z)=\prod_{g\in G}(Z-gf)
             =Z^d+\sum_{i=0}^{d-1}c_iZ^i.               \tag{11}
\] Its coefficients lie in \(B^G=A\); \(c_0\) is a unit. At an analytic
point \(y\), the roots \(gf(y)\) are actual nonzero elements of
\(\mathcal H(y)\) with multiplicities indexed by \textbf{all} group
elements. This includes repeated roots and arbitrary stabilizers.
Neither a separable orbit polynomial nor characteristic-zero
coefficients are required.

We prove the root-valuation rule rather than importing Newton-polygon
recognition. For a polynomial \(\prod_{j=1}^d(Z-r_j)\) over any valued
field with nonzero roots, sort \(u_j=-\log|r_j|\) in nondecreasing
order. The coefficient of \(Z^{d-k}\) is a sum of products of \(k\)
roots. Every summand has valuation at least \(u_1+\cdots+u_k\), so its
coefficient has valuation at least that number, allowing infinity. If
\(u_k<u_{k+1}\), the subset of the \(k\) roots of smallest valuation is
unique and its product has strictly greater norm than every other
summand. Ultrametricity makes the coefficient valuation exactly that
sum. Equality also holds at \(k=0,d\). These equality points are
precisely the ends of constant-valuation blocks. The convex polygon
joining them lies below all the coefficient points and has its
successive slopes equal to the root valuations in reversed sign when
indexed by coefficient exponent. Thus the lower convex hull of \[
          (i,v(c_i))\ (0\le i\le d),\qquad v(c_d)=0      \tag{12}
\] has each segment \(i<j\) contributing \(j-i\) roots of valuation
\((v(c_i)-v(c_j))/(j-i)\). Repeated roots, ties, and every cancellation
are covered by the strict-gap argument. No division occurs in the valued
field.

Near a fixed downstairs \(x\), choose \(M\) so all \(|t_{gf}|<M\) on its
finite fibre and retain this inequality by full-fibre shrinking. Then
\(|t_{c_0}|<dM\). If \(c_i(x)=0\) for \(0<i<d\), shrink further until \[
                      |c_i|<\exp(-dM-1).                \tag{13}
\] Its valuation lies strictly above the chord from \((0,t_{c_0})\) to
\((d,0)\), since that chord's height is at most \(dM\). Hence it cannot
belong to the lower hull and may be omitted. If \(c_i(x)\ne0\), shrink
to make it a downstairs unit: choose a strict affinoid neighbourhood
inside its nonzero locus, then use the unit criterion of Section 3. Keep
\(c_0,c_d=1\) in all cases. The retained index set
\(I\subseteq\{0,\ldots,d\}\) therefore includes both ends. Its other
coefficients are actual units; all omitted coefficients obey (13)
throughout this neighbourhood. Nilpotent coefficients are included in
the zero-seminorm case. No coefficient with value zero is logarithmed,
and small coefficients are not simply assumed harmless.

\subsection{\texorpdfstring{7. A \(C^1\) piecewise polynomial and its
tangential
mollification}{7. A C\^{}1 piecewise polynomial and its tangential mollification}}\label{a-c1-piecewise-polynomial-and-its-tangential-mollification}

For arbitrary real heights \(z_i\), \(i\in I\), with \(z_d=0\), form the
lower hull of the finitely many points \((i,z_i)\). If consecutive
vertices are \(i<j\), put \(m=j-i\) and \(u=(z_i-z_j)/m\). Define \[
               H_{I,k}(z)=\sum_{\text{hull segments }i<j}
                                 m\,u^k.                \tag{14}
\] Collinear vertices may be retained or omitted: (14) is unchanged.
There are finitely many possible hull chains. Each is selected by linear
weak inequalities in the heights, so (14) is a polynomial on each region
of a finite polyhedral subdivision of the whole height space. Equations
(11)--(13) prove at \textbf{every} point upstairs that \[
               T_{f,k}=d^{-1}H_{I,k}(t_{c_i}).            \tag{15}
\]

We need more than continuity of this formula. For a general smooth
one-variable \(\phi\), the segment term \(m\phi((z_i-z_j)/m)\) has
gradient \(\phi'(u)(e_i-e_j)\) and Hessian \[
            \frac{\phi''(u)}m(e_i-e_j)\otimes(e_i-e_j).  \tag{16}
\] If a vertex is inserted or deleted on a collinear segment, all the
affected \(u\) equal one common value. The gradients then telescope:
\(\sum\phi'(u)(e_i-e_j)\) is exactly the gradient of the single merged
segment. This remains true if several vertices are collinear. At a
height configuration on a chamber boundary, every limiting hull chain
refines the unique lower hull by such collinear vertices. Values and
gradients therefore agree for every limiting chamber. A finite piecewise
smooth function with these matching values and gradients is \(C^1\):
integrate the matching, continuous candidate gradient on short line
segments subdivided by the finitely many chamber faces; segments
contained in faces follow by a small translation and continuity. The
usual differentiability remainder follows from continuity of that
gradient. Thus (14) is \(C^1\).

On any compact neighbourhood, slopes \(u\) are bounded. Each Hessian
(16) is bounded by \(2\sup|\phi''|/m\) in operator norm; there are at
most \(d\) segments. For \(\phi(u)=u^k\), there is a finite bound for
all chamber Hessians on that neighbourhood. The gradient is locally
Lipschitz by integrating these Hessians along line segments, using its
continuity at faces as above. The weak Hessian is the bounded piecewise
Hessian, with no singular face measure: integrate by parts on each
polyhedral region; the boundary contributions from the continuous
gradient cancel in pairs. Lower-dimensional intersections give no
additional boundary term. This proves all assertions used in
mollification; bounded third or fourth derivatives are unnecessary.

Choose a nonnegative even \(C^\infty\) kernel \(\rho\) supported in the
unit ball with integral one. Let \(H_\varepsilon=H*\rho_\varepsilon\).
The local gradient Lipschitz bound yields uniformly on fixed compacts
\(H_\varepsilon\to H\) and \(\nabla H_\varepsilon\to\nabla H\); in fact
the errors are \(O(\varepsilon^2)\) and \(O(\varepsilon)\). The first
bound follows from the first-order remainder and the kernel's zero first
moment. Its Hessian is the convolution of the weak Hessian and is
uniformly bounded on those compacts.

\textbf{Tangential second derivative lemma.} Let an open segment \(L\)
lie in the relative interior of a face \(F\) of the subdivision of
\(H\), with tangent vector \(v\) in the affine direction of \(F\). At
every point \(z\in L\), \[
        \nabla^2H_\varepsilon(z)(v,v)
                  \longrightarrow (H|_L)''(z;v,v).       \tag{17}
\] To prove this, only chambers incident to \(F\) meet a sufficiently
small ball about \(z\). Let their polynomial formulas be \(Q_a\). Each
agrees with \(H\) on a relatively open neighbourhood in \(F\), so
\(\nabla^2Q_a(z)(v,v)\) is the same tangential derivative of that
restriction. For points \(w\) tending to \(z\) through any such chamber,
the corresponding tangential Hessian tends to this common value.
Convolving the bounded, piecewise Hessian proves (17). If the line is
mapped to one point in height space, the tangent is zero and the
statement is immediate. This proof deals with segments \textbf{entirely
contained in collision walls}; they are not dismissed as measure zero. A
compact moment segment maps affinely to height space and can be
subdivided at finitely many face preimages, so the conclusion holds
almost everywhere on every such segment, including those of deficient
height-map rank.

\subsection{8. Calibrated differentiated approximation of a fixed
trace}\label{calibrated-differentiated-approximation-of-a-fixed-trace}

Fix \(f,k\) and one neighbourhood on which (13)--(15) hold. The
functions \[
                b_\varepsilon=d^{-1}H_\varepsilon(t_{c_i})
                                                               \tag{18}
\] are ordinary smooth downstairs, using actual unit moments. On a
compact upstairs domain add to one moment tuple the units \(gf\), the
pullbacks of all retained coefficients, and any fixed cutoff's units.
Its image is a finite polyhedral set. Refine its one-dimensional cells
by the height-fan preimages and by a common skeletal calibration
subdivision. On each segment, (15) is an identity of functions of affine
coordinates. Its right side restricts to a polynomial on the segment,
whose first and second tangent derivatives are therefore those of the
ordinary upstairs polynomial \(T_{f,k}\).

The errors \(e_\varepsilon=T_{f,k}-\pi^*b_\varepsilon\) and their first
derivatives converge uniformly tangentially. Their second tangential
derivatives are uniformly bounded and converge almost everywhere on
every segment by (17). No differentiation of (15) as a nonsmooth
ordinary section is performed; the derivatives are compared on actual
moment cell tangents.

Here is why this is a total-variation statement on the actual carrier.
For the enlarged tuple \(t_1,\ldots,t_r\) put
\(\zeta=\sum_a d't_a\wedge d''t_a\), with the positive top-order
convention. On a segment with tangent \(v\), a coefficient matrix
\(A=(A_{ab})\) contracts to \(\sum A_{ab}v_av_b\) times its positive
calibration; \(\zeta\) contracts to \(\sum v_a^2\) times the same
calibration. Hence \(|A_{ab}|\le M\) gives \[
       |\eta_{\mathrm{Rad}}|\le rM\zeta_{\mathrm{Rad}},    \tag{19}
\] by \(\big(\sum|v_a|\big)^2\le r\sum v_a^2\). Equivalently use bounded
tangential coefficients in unit tangent frames on the finitely many
segments. The estimate is made on the supporting skeletal cells before
pushforward. Section 2(5) gives finite total mass of \(\zeta\) on the
compact domain and a constant finite positive density on each moment
segment. Thus bounded tangential error converging almost everywhere has
carrier total variation tending to zero by bounded convergence on those
finite calibrated segments. Zero-dimensional faces carry no top measure.
Multiplicities and generic lengths multiply both sides of (19); none is
removed or replaced by a point count.

For any \textbf{fixed} ordinary smooth cutoff \(\lambda\) supported in
the neighbourhood, the product rule for \(D(\lambda e_\varepsilon)\)
uses the scalar error, its two first derivatives, and its mixed second
derivative, together with fixed bounded derivatives of \(\lambda\). The
preceding convergence and (19) show \[
       \big\|(D((\lambda\circ\pi)e_\varepsilon))_{
                          \mathrm{Rad},E}\big\|_{\mathrm{TV}}
                                 \longrightarrow0.       \tag{20}
\] Cover a fixed compact cutoff support by finitely many compact moment
neighbourhoods to use this argument; their enlarged tuples and
calibration refinements are fixed during this limit. Signed fibre
cancellation is not used to infer (20).

For a fixed finite combination of traces (10), Sections 6--8 apply near
every downstairs point, with finitely many units, omitted coefficients
and strict inequalities at that point. If a fixed downstairs cutoff
\(\chi\) has compact support in the original unit domain, cover its
support by finitely many of these local neighbourhood interiors, and
choose an inner compact partition \(\theta_a\) by (7). Let
\(b_{a,\varepsilon}\) be the local combination of (18). Put \[
        a_\varepsilon=
              \sum_a\widetilde{\chi\theta_a b_{a,\varepsilon}},     \tag{21}
\] where tilde means extension by zero. This is ordinary smooth and
supported in \(\operatorname{supp}\chi\). A product supported compactly
inside its ordinary open is identically zero on a neighbourhood of each
point outside that open, so this extension is smooth and has no boundary
derivative term. The sum-one identity near \(\operatorname{supp}\chi\)
and (20), for this finite fixed cover, prove \[
 \big\|(D((\chi\circ\pi)T-\pi^*a_\varepsilon))_{
                      \mathrm{Rad},E}\big\|_{\mathrm{TV}}\to0       \tag{22}
\] for the given invariant scalar polynomial \(T\). All supports here
lie in \(\pi^{-1}(\operatorname{supp}\chi)\), independently of
\(\varepsilon\). There is no assertion of positive-test preservation.

\subsection{9. Every compact ordinary scalar test; the two limits and
support}\label{every-compact-ordinary-scalar-test-the-two-limits-and-support}

Let invariant \(\beta\in A_c^{0,0}(X_E)\) be arbitrary, and set
\(F=\pi(\operatorname{supp}\beta)\). It is compact. Cover \(F\) by
finitely many ordinary interiors \(U_i\) of strict affinoids \(V_i\) on
which (8) holds, and choose outer cutoffs \(\chi_i\) from (7), compactly
supported in \(U_i\), with sum one near \(F\). Put
\(C=\bigcup_i\operatorname{supp}\chi_i\). It is a fixed compact set
containing \(F\).

On \(V_{i,E}\) approximate the fixed Euclidean coefficient \(\Phi_i\)
through order two on a box containing the fixed unit moment image by
polynomials \(P_{i,j}\). Choose the box to contain all translated tuple
images as well. Let \(T_{i,j}\) be the real group average of
\(P_{i,j}(T_i)\). Since \(\beta\) is invariant, it is also the average
of \(\Phi_i(T_i)\), so \(\beta-T_{i,j}\) tends to zero in coefficient
\(C^2\) on a \textbf{fixed upstairs} moment tuple containing all the
translates. Include fixed moment descriptions of \(\chi_i\circ\pi\) on a
finite compact cover of its support. The product rule and (19) give \[
 \big\|(D((\chi_i\circ\pi)(\beta-T_{i,j})))_{
                     \mathrm{Rad},E}\big\|_{\mathrm{TV}}\to0.       \tag{23}
\]

For each fixed \(j\), use (9) to express \(T_{i,j}\) as finitely many
traces. Apply (21)--(22) with outer cutoff \(\chi_i\) to obtain
\(a_{i,j}\) supported in \(\operatorname{supp}\chi_i\) and with \[
 \big\|(D((\chi_i\circ\pi)T_{i,j}-\pi^*a_{i,j}))_{
                     \mathrm{Rad},E}\big\|_{\mathrm{TV}}<1/j       \tag{24}
\] (or \(1/(ij)\) after any fixed finite enumeration). Each
fixed-polynomial local cover, inner partition and moment refinement is
chosen \textbf{before} its smoothing limit. It may depend on \(j\). Its
derivative bounds may grow with \(j\); smoothing is chosen after those
bounds. This causes no problem because (23) is evaluated with the fixed
outer cutoffs, before the inner localization of \(T_{i,j}\). The varying
inner cutoffs never multiply the polynomial approximation error in (23),
and every approximate support remains inside the same outer \(C\).
Neither a degree-uniform estimate nor an infinite simultaneous unit
shrink is asserted.

Set \(\alpha_j=\sum_i a_{i,j}\). Globally
\(\beta=\sum_i(\chi_i\circ\pi)\beta\), since the sum is one near its
support and it vanishes off that support. Equations (23)--(24) and the
finite triangle inequality prove (2). The support of every
differentiated error is contained in \(\pi^{-1}(C)\), compact by Section
3. Thus for \textbf{every} continuous finite multiplier \(h\), \[
       B_{h\circ\pi}(\beta-\pi^*\alpha_j)\longrightarrow0,
                                                               \tag{25}
\] using its boundedness on that compact set. This proves the actual
ordinary scalar test gate for curves, not merely \(C^0\) density of
invariant functions.

\subsection{10. Forward Galois ascent and actual normal-closure
descent}\label{forward-galois-ascent-and-actual-normal-closure-descent}

For a general compact ordinary scalar test \(\beta\) upstairs, (6),
ordinary exterior functoriality on the restricted carrier, and
invariance of \(h\circ\pi\) show
\(B_{h\circ\pi}(g^*\beta)=B_{h\circ\pi}(\beta)\). Its real average
\(\beta^{\mathrm{av}}=d^{-1}\sum_g g^*\beta\) has the same pairing. Its
support is in a compact finite group saturation. Real averaging does not
divide by \(d\) in \(K\), including when \(\operatorname{char}K\)
divides \(d\).

For every downstairs compact test \(\alpha\), its pullback is compact.
The scalar projection formula (4), with \textbf{no degree factor}, gives
\[
  B_{h\circ\pi}(\pi^*\alpha)
     =\int_{X_E}(h\circ\pi)(\pi^*D\alpha)_{\mathrm{Rad},E}
     =\int_X h(D\alpha)_{\mathrm{Rad},K}=B_h(\alpha).     \tag{26}
\] Smooth exterior functoriality suffices; no continuous-current
pullback has been presumed. If the downstairs current is zero,
approximate \(\beta^{\mathrm{av}}\) by (2). Its approximate pairings
vanish by (26), and (25) proves that its pairing, hence that of every
upstairs test, vanishes. This proves forward ascent for every finite
Galois extension, including wild ones, on all curves in the stated
scope.

Conversely, (26) alone proves reverse descent whenever the upstairs
current is zero. This calculation works for every finite separable
extension, without a group, by compactness, purity, and (4). It is an
actual proof of the descent needed below, not a use of an earlier note.

For arbitrary finite separable \(E/K\), take the finite Galois normal
closure \(N/K\) constructed in Section 3, with its unique extending
absolute value and completeness. Forward Galois ascent gives zero
current on \(X_N\). The space \(X_E\) retains every ambient hypothesis
and pure dimension one by Section 3, and \(N/E\) is finite separable.
There is a canonical analytic identification \((X_E)_N=X_N\): on
affinoids it is the associative tensor isomorphism
\((A\otimes_K E)\otimes_E N=A\otimes_K N\) with the equivalent
finite-module topologies. The pullback of \(h\circ\pi_E\) is the
original pullback on \(X_N\). Apply the just-proved reverse calculation
to \(N/E\), testing each compact ordinary function on \(X_E\), to obtain
zero current on \(X_E\). This proves (1) at every finite separable
stage. It is the finite dependency, before the completed stage in
Section 13. Ordinary opens retain the stated assumptions and the same
argument applies to their scalar inverse images.

\subsection{11. Finite inseparability and arbitrary finite algebraic
stages}\label{finite-inseparability-and-arbitrary-finite-algebraic-stages}

\subsubsection{11.1. Frobenius, topology and exact scalar
tests}\label{frobenius-topology-and-exact-scalar-tests}

Let \(k\) be complete, nontrivially valued, of characteristic \(p>0\),
and let \(F/k\) be finite purely inseparable. Choose \(q=p^e\) with
\(F^q\subset k\). Section 3's bounded-coordinate argument applies to
\textbf{every} finite extension, and gives
\(A\widehat\otimes_k F=A\otimes_k F\) with its finite-module topology.
Consequently every \(u=\sum a_i\otimes\lambda_i\) satisfies \[
                    u^q=\sum a_i^q\lambda_i^q\in A.       \tag{27}
\] This is an identity in the actual algebra, including nilpotents.

Scalar projection is surjective by the affine foundation in Section 2;
Section 12.1 also gives a direct proof for all fields used here. If two
points of \(\mathcal M(A\otimes F)\) project to \(x\), (27) forces their
values on every \(u\) to equal \(|u^q(x)|^{1/q}\). They are the same
seminorm. The projection is thus a continuous bijection from a compact
space to a Hausdorff space, hence a homeomorphism. On a good separated
space these affine homeomorphisms glue on the ordinary interiors of
downstairs affinoid neighbourhoods. No covering-space assertion is used.

For a unit \(f\in(A\otimes F)^\times\), (27) gives \(f^q=g\in A\). The
function \(g\) has no analytic zero downstairs, by surjectivity, so it
is an actual unit by Section 3's unit criterion. On the upstairs carrier
\(t_f=q^{-1}t_g\). Now let \(\beta\in A_c^{0,0}(X_F)\) be arbitrary. Its
support has compact image \(C\) under the homeomorphism. For each point
of \(C\), select an ordinary local description of \(\beta\), then a
downstairs affinoid neighbourhood \(V\) whose full inverse image lies in
that description open. The inverse image is \(V_F\), so all coordinate
units restrict to elements of its actual affinoid algebra. Taking their
\(q\)th powers gives an ordinary downstairs scalar description
\(\varphi(q^{-1}t_{g_1},\ldots,q^{-1}t_{g_m})\) on
\(\operatorname{int}_X V\). Equality of numerical smooth scalar
functions is equality of their ordinary \((0,0)\) sections, as explained
in Section 12.3. Cover \(C\) by finitely many such opens, apply the
compact partition of Section 4, multiply these local descriptions by its
cutoffs, and extend by zero. This produces \(\alpha\in A_c^{0,0}(X)\)
with \(\beta=\pi_F^*\alpha\) everywhere. Compact support inside each
open makes the extension ordinary smooth; no derivative is extended
through an exposed boundary.

Goodness, strictness, separatedness and boundarylessness are preserved
as in Section 3. Purity is also preserved, and here is a proof that does
not assume equality of total dimensions gives local purity. Given an
upstairs point \(y\), a downstairs affinoid neighbourhood of its image
has dimension one, and its scalar extension has dimension one, giving
the upper local bound. If local dimension at \(y\) were zero, take this
affinoid upstairs and remove all its positive-dimensional irreducible
components, none of which contains \(y\). An affinoid neighbourhood
inside the remaining ordinary open has dimension zero by CLD 11.2.3's
component/domain formula. By the homeomorphism choose a downstairs
affinoid neighbourhood whose full inverse image is inside that open. Its
downstairs dimension is one (it is a domain in the pure curve), and its
upstairs total dimension is one by scalar dimension invariance, a
contradiction. Existence of the smaller downstairs affinoid
neighbourhood uses goodness. This argument only checks dimension; it
never reduces the measure.

On these pure curves \(D\beta=\pi_F^*D\alpha\), and the smooth
top-measure identity (4) proves \[
             B_{h\circ\pi_F}(\beta)=B_h(\alpha).         \tag{28}
\] All supports are compact. The converse is obtained by taking
\(\beta=\pi_F^*\alpha\) for an arbitrary compact downstairs scalar test.
Thus finite purely inseparable zero-Hessian ascent and descent are
proved here directly. No earlier inseparable manuscript is a premise.

\subsubsection{11.2. The finite mixed-field
decomposition}\label{the-finite-mixed-field-decomposition}

Let \(E/K\) be any finite algebraic extension. In characteristic zero
every irreducible polynomial has nonzero derivative, and is separable,
so Sections 3--10 apply. In characteristic \(p\), choose finitely many
generators \(a_1,\ldots,a_r\) of \(E/K\). For each irreducible minimal
polynomial, repeatedly remove common factors \(p\) from its exponents
until it has the form \[
                f_i(Z)=g_i(Z^{p^{e_i}}),\qquad g_i'\ne0.
                                                               \tag{29}
\] The polynomial \(g_i\) is irreducible: a factorization of \(g_i\)
would give a factorization of \(f_i\) by substitution. Since \(g_i\) is
irreducible and its derivative is nonzero of smaller degree,
\(\gcd(g_i,g_i')=1\); its roots are distinct. Therefore
\(b_i=a_i^{p^{e_i}}\) is separable over \(K\). The compositum
\(S=K(b_1,\ldots,b_r)\subset E\) is finite separable: adjoining a root
of a separable polynomial remains separable over an intermediate field
because its minimal polynomial divides that polynomial; extending
embeddings successively gives degree-many embeddings of the tower. An
intermediate simple field has degree-many embeddings as well: each
embedding has at most the remaining extension degree lifts, while the
entire tower has its full degree of embeddings. Thus its minimal
polynomial is separable. This proves the field facts used here without
assuming separability of an arbitrary algebraic stage.

With \(e=\max e_i\), every \(a_i^{p^e}\in S\), and taking \(p^e\)th
powers in a rational expression in the \(a_i\) shows
\(E^{p^e}\subset S\). Hence \(E/S\) is finite purely inseparable. Both
fields are complete under the prescribed absolute value by the finite
norm proof in Section 3. This gives the \textbf{actual embedded} tower
\(K\subset S\subset E\), not an abstract factorization of degrees.

The finite tensor associativity with bounded inverse maps, proved in
Section 12.1, identifies \((X_S)_E=X_E\) as analytic spaces, with the
actual algebra and not its reduction. \(X_S\) has all ambient hypotheses
and pure curve dimension by Section 3. Finite separable ascent to \(S\),
then the proved Frobenius ascent to \(E\), proves (1) at every finite
algebraic stage. Reverse descent uses (4) directly. Generic lengths are
those of the actual tensors throughout. No degree factor is introduced
into scalar projection.

\subsection{12. Completion, rational neighbourhoods and every compact
scalar
test}\label{completion-rational-neighbourhoods-and-every-compact-scalar-test}

\subsubsection{12.1. Fields, completed tensors, surjectivity and
compactness}\label{fields-completed-tensors-surjectivity-and-compactness}

Write \(M=\bigcup_E E\), where \(E\) ranges over its finite
subextensions of \(K\). It is a directed union: a finite list of
algebraic elements generates a finite field, by successive
multiplication of finite degrees. Each \(E\) is complete by Section 3.
We never regard \(M\) as a complete analytic field unless it happens
already to be complete. Its completion \(L\) is a complete valued field.
For a \(K\)-affinoid algebra \(A\), the image of \(A\otimes_K M\) is
dense in \(A_L=A\widehat\otimes_K L\). Indeed finite algebraic tensors
are dense by definition of completion; a tensor
\(\sum a_j\otimes\lambda_j\) can be approximated by replacing its
finitely many \(\lambda_j\in L\) with \(m_j\in M\), with tensor error at
most \(\max_j\|a_j\||\lambda_j-m_j|\). Approximation statements refer to
the \textbf{image}, so no injectivity or isometry of an uncompleted
infinite tensor is presumed.

Here is the associative tensor argument actually required. For finite
\(E/K\) with a basis \(e_i\), the bounded coordinate maps give
\(A_E=A\otimes_K E\) with its finite-module topology. The maps on
algebraic tensors \[
 (a\otimes e)\otimes\lambda\mapsto a\otimes e\lambda,
 \qquad a\otimes\lambda\mapsto(a\otimes1)\otimes\lambda
                                                               \tag{30}
\] extend continuously to inverse maps
\((A_E)\widehat\otimes_E L\simeq A\widehat\otimes_K L\). For the first
map, the bounded multiplication map
\((b,\lambda)\mapsto\sum_i a_i\otimes e_i\lambda\), where
\(b=\sum_i a_i\otimes e_i\), is bounded in \(b\) by the coordinate
estimates and in \(\lambda\) by the tensor inequality. It is balanced
over \(E\). The reverse map is bounded by the tensor inequality and the
bounded inclusion \(A\to A_E\). These maps are inverse on dense
elementary tensors, hence inverse on the completions. The same proof
works when \(A\) is affinoid with arbitrary radii, and for two finite
stages. Thus all stage projections and inverse images of rational
domains used below are actual analytic base changes. No associativity
claim for an infinite uncompleted Banach tensor is needed.

For clarity, scalar surjectivity in the present completed-algebraic
scope also has a direct seminorm proof. Given \(x\in\mathcal M(A)\), put
\(H=\mathcal H(x)\). Embed the algebraic extension \(M\) into an
algebraic closure \(\overline H\) over \(K\). Such an embedding is
obtained by successively adjoining roots, or by maximizing partial
embeddings and adjoining a root of the minimal polynomial if the domain
is not all of \(M\). Equip \(\overline H\) with the extending absolute
value: the finite extension existence and uniqueness proof of Section 3
makes the norms on its finite \(H\)-subfields compatible. On every
finite \(K\)-subfield of \(M\) this value is the specified one, again by
uniqueness. Completing \(\overline H\) gives a complete field \(F\) into
which the isometric embedding \(M\to F\) extends to \(L\to F\). For
\textbf{every} tensor representation with \(\lambda_j\in L\), evaluation
satisfies \[
 \left|\sum_j a_j(x)\lambda_j\right|
       \le \max_j\|a_j\||\lambda_j|.
                                                               \tag{31}
\] Taking the infimum over all representations bounds it by the full
tensor seminorm. It extends to \(A_L\) and gives an analytic point over
\(x\). This proof handles the separated quotient in the tensor
completion and uses no claimed isometry of \(A\otimes M\). The identical
proof over \(E\) gives surjectivity \(X_L\to X_E\) and surjectivity on
every scalar-extended affinoid domain. On good spaces these affine facts
cover the whole space.

The inverse image of a compact \(C\subset X_E\) is compact: finitely
many ordinary interiors of affinoid neighbourhoods cover \(C\). Its
inverse image is closed in the finite union of their compact scalar
extensions, since the spaces are Hausdorff. This works for the infinite
completed projection. It does not assert a finite fibre, a finite Galois
cover or any infinite group averaging.

For a nonzero \(\lambda\in L\), choose \(m\in M\) with
\(|m-\lambda|<|\lambda|\). Then \(|m|=|\lambda|\). Thus
\(|L^\times|=\bigcup_E|E^\times|\). This observation handles any
value-group radius if desired, but the rational argument below allows
\textbf{all positive real radii}, so it needs no strictifying extension
and no approximation of a radius.

\subsubsection{12.2. Exact rational and strict-locus
descent}\label{exact-rational-and-strict-locus-descent}

Let \(A\) be affinoid over a finite stage \(E_0\), set
\(B=A\widehat\otimes_{E_0}L\), and let a rational domain of
\(\mathcal M(B)\) be defined by \[
             |f_i|\le r_i|g|\quad(1\le i\le s),
             \qquad r_i>0,                              \tag{32}
\] where \(g,f_1,\ldots,f_s\) have no common analytic zero. Empty
systems can instead be handled by the whole affinoid with the harmless
condition \(|1|\le2|1|\). Compactness gives \[
 \delta=\min_{y\in\mathcal M(B)}
             \max\{|g(y)|,r_1^{-1}|f_1(y)|,\ldots,
                                  r_s^{-1}|f_s(y)|\}>0.
                                                               \tag{33}
\] Choose simultaneous finite-stage approximants \(g',f_i'\) in one
\(A_E\) with Banach errors less than \(\delta\min(1,r_1,\ldots,r_s)\).
This is possible by Section 12.1. Their evaluation errors satisfy \[
                 |g'-g|<\delta,
                 |f_i'-f_i|<r_i\delta                 \tag{34}
\] at every point upstairs. Neither vanished nor nilpotent functions are
logarithmed, and none is assumed to have small Banach norm merely
because its spectral norm is zero. They are approximated in the actual
Banach algebra just like the other elements.

Both conjunctions, (32) and the conjunction of its strict versions, are
\textbf{unchanged} by (34). To check all boundary cases, first suppose
\(|g(y)|\ge\delta\). Then \(|g'(y)|=|g(y)|\). For each \(i\) with
\(|f_i(y)|\ge r_i|g(y)|\), (34) forces \(|f_i'(y)|=|f_i(y)|\). For each
\(i\) with \(|f_i(y)|<r_i|g(y)|\), it instead forces
\(|f_i'(y)|<r_i|g'(y)|\). Hence each weak inequality, each equality, and
each strict inequality is preserved at that point. If \(|g(y)|<\delta\),
(33) supplies an \(i\) with \(r_i^{-1}|f_i(y)|\ge\delta\). Its numerator
norm is unchanged, while \(|g'(y)|<\delta\). This inequality remains
strictly false, so both conjunctions remain false. These arguments also
cover \(g(y)=0\) and any number of zero numerators.

The weighted maximum for the approximants is everywhere positive
upstairs. Surjectivity to \(\mathcal M(A_E)\) then shows the
denominator/numerators have no common zero there. They generate the unit
ideal: otherwise a maximal ideal containing them gives an analytic
common zero. Thus they define a genuine rational analytic domain \(R_E\)
by CLD 11.1.5. Its scalar extension has exactly the set (32). Analytic
domain uniqueness, or the universal rational-domain presentation by
adjoining the bounded ratios, identifies it with the given domain
\(R_L\) as an analytic domain. The actual rational localization algebra,
with all nilpotents, is retained. The strict locus
\(O_E=\{|f_i'|<r_i|g'|\ \forall i\}\) is ordinary open in the ambient
affinoid, lies in its rational domain, and has inverse image
\textbf{exactly} the original strict locus \(O_L\). No invariance of
absolute analytic interiors or equality of all topological interiors is
asserted.

The stronger rational-neighbourhood statement of CLD 11.1.5 says that
for any \(y\) and any prescribed ordinary open in an affinoid, one can
choose (32) with \(R_L\) inside that open and \(y\) in its strict locus.
Thus the preceding construction supplies finite-stage ordinary
neighbourhoods of each chosen point, without controlling an entire
infinite scalar fibre. Only finitely many such choices are subsequently
needed on a compact support.

\subsubsection{12.3. Exact finite-stage local
descriptions}\label{exact-finite-stage-local-descriptions}

Let \(\beta\) be ordinary smooth on \(X_L\) near \(y\), and let \(x\) be
its image in \(X\). Choose a downstairs affinoid neighbourhood \(V_0\)
and an ordinary open \(U_0\subset V_0\) containing \(x\). The inverse
image \(U_{0,L}\) is open. Intersect it with an ordinary moment
description open of \(\beta\); on that intersection write
\(\beta=\varphi(t_{f_1},\ldots,t_{f_m})\), with actual analytic units
and a smooth real coefficient defined on an open Euclidean set
containing the coordinate image. This can always be arranged locally by
shrinking the description near \(y\).

Inside \(V_{0,L}\) select a rational neighbourhood \(R_L\) of \(y\)
contained in this intersection, with \(y\) in its strict locus. Section
12.2 gives \(R_L=(R_{E_0})_L\) for a finite stage \(E_0\), and a
descended strict locus. Because \(R_L\subset U_{0,L}\), surjectivity on
\(R_{E_0}\) gives \(R_{E_0}\subset U_{0,E_0}\). Its strict locus is
therefore an ordinary open of the whole \(X_{E_0}\), not just an open of
an exposed affinoid boundary.

On the compact rational affinoid \(R_L\), put \(\delta_j=\min|f_j|>0\).
Apply density in the actual rational affinoid algebra, now viewed over
\(E_0\), to select a common finite \(E\supset E_0\) and
\(u_j\in\mathcal O(R_E)\) with \(\|u_j-f_j\|<\delta_j\) upstairs.
Ultrametricity gives \[
                     |u_j(y')|=|f_j(y')|
                                 \quad(y'\in R_L).     \tag{35}
\] By surjectivity \(u_j\) has no zero on \(R_E\), so it is an actual
unit by the unit criterion. Its inverse exists in that affinoid algebra;
no approximation of an inverse or reduction of the algebra is necessary.
Surjectivity also makes the finite-stage moment image of the \(u_j\)
equal to the upstairs image of the \(f_j\), so the same Euclidean
coefficient \(\varphi\) is defined on it.

On the inverse image \(O_E\) of the strict locus define the ordinary
smooth function \(b_E=\varphi(t_{u_1},\ldots,t_{u_m})\). Then
\(q_E^*b_E=\beta\) on \(O_L=q_E^{-1}O_E\). This is an exact
scalar-section identity. To justify the section claim precisely, CLD
3.4.6 embeds the \((0,0)\) G-form sheaf in numerical functions: on each
compact chart the nullity criterion 3.4.1/3.4.4 says that a zero
numerical coefficient is the zero section; compact charts G-cover the
space. Hence equal numerical scalar sections coincide, and ordinary
scalar sections are included in this sheaf. Alternatively, the combined
moment \((f,u)\) has image in the diagonal \(\{t_f=t_u\}\) by (35); the
difference of the two scalar coefficient functions restricts to zero
there, and moment nullity gives the same conclusion. No universal
moment-dimension equality is used. Smooth differentiation consequently
commutes with this \textbf{exact identity}, giving \(D\beta=q_E^*Db_E\)
on that open.

\subsubsection{12.4. Compact scalar tests and extension by
zero}\label{compact-scalar-tests-and-extension-by-zero}

Now let \(\beta\in A_c^{0,0}(X_L)\) be arbitrary and put
\(F=\operatorname{supp}\beta\). If \(F\) is empty, take \(\alpha=0\) at
any finite stage. Otherwise Section 12.3 gives a neighbourhood
\(O_{i,L}=q_{E_i}^{-1}O_{i,E_i}\) and a descended ordinary function
\(b_{i,E_i}\) near each point of \(F\). Choose finitely many covering
\(F\). Enlarge to a common finite field \(E\subset M\) and base change
all these rational domains, strict loci and unit lists to \(E\).
Associativity (30) shows \(O_{i,L}=q_E^{-1}O_{i,E}\) and
\(\beta=q_E^*b_{i,E}\) there.

The compact set \(F_E=q_E(F)\) is covered by the finitely many
\(O_{i,E}\). Apply Section 4's explicit compact partition, built only
from CLD 15.2.5, to obtain ordinary smooth \(\theta_i\) with compact
support in \(O_{i,E}\) and sum one near \(F_E\). Put \[
                \alpha=\sum_i\widetilde{\theta_i b_{i,E}},
                                                               \tag{36}
\] where tilde means extension by zero to \(X_E\). Each product is
ordinary smooth on its open, with support in the compact
\(\operatorname{supp}\theta_i\subset O_{i,E}\). Since that compact set
is closed in the Hausdorff \(X_E\), every point outside it has an
ordinary zero neighbourhood. Therefore its extension is ordinary smooth
at every exterior point as well. This proves (36) lies in
\(A_c^{0,0}(X_E)\), with no boundary derivative term.

On the inverse image of its defining open each summand pulls back to
\((\theta_i\circ q_E)\beta\). Off that open its pullback is zero, and
\(\theta_i\circ q_E\) is also zero. Thus, \textbf{globally}, \[
 q_E^*\alpha=(\sum_i\theta_i\circ q_E)\beta=\beta.       \tag{37}
\] On \(F\) the sum is one. Outside \(F\) the function \(\beta\) is
zero, so the equality also holds there regardless of that sum. No
descent or saturation of the original support is assumed. No agreement
of the separate \(b_{i,E}\) on downstairs overlaps needs to be presumed;
the finite sum is itself a smooth section, and (37) proves precisely its
required pullback. Its support is contained in a finite union of compact
cutoff supports. The inverse image of that union is compact by Section
12.1. Differentiation commutes with the exact ordinary smooth identity
(37), so every compact upstairs scalar test and its Hessian are covered.
This is stronger than \(C^0\) density and does not use a differentiated
approximation topology at completion.

\subsubsection{12.5. Ambient geometry and local purity after
completion}\label{ambient-geometry-and-local-purity-after-completion}

Goodness follows because inverse images of downstairs affinoid
neighbourhoods are affinoid and contain inverse images of their ordinary
interiors. Strict presentations base change to strict presentations. The
diagonal base change preserves separatedness. Boundarylessness follows
from GR 2.11(5'), as stated in Section 2; ordinary open restrictions
retain it by GR 2.11(3),(4).

We check local purity for \(X_L\), rather than just total dimension.
Around the image of \(y\in X_L\) take a downstairs strict affinoid
neighbourhood \(V_0\) of dimension one. Then \(V_{0,L}\) has dimension
one, giving the upper local bound. Suppose local dimension at \(y\) were
zero. In \(V_{0,L}\) discard its finitely many dimension-one
minimal-prime components, none of which contains \(y\). In the remaining
ordinary open choose a rational affinoid neighbourhood \(R_L\) of \(y\)
inside the inverse image of \(\operatorname{int}_X V_0\), with strict
inequalities at \(y\). The component/domain dimension formula in CLD
11.2.3 gives \(\dim R_L=0\). Section 12.2 descends \(R_L\) to a nonempty
\(R_E\) at a finite field, inside \(X_E\). The actual \(X_E\) is pure
one-dimensional by Sections 3 and 11. Its domain \(R_E\) has dimension
one, including if its chosen real radii make the domain non-strict.
Scalar dimension invariance and analytic domain identification give
\(\dim R_L=\dim(R_E)_L=\dim R_E=1\), a contradiction. Hence every
upstairs point has dimension one. This argument keeps the original
strict ambient space and does not infer an endpoint theorem from a
non-strict radial core.

All top measures below are on the \textbf{actual} pure curves just
proved to exist. Their multiplicities need not remain numerically
unchanged point by point. They are the finite-flat fibre-algebra
dimensions in the canonical calibration. For reference, the usual
generic-length identity is obtained by tensoring a composition series of
a generic Artinian stalk with the flat upstairs generic stalk: each
residue-field quotient becomes the generic Artinian algebra of the
base-changed reduction. If its length is \(\mu\) and the original length
is \(\ell\), the length upstairs is \(\mu\ell\). This is CLD 13.2.8 with
\(\mathcal F=\mathcal O_X\); its proof body is included in Section 15.
We use the calibration and the Artinian degree-sum proof of (4), not a
replacement by geometric point counts or a supposed constant extension
degree for an infinite extension.

\subsection{13. Completed forward ascent and actual reverse
descent}\label{completed-forward-ascent-and-actual-reverse-descent}

Assume \(D[h]=0\) downstairs. Given \textbf{any}
\(\beta\in A_c^{0,0}(X_L)\), Section 12.4 supplies finite \(E\) and
\(\alpha\in A_c^{0,0}(X_E)\) with \(\beta=q_E^*\alpha\) exactly. Put
\(h_E=h\circ\pi_E\). Finite algebraic ascent, proved in Sections 3--11,
gives \(D[h_E]=0\). Section 12.5 verifies the pure curve and ambient
hypotheses for both \(X_E\) and \(X_L\). Smooth functoriality and the
\textbf{\(L\)- and \(E\)-calibrated} projection formula (4) yield \[
 \begin{aligned}
 B_{h\circ\pi}(\beta)
 &=\int_{X_L}(h_E\circ q_E)(q_E^*D\alpha)_{\mathrm{Rad},L}\\
 &=\int_{X_E}h_E(D\alpha)_{\mathrm{Rad},E}
 =0.
 \end{aligned}                                             \tag{38}
\] The relevant supports lie in the compact inverse image of
\(\operatorname{supp}\alpha\); the continuous finite multiplier is
bounded there. The signed-measure pushforward formula is therefore
legitimate. Since \(\beta\) was arbitrary, the upstairs zero-current
equation follows. There is no pullback of an arbitrary continuous
current in this proof and no field-degree factor in (38).

Conversely, suppose the upstairs equation holds. For an arbitrary
\(\psi\in A_c^{0,0}(X)\) its pullback has compact support by Section
12.1. Apply smooth functoriality and (4) directly for \(L/K\): \[
             B_h(\psi)=B_{h\circ\pi}(\pi^*\psi)=0.        \tag{39}
\] Thus the downstairs equation holds. This reverse calculation is
included for completeness, not counted as new forward progress.
Ordinary-open localization follows from the ambient checks and the same
arguments on that open; no global quasicompactness or paracompactness is
added. This proves (1) with all its stated scope.

\subsection{14. Comparison, fallibility and
reproduction}\label{comparison-fallibility-and-reproduction}

The controlling P06.2 chapter asks for \textbf{classical pluriharmonic
recognition} of arbitrary continuous zero-current functions on good
separated strict boundaryless spaces, or a counterexample with all those
ambient hypotheses. Neither conclusion is claimed here. The new gate is
unconditional ordinary CLD scalar-current ascent on arbitrary pure
curves under the completion of an arbitrary algebraic extension,
including singular and nonreduced curves and the completed algebraic
closure. It goes beyond merely restating the finite separable curve
draft or imposing scalar stability.

The finite separable dependency has been supplied with its entire
scalar-polarization, root-valuation, tangential-convolution and
calibrated carrier proof. The finite inseparable dependency has been
supplied with its actual tensor, Frobenius and compact scalar-test
proof. Neither predecessor's author audit, publication, native note,
hash, or review verdict is accepted as a mathematical premise. The
finite-stage support and domain gate for completion has been proved
directly. No arbitrary infinite fibre is treated as finite.

In dimension at least two, complementary Hessian tests have positive
bidegree, and the finite separable reconstruction here is only scalar.
Therefore even an extension of the completed-stage argument to all form
degrees would not by itself prove the missing finite differential-form
theorem. Transcendental scalar extensions, universal curve probes,
normalization transport for recognition, and two-sign regularization
remain separate unresolved obligations. Zero opposite currents after
\(L/K\) do not put either continuous sign into the generated CLD-psh
hierarchy. Ordinary smoothness, weak smoothness, real-PL regularity and
GJR harmonicity of \(h\) are not asserted. The catalogue's non-strict
arbitrary-radius invisible boundary example is outside the original
strict boundaryless scope and is not a counterexample supplied here.

The authoritative file is \passthrough{\lstinline!P06.2-argument.md!},
with the complete standalone \passthrough{\lstinline!P06.2-partial.tex!}
and twice-compiled \passthrough{\lstinline!P06.2-partial.pdf!}.
\passthrough{\lstinline!self-audit.md!} and
\passthrough{\lstinline!research-notes.md!} record the fallible audit
and discarded shortcuts. \passthrough{\lstinline!status.json!} has
exactly
\passthrough{\lstinline!id,status,exact\_result,proof\_file,concrete\_gaps!}.
\passthrough{\lstinline!source-manifest.json!} records copied input
origins and hashes;
\passthrough{\lstinline!evidence/primary-passages.txt!} preserves the
primary bytes and
\passthrough{\lstinline!evidence/artifact-hashes.json!} records build
hashes. The two parent native notes are preserved as data, not
dependencies.

Run \passthrough{\lstinline!python3 prepare-sources.py!} followed by
\passthrough{\lstinline!python3 build-and-check.py!} in this fresh
workspace. Only installed tools and bounded standard-library diagnostics
are used. The finite diagnostic tests perturb exact rational numbers in
the 2-adic norm and check (32)--(34), retaining zero denominators, zero
numerators, equalities and both conjunctions. They do not verify an
analytic theorem. Two successful TeX passes likewise establish
typesetting, not correctness. No heavy computation was required; no
supervised compute service was exposed or used. There is no independent
model review, formal certification, inherited verdict or priority claim.
The manuscript remains an unverified draft,
\passthrough{\lstinline!status=partial!},
\passthrough{\lstinline!candidate=false!}.

\subsection{Bibliography}\label{bibliography}

\begin{itemize}
\tightlist
\item
  A. Chambert-Loir and A. Ducros, \emph{Formes différentielles réelles
  et courants sur les espaces de Berkovich}, arXiv:1204.6277v2, 25 July
  2025. Used statements and proof bodies are specified in Section 2 and
  included below. The rational-domain completion lemma is proved in
  Section 12, not attributed to a source slogan.
\item
  W. Gubler and J. Rabinoff, \emph{Classically psh and pluriharmonic
  functions on Berkovich spaces}, arXiv:2506.13548v2, 17 September 2025,
  Section 2.11(3)--(5).
\item
  Controlling chapter and family dossier: scope data only. Prior finite
  separable and inseparable manuscripts: fallible evidence; their needed
  arguments are reconstructed in this file.
\end{itemize}

\subsection{15. Included primary statements and proof
bodies}\label{included-primary-statements-and-proof-bodies}

These are the inspected primary snapshot passages, included here so a
reader can examine the exact cited statements and the available proof
bodies without filesystem access. They are source data, with their
text-extraction artefacts retained; Section 2 and the direct arguments
give the precise English statements used. Basic affinoid construction,
noetherianity and the analytic dimension foundations are cited in CLD's
foundational recall paragraphs; this packet does not claim to rederive
all of Berkovich geometry from set theory. Their exact recalled
statements are included. The new finite-stage domain/test descent is
proved in full in Section 12, and the needed strong neighbourhood and
compact bump constructions are supplied directly in Section 2.1.

In the quoted 12.2.4 wording, the universal dimension-equality phrase is
\textbf{not used}. Only compact polyhedrality is used there; the needed
dimension upper bound is the separately included 12.2.10. Constant
moments disprove a universal equality. No broad 15.2.9 partition
assertion is used. The only partition is the finite compact construction
of Section 4 from 15.2.5, whose local cutoff proof was also
reconstructed in Section 2.1. Other hypotheses in these passages are not
silently relaxed.

The PDF transliterates the source-excerpt code blocks to ASCII for the
installed TeX engine, explicitly naming any symbol without a standard
ASCII equivalent. The authoritative Markdown retains the exact Unicode
text. The English mathematical argument is unchanged.

\subsubsection{CLD.txt: cached text lines
14838--15060}\label{cld.txt-cached-text-lines-1483815060}

\begin{lstlisting}
s[RIGHT SINGLE QUOTATION MARK]il possede un pre-modele (V, F ) ou F est engendre par ses sections globales.
Remarque 10.7.8. -- Supposons que L soit de type fini et degre de transcendance 1
sur K. C[RIGHT SINGLE QUOTATION MARK]est alors comme dans le cas non gradue le corps gradue des fonctions d[RIGHT SINGLE QUOTATION MARK]une
K-courbe graduee projective integre et normale C qui est uniquement determinee, et
la fleche naturelle rho : PL/K -> C est un isomorphisme d[RIGHT SINGLE QUOTATION MARK]espaces localement anneles
gradues ; tout fibre vectoriel E sur un ouvert quasi-compact U de PL/K admet des lors
un modele canonique, a savoir le couple (rho(U ), rho* E).
Lemme 10.7.9. -- Soit U un ouvert quasi-compact de PL/K et soit E un fibre en
droites sur U . Il existe deux fibres en droites D1 et D2 sur U engendres par leurs sections
globales et un isomorphisme E ~= D1  tensor  D2[LOGICAL OR] .
Demonstration. -- Choisissons un pre-modele (V, G) de (U, E). D[RIGHT SINGLE QUOTATION MARK]apres le corollaire
10.7.4 on peut supposer V quasi-projectif. Il resulte alors de 10.3.7 qu[RIGHT SINGLE QUOTATION MARK]il existe deux fibres
en droites F1 et F2 sur V , engendre par leurs sections globales, et tels que G ~= F1  tensor  F2[LOGICAL OR] .
Il suffit pour conclure de prendre pour D1 et D2 les images reciproques respectives de
F1 et F2 sur U .

292        CHAPITRE 10. ALGEBRE GRADUEE ET REDUCTION DES ANNEAUX NORMES

Exemple 10.7.10. -- Soit D une droite vectorielle graduee sur L et soit t un element
homogene non nul de D, de degre r. Soit (ai )1<=i<=n une famille finie d[RIGHT SINGLE QUOTATION MARK]elements de L x 
algebriquement independants sur K et soit E le fibre en droites sur PL/K dont la fibre
generique est D et dont la metrique est telle que v(t-1 ) = max(1, v(a1 ), . . . , v(an ))
(lemme 10.6.18).
                 (1,|a |,...,|an |)
   Identifions PK 1                 a un premodele de L sur K via le plongement K(Ti /T0 )1<=i<=n ,->
L, (Ti /T0 ) 7-> ai dont nous notons F l[RIGHT SINGLE QUOTATION MARK]image. Soit [INCREMENT] l[RIGHT SINGLE QUOTATION MARK]image reciproque sur PF/K du
                                       1,|a |,...,|an |
fibre en droites O(1)[r] de PK 1                        . La fibre generique de [INCREMENT] est egale a F * T0 ou T0
est considere comme homogene de degre r. Les elements T0 et Ti = ai T0 constituent
une famille de sections globales de [INCREMENT] qui engendrent ce dernier. Il resulte alors du
lemme 10.6.15 que la metrique sur F * T0 au-dessus de PF/K induite par [INCREMENT] est telle
que v(T0 )-1 = max(1, v(a1 ), . . . v(an )) pour tout v. Il s[RIGHT SINGLE QUOTATION MARK]ensuit que l[RIGHT SINGLE QUOTATION MARK]image reciproque de
                                                                               (1,|a |,...,|an |)
[INCREMENT] sur PL/K est isomorphe a E, c[RIGHT SINGLE QUOTATION MARK]est-a-dire encore que (PK 1                                       , O(1)[r]) est un
pre-modele de (PL/K , E). C[RIGHT SINGLE QUOTATION MARK]en est un modele si L = K(ai )1<=i<=n .

10.7.11. -- Soit K ,-> L une extension de corps classiques gradues ; soit n un entier
superieur ou egal au degre de transcendance de L sur K. Soit (E1 , . . . , En ) une famille
de fibres en droites sur PL/K . Il existe un pre-modele propre (V, (Fi )) de (PL/K , (Ei ))
(on peut meme si on le souhaite prendre pour V un modele propre de L sur K, mais
ce n[RIGHT SINGLE QUOTATION MARK]est pas necessaire pour ce qui suit). La formule de projection (lemme 10.4.8 d ))
entraine que l[RIGHT SINGLE QUOTATION MARK]entier defini par
                    (            R
                      [L : K(V )] V c1 (L1 ) * * * c1 (Ln ) si dim(V ) = n ;
                      0                                     si dim(V ) < n
ne depend pas du choix de V et des Fi . On le notera
                              Z
                                    c1 (E1 ) . . . c1 (En ).
                                           PL/K
                          '
10.7.12. -- Soit L une extension finie de L et soit pi le morphisme naturel de PL' /K
vers PL/K . Pour tout i, notons Ei' le fibre en droites pi * Ei sur PL' /K . Soit (V, (Fi )) un
modele propre de la famille des Li . C[RIGHT SINGLE QUOTATION MARK]est alors egalement un pre-modele propre de la
famille des Ei' , si bien qu[RIGHT SINGLE QUOTATION MARK]on a les egalites
                  Z                                         Z
                               '              '       '
                          c1 (E1 ) . . . c1 (En ) = [L : L]   c1 (F1 ) * * * c1 (Fn )
                     PL' /K                                      V
                                                                Z
                                                   = [L' : L]           c1 (E1 ) . . . c1 (En ).
                                                                 PL/K

10.7.13. -- Soit K ' une extension de K. Choisissons un modele propre (V, (Fi )) de la
famille des Ei . Le K ' -schema V ' = K '  tensor k L possede alors r composantes irreductibles
V1 , . . . , Vr (considerees comme munies de la structure reduite), qui sont toutes de dimen-
sion n (on se ramene a l[RIGHT SINGLE QUOTATION MARK]enonce non gradue analogue grace au lemme 10.3.5). Notons
Oj local de V ' au point generique de Vj , notons mj sa longueur et Lj son corps gradue
residuel, c[RIGHT SINGLE QUOTATION MARK]est-a-dire le corps gradue des fonctions de Vj . Remarquons que K '  tensor K L est
noetherien (en ecrivant L comme corps gradue des fractions de OV (U ) ou U est un ouvert
affine non vide de V , on voit en effet que K '  tensor K L est un localise d[RIGHT SINGLE QUOTATION MARK]un anneau gradue
noetherien) et que les Oj sont ses localises en ses differents ideaux premiers gradues

           [SECTION SIGN] 10.7. NOMBRES D[RIGHT SINGLE QUOTATION MARK]INTERSECTION SUR UN ESPACE DE ZARISKI-RIEMANN                                    293

minimaux. Les Oj , les mj et les Lj ne dependent donc que de K ' et pas du choix du
modele V . Pour tout (i, j), notons Eij l[RIGHT SINGLE QUOTATION MARK]image reciproque de Ei sur PLj /K ' .
   Pour tout i et tout j notons Fi' l[RIGHT SINGLE QUOTATION MARK]image reciproque de Fi sur V ' et Fij son image
reciproque sur Vj . Pour tout indice j la famille (Vj , (Fij )i ) est alors un modele propre de
la famille (Eij )i . Il vient

               Z                                     Z
                         c1 (E1 ) . . . c1 (En ) =           c1 (F1 ) * * * c1 (Fn )
                 PL/K                                V
                                                     Z
                                               =   c1 (F1' ) * * * c1 (Fn' )
                                                     V'
                                                 X     Z
                                               =   mj         c1 (F1j ) * * * c1 (Fnj )
                                                         j          Vj
                                                     X             Z
                                               =             mj                 c1 (E1j ) . . . c1 (Enj ).,
                                                         j          PL /K '
                                                                         j




ou la troisieme egalite resulte du lemme 10.4.8 c).
  Mentionnons un cas particulier important : si K '  tensor K L est integre, si l[RIGHT SINGLE QUOTATION MARK]on designe par
L' son corps gradue des fractions et par Ei' l[RIGHT SINGLE QUOTATION MARK]image reciproque de Ei sur PL' /K ' on a
alors
                     Z                                         Z
                               c1 (E1 ) . . . c1 (En ) =                     c1 (E1' ) . . . c1 (En' ).
                        PL/K                                     PL' /K '

10.7.14. -- Soit (E1 , . . . , Em ) une famille finie de fibres en droites sur PL/K et soit
(V, F1 , . . . , Fm ) un pre-modele de (PL/K , E1 , . . . , Em ) ; remarquons que V est propre.
Nous dirons que c1 (E1 ) . . . c1 (Em ) est numeriquement effectif si c1 (F1 ) . . . , c1 (Fm ) est
numeriquement effectif (10.4.11), ce qui ne depend pas du choix du pre-modele V . Nous
dirons qu[RIGHT SINGLE QUOTATION MARK]un fibre en droites E sur PL/K est numeriquement effectif si c1 (E) est nume-
riquement effectif. Tout fibre en droites sur PL/K engendre par ses sections globales est
numeriquement effectif.
   Soit L' une extension de L et soit K ' un sous-corps gradue de L' contenant K. Pour
tout i notons Ei' l[RIGHT SINGLE QUOTATION MARK]image reciproque de Ei sur PL' /K ' . Si c1 (E1 ) . . . c1 (Em ) est numeri-
quement effectif il en va de meme de c1 (E1' ) . . . c1 (Em   '
                                                                ) : cela resulte de la preservation
du caractere numeriquement effectif par morphismes propres et extension des scalaires
(voir 10.4.11).



Exemple 10.7.15. -- Soit D une droite vectorielle graduee sur L et soit t un element
homogene non nul de D. Soit (ai )1<=i<=n une famille finie d[RIGHT SINGLE QUOTATION MARK]elements de L x  algebriquement
independants sur K telle que K = K(ai )i et soit E le fibre en droites sur PL/K dont la
fibre generique est D et dont la metrique est telle que v(t-1 ) = max(1, v(a1 ), . . . , v(an ))
(lemme 10.6.18). L[RIGHT SINGLE QUOTATION MARK]exemples 10.7.10 exhibe un modele de (PL/K , E) de la forme
(PrK , O(1)[s]) (pour un certain polyrayon s et un certain
                                                         R s > 0). On deduit alors de
l[RIGHT SINGLE QUOTATION MARK]exemple 10.4.10 et de la remarque 10.4.9 on voit que PL/K c1 (E)n = 1.

294       CHAPITRE 10. ALGEBRE GRADUEE ET REDUCTION DES ANNEAUX NORMES

      [SECTION SIGN] 10.8. REDUCTION GRADUEE DES ANNEAUX ET MODULES
                              NORMES

  Nous allons ici presenter la theorie de la reduction graduee des anneaux et modules
normes telle qu[RIGHT SINGLE QUOTATION MARK]introduite par Temkin dans [127], en insistant sur le cas particulier
majeur des corps ultrametriques.
10.8.1. Reduction graduee. -- Soit M un groupe abelien muni d[RIGHT SINGLE QUOTATION MARK]une semi-norme
ultrametrique || * ||, c[RIGHT SINGLE QUOTATION MARK]est-a-dire d[RIGHT SINGLE QUOTATION MARK]une application || * || : M -> R+ telle que

               ||0|| = 0,     ||(-m)|| = ||m|| et ||m + n|| <= max(||m||, ||n||)

pour tout (m, n). On pose
                       M
                   M
                   f=     {m  in  M, ||m|| <= r}/{m  in  M, ||m|| < r}.
                            r>0


La loi interne de M induit une loi interne sur M   f qui en fait un groupe abelien gradue.
Dans ce contexte, si [MICRO SIGN] est un element homogene de M          f nous poserons ||[MICRO SIGN]|| = 0 si [MICRO SIGN]
est nul, et nous designerons sinon par ||[MICRO SIGN]|| le degre de l[RIGHT SINGLE QUOTATION MARK]unique composante homogene
contenant [MICRO SIGN] (si la semi-norme de M est notee |*|, nous ecrirons |[MICRO SIGN]| plutot que ||[MICRO SIGN]||). Si
m est un element de M et si r est un reel strictement positif tel que r >= m, on notera
me r l[RIGHT SINGLE QUOTATION MARK]image de a dans A er . Nous designerons par m e (sans indice) l[RIGHT SINGLE QUOTATION MARK]element homogene de
M egal a m
f           e ||m|| si ||m|| est non nul, et a 0 sinon. Si N est un sous-groupe dense de M
pour la topologie definie par || * || il resulte des definitions que N e =M f. .
   Soit A un anneau muni d[RIGHT SINGLE QUOTATION MARK]une semi-norme ultrametrique et sous-multiplicative || * ||.
(nous incluons dans la definition le fait que ||1|| <= 1, ce qui entraine que ||1|| = 1
ou ||1|| = 0, ce dernier cas n[RIGHT SINGLE QUOTATION MARK]etant possible que si || * || est identiquement nulle). La
multiplication de A induit alors une loi interne sur le groupe abelien gradue A       e qui en
fait un anneau gradue ; nous l[RIGHT SINGLE QUOTATION MARK]appellerons l[RIGHT SINGLE QUOTATION MARK]anneau gradue residuel de A.
   L[RIGHT SINGLE QUOTATION MARK]anneau gradue A  e est integre si et seulement si || * || est multiplicative (nous incluons
dans la definition le fait que ||1|| = 1).
   Remarquons que A   e1 est l[RIGHT SINGLE QUOTATION MARK]anneau residuel classique {a  in  A, ||a|| <= 1}/{a  in  A, ||a|| < 1}
de A, et que l[RIGHT SINGLE QUOTATION MARK]application a 7-> e a1 de A vers A e1 est l[RIGHT SINGLE QUOTATION MARK]application de reduction classique.
   Si k est un corps ultrametrique complet, son anneau gradue residuel e        k est un corps
gradue, que nous appellerons le corps gradue residuel de k.
   Soit M un A-module muni d[RIGHT SINGLE QUOTATION MARK]une semi-norme || * || telle que ||am|| <= ||a|| * ||m|| pour
tout (a, m). La loi externe de M induit alors une loi de A      e x M f vers M  f qui fait de ce
dernier un A-module
             e          gradue.
10.8.2. -- Soit A un anneau muni une semi-norme sous-multiplicative || * ||, et soit M
un A-module muni d[RIGHT SINGLE QUOTATION MARK]une semi-norme || * || telle que ||am|| <= ||a|| * ||m|| pour tout (a, m).
Soit m = (mi ) une famille d[RIGHT SINGLE QUOTATION MARK]elements de A tels que ||mi || =     / 0 pour tout i.
   Soit (ai )i in I une famille d[RIGHT SINGLE QUOTATION MARK]elements presque tous nuls de A. Supposons que les ||ai || sont
non toutes nulles, soit r le maximumP    des ||ai || * ||mi ||, et soit I ' l[RIGHT SINGLE QUOTATION MARK]ensemble des
                                                                                    P indices i
tels que ||ai || * ||mi || = r. Posons Pn = i in I ai mi . On a alors ||n|| <= r et n     r =
                                                                               eP      i in I ' a
                                                                                             ei m
                                                                                                fi .
On a ||n|| < r si et seulement si i in I ' e      e i = 0, et si ||n|| = r alors n
                                          ai m                              e = i in I ' e
                                                                                      ai m
                                                                                         fi .
   Il en resulte que les assertions suivantes sont equivalentes :

              [SECTION SIGN] 10.8. REDUCTION GRADUEE DES ANNEAUX ET MODULES NORMES                           295
               P
   (i) On a || ai mi || = supi ||ai || * ||mi || pour toute famille (ai ) de scalaires presque tous
nuls (notons que supi ||ai || * ||mi || = maxi ||ai || * ||mi || des que I est non vide) ;
   (ii) La famille (f
                    mi )i d[RIGHT SINGLE QUOTATION MARK]elements homogenes de M   f est libre sur A.
                                                                       e
   Et si elles sont satisfaites alors pour toute famille (ai ) d[RIGHT SINGLE QUOTATION MARK]elements de A presque tous
                                                   fi ou I ' designe l[RIGHT SINGLE QUOTATION MARK]ensemble des indices i
                     P                     P
nuls, la reduction ^    ai mi est egale a i in I ' e
                                                ai m
tels que ||ai || * ||mi || = supj ||aj || * ||mj ||.
10.8.3. -- Soit k un corps ultrametrique complet et soit A une k-algebre munie d[RIGHT SINGLE QUOTATION MARK]une
semi-norme multiplicative || * ||. Soit (ai ) une famille d[RIGHT SINGLE QUOTATION MARK]elements de A tels que les ||ai ||
soient tous non nuls. On a ||aJ || = ||a||J et donc a     fJ = e aJ pour tout multi-exposant J.
                                                    J
Il resulte alors de 10.8.2 applique a la famille (a )J du k-espace vectoriel semi-norme A
que les assertions suivantes sont equivalentes :
   (i) On a || lamdaJ aJ || = max|lamdaJ | * ||aJ || pour toute famille (lamdaJ ) de scalaires presque tous
               P
nuls ;
   (ii) La famille (aei )i d[RIGHT SINGLE QUOTATION MARK]elements homogenes de A
                                                   e est algebriquement independante sur e k.
   Et si elles sont satisfaites alors pour toute famille (lamdaJ ) de scalaires presque tous nuls,
                                             aJ ou E designe l[RIGHT SINGLE QUOTATION MARK]ensemble des multi-indices J
               P                     P
la reduction ^    lamdaJ aJ est egale a J in E lamda eJ e
tels que |lamdaJ | * ||a||J = maxN |lamdaN | * ||a||N .
Exemple 10.8.4. -- Soit k un corps ultrametrique et soit r = (r1 , . . . , rn ) un polyrayon
                       P T I= (T1 , . . . , Tn ) I et notons A la k-algebre k[T ] munie de la
de longueur finie. Posons
norme multiplicative      aI T 7-> max|aI |r . On deduit de l[RIGHT SINGLE QUOTATION MARK]equivalence (i) <==> (ii)
\end{lstlisting}

\subsubsection{CLD.txt: cached text lines
15287--15355}\label{cld.txt-cached-text-lines-1528715355}

\begin{lstlisting}
                           S         un G-recouvrement de X si tout point x de X possede
un voisinage de la forme i in J Xi ou J est fini et ou x  in  Xi pour tout i  in  J. Si l[RIGHT SINGLE QUOTATION MARK]on
ne s[RIGHT SINGLE QUOTATION MARK]astreint a recoller que des espaces strictement affinoides, on obtient la classe des
espaces strictement k-analytiques.
   Soit X un espace k-analytique. C[RIGHT SINGLE QUOTATION MARK]est en particulier un espace topologique dont tout
point possede une base de voisinages compacts (mais X n[RIGHT SINGLE QUOTATION MARK]est pas necessairement lo-
calement compact puisqu[RIGHT SINGLE QUOTATION MARK]un espace topologique localement compact est par definition
separe, et que ce n[RIGHT SINGLE QUOTATION MARK]est pas forcement le cas de X). L[RIGHT SINGLE QUOTATION MARK]espace X possede par ailleurs une
classe distinguee de partie compactes, ses domaines affinoides (si X = M (A ), les do-
maines affinoides de X sont precisement ceux decrits au 11.1.5 ci-dessus). Un domaine
analytique de X est une partie de X qui est G-recouverte par les domaines affinoides
qu[RIGHT SINGLE QUOTATION MARK]elle contient. Tout domaine affinoide de X est un domaine analytique de X ; plus
generalement, tout compact de X qui est une union finie de domaines analytiques de
X est un domaine analytique. Et tout ouvert de X (et notamment X lui-meme) est un
domaine analytique. En particulier, tout point x de X possede un voisinage compact
qui est une reunion finie de domaines affinoides contenant x. Les espaces definis dans [8]
sont ceux qui sont qualifies de bons dans [9], c[RIGHT SINGLE QUOTATION MARK]est-a-dire ceux dont tout point admet un
voisinage affinoide (et donc une base de voisinages affinoides.) Tout domaine analytique
de X est une partie localement fermee de X : l[RIGHT SINGLE QUOTATION MARK]intersection d[RIGHT SINGLE QUOTATION MARK]une famille finie de do-
maines analytiques est un domaine analytique ; la reunion d[RIGHT SINGLE QUOTATION MARK]une famille localement finie
de domaines analytiques fermes est un domaine analytique.
   Tout domaine analytique V de X herite d[RIGHT SINGLE QUOTATION MARK]une structure canonique d[RIGHT SINGLE QUOTATION MARK]espace analy-
tique, et d[RIGHT SINGLE QUOTATION MARK]un morphisme V ,-> X dont l[RIGHT SINGLE QUOTATION MARK]application sous-jacente est l[RIGHT SINGLE QUOTATION MARK]inclusion. Tout
morphisme f : Y -> X se factorisant ensemblistement par V se factorise en fait unique-
ment par la structure analytique de V .
11.1.8. La G-topologie. -- Soit X un espace k-analytique. En plus du site topolo-
gique sur X on definit un second site plus fin, note XG . Ses objets sont les domaines
analytiques de X, ses fleches sont les inclusions, et ses familles couvrantes sont les G-
recouvrements. La topologie de ce site est appelee la G-topologie de X.
  Comme tout ouvert de X est un domaine analytique et comme tout recouvrement
ouvert est un G-recouvrement, on dispose d[RIGHT SINGLE QUOTATION MARK]un morphisme de site canonique pi : XG ->
X. L[RIGHT SINGLE QUOTATION MARK]unite d[RIGHT SINGLE QUOTATION MARK]adjonction id -> pi* pi * est un isomorphisme ([9], p. 26).
  Le site XG est annele. Nous noterons OX le faisceau structural de XG . Sauf mention du
contraire, lorsque nous parlerons de faisceaux coherents sur X, il s[RIGHT SINGLE QUOTATION MARK]agira plus precisement
de OX -module coherents sur XG . Precisons que le faisceau OX ainsi que sa restriction
au site topologique usuel de X sont coherents ([49], lemme 0.1 ; pour la correction d[RIGHT SINGLE QUOTATION MARK]une
petite erreur dans la preuve, signalee par Jerome Poineau, voir [52], 1.1.1). Si X = M (A)
alors OX (X) = A.
  Tout morphisme f : Y -> X entre espaces analytiques est en particulier un morphisme
de sites anneles : si V est un domaine analytique de Y alors f -1 (V ) est un domaine
analytique de X et l[RIGHT SINGLE QUOTATION MARK]on dispose d[RIGHT SINGLE QUOTATION MARK]un morphisme f * : OX (V ) -> OY (f -1 (V )). Si X est
un espace k-affinoide M (A) et si Y est un espace affinoide M (B) sur une extension de

     [SECTION SIGN] 11.1. GEOMETRIE ANALYTIQUE : CONVENTIONS, NOTATIONS, PREMIERS RAPPELS               301

k, l[RIGHT SINGLE QUOTATION MARK]application f 7-> f * etablit une bijection entre l[RIGHT SINGLE QUOTATION MARK]ensemble des k-morphismes de Y
vers X et l[RIGHT SINGLE QUOTATION MARK]ensemble des morphismes bornes de k-algebres de A vers B.
   Si X = M (A) est un espace k-affinoide et si M est un A-module de type fini, le
faisceau Mf sur XG associe au prefaisceau V 7-> M  tensor A OX (V ) est coherent. Le foncteur
M 7-> M f etablit une equivalence entre la categorie des A-modules de type fini et celle
des faisceaux coherents sur X, de quasi-inverse F 7-> F (X) (cf. [9] ; ce resultat est du
a Kiehl dans le cas strictement affinoide).
   Lorsque X est bon, le foncteur pi* induit une equivalence de categories entre faisceaux
coherents sur XG et faisceaux coherents sur X, et qu[RIGHT SINGLE QUOTATION MARK]un faisceau coherent F sur XG est
localement libre si et seulement si c[RIGHT SINGLE QUOTATION MARK]est le cas de pi* F ([9], 1.3.4 ; notons que la definition
de faisceau coherent qu[RIGHT SINGLE QUOTATION MARK]il utilise dans cette proposition est ad hoc, mais qu[RIGHT SINGLE QUOTATION MARK]elle coincide
bien en fait avec la notion usuelle au vu de la coherence des faisceaux structuraux sur
les deux sites anneles en jeu).
11.1.9. Produits fibres, extension des scalaires. -- La categorie des espaces k-
analytiques admet des produits fibres. Si A -> B et A -> C sont deux morphismes
d[RIGHT SINGLE QUOTATION MARK]algebres k-affinoides, M (B)  x M (A) M (C) = M (B  tensor      b A C).
   Si L est un extension complete de k on dispose d[RIGHT SINGLE QUOTATION MARK]un foncteur d[RIGHT SINGLE QUOTATION MARK]extension des scalaires
X 7-> XL de la categorie des espaces k-analytiques vers celle des espaces L-analytiques.
Si X = M (A) alors XL = M (A tensor        b k L). En general on dispose d[RIGHT SINGLE QUOTATION MARK]un morphisme naturel
de sites anneles de XL vers X, qui est surjectif (cf. [48], 0.5).
   Supposons X bon, soit y un point de XL et soit x son image sur X. Le morphisme
naturel de Spec OXL ,y vers Spec OX,x est plat a fibres localement d[RIGHT SINGLE QUOTATION MARK]intersection complete,
et a fibres geometriquement regulieres si L est analytiquement separable sur k ; c[RIGHT SINGLE QUOTATION MARK]est le
theoreme [52], th. 2.6.5. La notion d[RIGHT SINGLE QUOTATION MARK]extension analytiquement separable est definie
\end{lstlisting}

\subsubsection{CLD.txt: cached text lines
15567--15604}\label{cld.txt-cached-text-lines-1556715604}

\begin{lstlisting}
   Un sous-espace analytique ferme de X est un ferme de Zariski Y de X muni d[RIGHT SINGLE QUOTATION MARK]une
structure d[RIGHT SINGLE QUOTATION MARK]espace k-analytique au-dessus de X induite selon le procede decrit ci-dessus
par un faisceau coherent d[RIGHT SINGLE QUOTATION MARK]ideaux I tel que Y = V (I ) ; notons que I se reconstitue
alors a partir de cette structure comme le noyau de OX -> i* OY , ou i : Y -> X est le
morphisme structural, qui est topologiquement l[RIGHT SINGLE QUOTATION MARK]inclusion.
   Tout ferme de Zariski Y possede au moins une structure de sous-espace analytique
ferme : celle donnee par le faisceau

                 J := V 7-> {f  in  O(V ), f (x) = 0 pour tout  in  Y  intersect  V

qui est bien coherent ([49], Prop. 4.2.1). C[RIGHT SINGLE QUOTATION MARK]est la plus petite structure de sous-espace
analytique ferme sur Y , au sens ou J est le plus grand faisceau coherent d[RIGHT SINGLE QUOTATION MARK]ideaux
sur X definissant le ferme Y ; c[RIGHT SINGLE QUOTATION MARK]est aussi la seule structure de sous-espace analytique
fermee sur Y dont le faisceau structural n[RIGHT SINGLE QUOTATION MARK]a pas de sections nilpotentes non nulles ;
on l[RIGHT SINGLE QUOTATION MARK]appelle la structure reduite sur Y . Si Y = X la structure reduite est associee au
faisceau d[RIGHT SINGLE QUOTATION MARK]ideaux dont les sections sont les fonctions localement nilpotentes ; le sous-
espace analytique ferme correspondant est note Xred . Le morphisme Xred ,-> X induit
un homeomorphisme entre Xred et X pour leur topologies usuelles, mais aussi pour leurs
topologies de Zariski.
   Si Y est un sous-espace analytique ferme de X, il resulte du theoreme de Gerritzen-
Grauert dans sa version due a Temkin ([128], Th. 3.1) que les domaines analytiques de
Y sont les sous-ensembles de Y qui admettent un G-recouvrement par des parties de la
forme V  intersect  Y ou V est un domaine affinoide de X.
   Si V est un domaine analytique de X la topologie de Zariski de V est plus fine que la
topologie de V induite par la topologie de Zariski de X, et en general strictement plus
fine.
   Si E est une partie de X et si (Xi ) est un G-recouvrement de X par des domaines
analytiques, alors E est un ferme (resp. un ouvert) de Zariski de X si et seulement si
E  intersect  Xi est un ferme (resp. un ouvert) de Zariski de Xi pour tout i ([49], Prop. 4.2.1).

306                  CHAPITRE 11. RAPPELS DE GEOMETRIE ANALYTIQUE

   Une immersion fermee (1) est un morphisme d[RIGHT SINGLE QUOTATION MARK]espaces k-analytiques f : Y -> X se
factorisant par un isomorphisme entre Y et un sous-espace analytique ferme de X. Une
immersion fermee preserve les corps residuels completes.
   Un morphisme d[RIGHT SINGLE QUOTATION MARK]espaces k-analytiques f : Y -> X est dit separe si sa diagonale
delta : Y -> Y  x X Y est une immersion fermee, et X est dit separe si le morphisme structural
de X vers M (k) est separe. La separation dans ce sens entraine la separation topolo-
\end{lstlisting}

\subsubsection{CLD.txt: cached text lines
16122--16143}\label{cld.txt-cached-text-lines-1612216143}

\begin{lstlisting}
Demonstration. -- Quitte a restreindre X, il suffit de prouver qu[RIGHT SINGLE QUOTATION MARK]il existe un tel voi-
sinage. Soit U le lieu de normalite de Xred ; c[RIGHT SINGLE QUOTATION MARK]est un ouvert de Zariski dense de X
contenant x. On peut donc supposer que Xred est normal. Tout voisinage analytique
(resp. voisinage affinoide) connexe de X est alors irreductible.

Lemme 11.5.8. -- Soit n et d des entiers. Soit f : Y -> X un morphisme fini entre
espaces k-analytiques purement de dimension n. Supposons que X est regulier et que Y
verifie la propriete (Sd ). Soit x un point de X tel que dk (x) >= n - d. Le morphisme f
est plat au dessus d[RIGHT SINGLE QUOTATION MARK]un voisinage de x.
Demonstration. -- Le morphisme f etant fini, il est topologiquement propre. Il suffit
donc pur prouver l[RIGHT SINGLE QUOTATION MARK]assertion de montrer que le lieu de platitude de f contient tous les
antecedents de x. Pour etablir ce fait, on peut remplacer X par n[RIGHT SINGLE QUOTATION MARK]importe quel de ses
domaines affinoides contenant x, et donc supposer X affinoide ; dans ce cas, Y l[RIGHT SINGLE QUOTATION MARK]est
egalement puisque f est fini. Soit y un antecedent de x. La finitude de f entraine que
les les dimensions centrales de (X, x) et (Y, y) sont egales ([52], 3.2.4) ; par ailleurs,
l[RIGHT SINGLE QUOTATION MARK]hypothese faite sur x assure que la dimension centrale de (X, x) vaut au moins n - d.
Comme on a d[RIGHT SINGLE QUOTATION MARK]autre part suppose que dimx (X) = dimy (Y ) = n, la formule (11.5.5.1)
ci-dessus implique que
                          dimKrull (OX,x ) = dimKrull (OY,y ) <= d.
Puisque OY,y satisfait la propriete (Sd ), il est de Cohen-Macaulay. La OX,x -algebre fi-
nie locale OY,y est de Cohen-Macaulay, et de meme dimension de Krull que l[RIGHT SINGLE QUOTATION MARK]anneau
local regulier OX,x . En consequence OY,y est une OX,x -algebre plate ([71], EGA IV2 ,
\end{lstlisting}

\subsubsection{CLD.txt: cached text lines
3750--3920}\label{cld.txt-cached-text-lines-37503920}

\begin{lstlisting}
Alors l[RIGHT SINGLE QUOTATION MARK]application p|T : T -> sigma est paralineaire.

Demonstration. -- L[RIGHT SINGLE QUOTATION MARK]assertion est G-locale sur X et Y . Par definition d[RIGHT SINGLE QUOTATION MARK]une partie pa-
ralineaire, on peut donc supposer que sigma est fidelement paralineaires. Soit f  in  lamda(sigma).
Comme p est un morphisme d[RIGHT SINGLE QUOTATION MARK]espaces G-tropicaux, f [WHITE BULLET] p est une application paralineaire
sur Y . Par suite, f [WHITE BULLET] p|T est paralineaire (lemme 3.2.8). Comme f identifie sigma a une
partie paralineaire d[RIGHT SINGLE QUOTATION MARK]un espace numerique, cela demontre que p induit une application
paralineaire de T vers sigma.

                 [SECTION SIGN] 3.3. FORMES DE TYPE (p, q) SUR UN ESPACE G-TROPICAL                    79

     [SECTION SIGN] 3.3. FORMES DE TYPE (p, q) SUR UN ESPACE G-TROPICAL

  Dans ce paragraphe, on fixe une classe de coefficients G .
3.3.1. -- Soit X un espace G-tropical.
    Les cartes G-tropicales sur les domaines de X forment une categorie essentiellement
petite, un morphisme d[RIGHT SINGLE QUOTATION MARK]une carte G-tropicale (f : X -> E, P ) dans une autre (f ' : X ->
E ' , P ' ) etant une application affine q : E -> E ' tel que f ' = q [WHITE BULLET] f et q(P )  subset  P ' .
    Cette categorie a un objet final pour lequel E = {0} et P est un point. Elle est
filtrante : si (f : X -> E, P ) et (f ' : X -> E ' , P ' ) sont deux cartes G-tropicales sur X,
alors (f  x  f ' : X -> E  x  E ' , P  x  P ' ) est une carte G-tropicale sur X qui domine les deux
cartes G-tropicales donnees.

Remarque 3.3.2. -- Soit X un espace G-tropical compact et soit (f : X -> E, P ) une
G-carte tropicale sur X.
  a) L[RIGHT SINGLE QUOTATION MARK]ensemble f (X) est un sous-espace paralineaire compact de E contenu dans P
(theoreme. 12.2.4) et la carte donnee se factorise par (f : X -> E, f (X)). On dira parfois,
par abus, que f est une carte G-tropicale sur X.
  b) Pour toute decomposition cellulaire (Pi ) de f (X), la famille (f -1 (Pi )) est un G-
recouvrement fini de X par des domaines compacts dont l[RIGHT SINGLE QUOTATION MARK]image par f est une cellule.

3.3.3. -- Pour tout domaine V de X et tout couple (p, q) d[RIGHT SINGLE QUOTATION MARK]entiers naturels, notons
Gpfp,q (V ) la limite inductive des R-espaces vectoriels GPp,q (P ), lorsque (f : V -> E, P )
parcourt l[RIGHT SINGLE QUOTATION MARK]ensemble des cartes G-tropicales sur V .
    Lorsque V parcourt l[RIGHT SINGLE QUOTATION MARK]ensemble des domaines de X, les espaces vectoriels Gpfp,q (V )
definissent un prefaisceau sur le site XG .
    Soit V un domaine de X et soit (f : V -> E, P ) une carte G-tropicale sur V . Toute
forme alpha appartenant a GPp,q (P ) definit un element de Gpfp,q (V ) qui est note f * alpha.
    Par definition, toute section de Gpfp,q (V ) est de cette forme ; on peut meme supposer
que P est ouvert. Plus precisement, soit (f : V -> E, P ) une carte G-tropicale sur V ,
soit alpha  in  GPp,q (P ) ; il existe un ouvert U de E dont P est une partie fermee ; d[RIGHT SINGLE QUOTATION MARK]apres le
lemme 1.9.4, il existe beta  in  GEp,q (U ) telle que alpha = beta|P ; on a alors f * alpha = f * beta.

Definition 3.3.4. -- On note GXp,qG le faisceau sur XG associe au prefaisceau Gpfp,q . Ses
sections sont appelees G-formes de type (p, q) a coefficients dans G .
  Lorsque la classe G est la classe A des fonctions lisses, on parle de G-forme lisse de
type (p, q).

   Si X est l[RIGHT SINGLE QUOTATION MARK]espace G-tropical associe a un espace paralineaire, on retrouve ainsi le
faisceau des G-formes de type (p, q) a coefficients dans G .
3.3.5. -- On definit le support d[RIGHT SINGLE QUOTATION MARK]une G-forme de type (p, q) sur X comme le comple-
mentaire du plus grand ouvert de X en restriction auquel cette forme est nulle.
  Si (Xi ) est un G-recouvrement de X, le support d[RIGHT SINGLE QUOTATION MARK]une G-forme omega coincide avec la
reunion des supports des G-formes omega|Xi . Il est tout d[RIGHT SINGLE QUOTATION MARK]abord evident que supp(omega|Xi )  subset 
supp(omega)  intersect  Xi pour tout i. Inversement, soit x un point de X qui n[RIGHT SINGLE QUOTATION MARK]appartient au support
d[RIGHT SINGLE QUOTATION MARK]aucune des formes omega|Xi ; demontrons T que x n[RIGHT SINGLE QUOTATION MARK]appartient
                                                   S       pas au support de omega. Soit I
un ensemble fini d[RIGHT SINGLE QUOTATION MARK]indices tel que x  in  i in I Xi et i in I Xi soit un voisinage de x. Pour
tout i  in  I, soit Vi un voisinage de x dans Xi qui est un domaine compact et sur lequel

80                            CHAPITRE 3. ESPACES G-TROPICAUX
                    S
omega|Vi = 0 ; soit V = i in I Vi ; c[RIGHT SINGLE QUOTATION MARK]est un domaine compact de X et un voisinage de x. Par
construction, la G-forme omega|V est G-localement nulle, donc nulle.

Proposition 3.3.6. -- Soit X un espace G-tropical et soit omega une G-forme de type (p, q)
sur X. Le support de omega est contenu dans l[RIGHT SINGLE QUOTATION MARK]ensemble des points x  in  X tels que dtrop (x) >=
max(p, q).

Demonstration. -- Soit x  in  X un point tel que dk (x) < max(p, q) ; demontrons que omega
est nulle au voisinage de x. Par hypothese, il existe des domaines compacts V1 , . . . , Vn
de X dont la reunion est un voisinage V de x dans X, des cartes G-tropicales fi : Vi -> Ei
et des formes alphai sur fi,trop (Vi ) tels que omega|Vi = fi* alphai pour tout i. La dimension tropicale
de fi en x est majoree par dtrop (x) ; quitte a diminuer les Vi , on peut supposer que
dim(fi (Vi )) <= dtrop (x) ; alors, alphai = 0. On a donc omega|V = 0.

Corollaire 3.3.7. -- Si dtrop (X) < max(p, q), alors les faisceaux FXp,q et FXp,qG sont
nuls.

3.3.8. -- Soit omega une G-forme (resp. une forme) de type (p, q) a coefficients dans G
sur X. Si V est un domaine de X, une G-tropicalisation de omega sur V est la donnee d[RIGHT SINGLE QUOTATION MARK]une
carte G-tropicale (f : V -> E, P ) et d[RIGHT SINGLE QUOTATION MARK]une forme alpha  in  GPp,q (P ) telle que omega = f * alpha.
  On dira que omega est G-tropicale si elle admet une G-tropicalisation sur X.
  Si V est un domaine de X tel que omega possede une G-tropicalisation sur V , nous dirons
aussi que V G-tropicalise omega.
  Soit V un domaine compact de X qui G-tropicalise omega, et soit (f V -> E, P, alpha) une
G-tropicalisation de omega|V . Comme f (V ) est un sous-espace paralineaire de E, (f : V ->
E, f (V ), alpha|f (V ) ) est encore une G-tropicalisation de omega|V ; une telle G-tropicalisation sera
simplement notee (f : V -> E, alpha) : lorsque V est compact, l[RIGHT SINGLE QUOTATION MARK]omission du sous-espace
paralineaire P signifie qu[RIGHT SINGLE QUOTATION MARK]il est egal a f (V ).

Remarque 3.3.9. -- Soit H une classe de coefficients contenant G . Observons que
les G-formes a coefficients dans G s[RIGHT SINGLE QUOTATION MARK]identifient a des G-formes a coefficients dans H .
   En effet, le G-prefaisceau Gpfp,q sur XG s[RIGHT SINGLE QUOTATION MARK]identifie a un sous-G-prefaisceau de Hpfp,q , de
sorte que le G-faisceau GXp,qG sur X s[RIGHT SINGLE QUOTATION MARK]identifie alors a un sous-G-faisceau de HXp,qG
                                                                                       . Ces
identifications sont compatibles avec les constructions faites dans ce paragraphe.

3.3.10. -- Le faisceau GXG := GX0,0G est un faisceau en R-algebres et les faisceaux GXp,qG
sont des GXG -modules. Leur somme directe GX*,*G = p,q GXp,qG est un faisceau en algebres
                                                   L

bigraduees commutatives : si omega est de type (p, q) et omega ' de type (p' , q ' ), alors
                                                    '     '
                                omega '  wedge  omega = (-1)(p+p )(q+q ) omega  wedge  omega ' .
3.3.11. -- Le G-prefaisceau Gpf*,* est naturellement muni d[RIGHT SINGLE QUOTATION MARK]un operateur J qui s[RIGHT SINGLE QUOTATION MARK]etend
au faisceau GX*,*G . Ce sont des morphismes d[RIGHT SINGLE QUOTATION MARK]algebres. On dit qu[RIGHT SINGLE QUOTATION MARK]une section omega de GXp,qG
est symetrique si p = q et si Jomega = omega.
    Soit H une classe de coefficients telles que les derivees partielles des fonctions de G
existent et appartiennent a H . On dispose alors d[RIGHT SINGLE QUOTATION MARK]operateurs differentiels d' et d'' de
Gpf*,* dans Hpf*,* ; ils induisent des operateurs differentiels de GX*,*G dans HX*,*
                                                                                  G
                                                                                     , caracterises
par les relations
                               d' f * alpha = f * d' alpha, d'' f * alpha = f * d'' alpha,

                                 [SECTION SIGN] 3.4. NULLITE DES G-FORMES                             81

pour toute carte G-tropicale (f : U -> E, P ) et toute forme alpha de type (p, q) a coeffi-
cients dans G sur P . Les differentielles verifient la regle de Leibniz : si omega et omega ' sont
respectivement de type (p, q) et (p' , q ' ), on a
(3.3.11.1)               d' (omega  wedge  omega ' ) = d' omega  wedge  omega ' + (-1)p+q omega  wedge  d' omega ' ,
et de meme pour d'' (omega  wedge  omega ' ).
   Si G  subset  C 2 , ces operateurs verifient d' d' = d'' d'' = 0.
3.3.12. -- Cette construction est fonctorielle. Plus precisement, si f : Y -> X est un
morphisme d[RIGHT SINGLE QUOTATION MARK]espaces G-tropicaux, on dispose de morphismes de faisceaux
                                     f * : GX*,*G -> f* (GY*,*
                                                           G
                                                              )
sur XG qui sont des morphismes d[RIGHT SINGLE QUOTATION MARK]anneaux gradues, commutent aux derivations d' et
d'' (lorsqu[RIGHT SINGLE QUOTATION MARK]elles existent) et a l[RIGHT SINGLE QUOTATION MARK]operateur J.
   Lorsque le contexte rend evident le morphisme f , on se permettra parfois de noter
omega|Y l[RIGHT SINGLE QUOTATION MARK]image dans GYp,q G
                         (Y ) d[RIGHT SINGLE QUOTATION MARK]une G-forme omega  in  GXp,qG .
3.3.13. -- Soit (p, q) un couple d[RIGHT SINGLE QUOTATION MARK]entiers naturels. Soit sigma une partie paralineaire de X
et soit j l[RIGHT SINGLE QUOTATION MARK]inclusion de sigma dans X. C[RIGHT SINGLE QUOTATION MARK]est un morphisme d[RIGHT SINGLE QUOTATION MARK]espaces G-tropicaux. On dispose
donc d[RIGHT SINGLE QUOTATION MARK]un morphisme de faisceaux de FXp,qG dans j* Fsigmap,qG . On notera omega 7-> omega|sigma l[RIGHT SINGLE QUOTATION MARK]image
d[RIGHT SINGLE QUOTATION MARK]une forme omega et on dira que c[RIGHT SINGLE QUOTATION MARK]en est la restriction a sigma.
   Concretement, cette restriction se calcule de la facon suivante.
   Soit V un domaine de X, soit (f : V -> E, P ) une carte G-tropicale sur V , et soit
alpha  in  FPp,q (P ). Alors sigma  intersect  V est une partie paralineaire de V , l[RIGHT SINGLE QUOTATION MARK]application f : sigma  intersect  V -> E
est paralineaire et son image est contenue dans P , l[RIGHT SINGLE QUOTATION MARK]image reciproque f * alpha de alpha est
une forme de type (p, q) sur l[RIGHT SINGLE QUOTATION MARK]espace paralineaire sigma  intersect  V . Par definition de l[RIGHT SINGLE QUOTATION MARK]ensemble
Fpfp,q (V ) comme limite inductive, la forme f * alpha ne depend que de la classe de f * alpha dans
cet ensemble. Cette construction definit un morphisme de prefaisceaux de FX,pf     p,q
                                                                                       dans
j* Fsigmap,q .
   Par passage aux faisceaux associes, on en deduit le morphisme de restriction de FXp,qG
dans j* Fsigmap,qG .
3.3.14. -- Soit f : Y -> X un morphisme d[RIGHT SINGLE QUOTATION MARK]espaces G-tropicaux, soit T une partie
paralineaires de Y et soit sigma une partie paralineaire de X telle que f (sigma)  subset  T. Pour toute
forme omega  in  FXp,qG (X), on a (f * omega)|T = (f |T )* (omega|sigma ).
   Nous appliquerons notamment cette remarque lorsque Y est un domaine de X et
T = sigma  intersect  Y , ou bien lorsque Y = X et T est une partie polyedrale de sigma.



                          [SECTION SIGN] 3.4. NULLITE DES G-FORMES

Proposition 3.4.1. -- Soit X un espace G-tropical compact, soit f : X -> E une carte
G-tropicale, soit P un sous-ensemble paralineaire de E contenant f (X) et soit alpha  in 
GPp,q (P ). Les assertions suivantes sont equivalentes :
   a) La restriction de alpha au sous-espace paralineaire f (X) est nulle ;
   b) La section globale f * alpha de Gpfp,q est nulle.
Demonstration. -- L[RIGHT SINGLE QUOTATION MARK]hypothese (1) entraine que l[RIGHT SINGLE QUOTATION MARK]image de alpha dans la carte tropicale
(f, f (X)) est nulle, de sorte que (2) vaut. Supposons que (2) est verifiee et prouvons (1).

82                          CHAPITRE 3. ESPACES G-TROPICAUX

   On peut remplacer P par f (X). Par definition, il existe une carte G-tropicale g : X ->
F , une morphisme affine p : F -> E tel que p [WHITE BULLET] g = f , et un sous-espace paralineaire Q
\end{lstlisting}

\subsubsection{CLD.txt: cached text lines
5289--5314}\label{cld.txt-cached-text-lines-52895314}

\begin{lstlisting}
que les classes de omega|sigmai et de omegai coincident dans Vpfn (Xi ). Notons f l[RIGHT SINGLE QUOTATION MARK]injection de sigmai
dans sigma intersect Xi ; il suffit de demontrer l[RIGHT SINGLE QUOTATION MARK]egalite omega|sigma intersect Xi = f* omegai de n-formes-volume sur sigma intersect Xi .
Comme le prefaisceau Vpfn est separe, il suffit de verifier cette egalite apres restriction
a sigmaj  intersect  Xi , pour tout j.
   Notons alpha la n-forme-volume omega|sigmai  intersect sigmaj . Comme les classes de omegai et omegaj dans V n (Xi  intersect  Xj )
coincide, la forme-volume omegai |sigmai  intersect Xj est le prolongement par zero de omegai |sigmai  intersect sigmaj = alpha, et
omegaj |sigmaj  intersect Xi est le prolongement par zero de alpha. Ainsi, omega|sigmaj  intersect Xi = omegaj |sigmaj  intersect Xi est le prolongement
par zero de omegaj |sigmai  intersect sigmaj = alpha. D[RIGHT SINGLE QUOTATION MARK]autre part, (f* omegai )|sigmaj  intersect Xi est le prolongement par zero de
omegai |sigmai  intersect sigmaj = alpha. Cela demontre l[RIGHT SINGLE QUOTATION MARK]egalite requise et termine la demonstration du lemme.

4.3.9. -- Soit f : X -> Y une immersion par morceaux d[RIGHT SINGLE QUOTATION MARK]espaces lineaires par mor-
ceaux. On a alors un morphisme d[RIGHT SINGLE QUOTATION MARK]image reciproque f * : VYnG -> f* VXnG , deduit du mor-
phisme defini dans le paragraphe 4.3.5 lorsque dim(X) <= n. Si, de plus, f est compacte,
on a un morphisme d[RIGHT SINGLE QUOTATION MARK]image directe f* : f* VXnG -> VYnG , deduit du morphisme defini dans
le paragraphe 4.3.6 lorsque dim(X) <= n. Dans les deux cas le point cle est que ces
operations commutent au prolongement par zero.




           [SECTION SIGN] 4.4. MESURE ASSOCIEE A UNE FORME-VOLUME

4.4.1. -- Soit X un espace paralineaire de niveau n. Soit omega une forme-volume dimen-
sionnelle a coefficients boreliens sur X. Le but de ce paragraphe est de lui associer une
mesure (positive) de Borel reguliere omegaBor sur X et, si cette mesure est localement finie,
\end{lstlisting}

\subsubsection{CLD.txt: cached text lines
5449--5483}\label{cld.txt-cached-text-lines-54495483}

\begin{lstlisting}
   En general, le G-faisceau des calibrages sur X est G-localement libre de rang 1
sur GXG /In GXG . On dira qu[RIGHT SINGLE QUOTATION MARK]un calibrage est G-localement constant (resp. posi-
tif, resp. ...) s[RIGHT SINGLE QUOTATION MARK]il est G-localement multiple d[RIGHT SINGLE QUOTATION MARK]un calibrage standard par une section
G-localement constante (resp. positive, resp. ...) de GXG /In GXG .

Exemple 4.5.2. -- Soit X un espace paralineaire de niveau n, paracompact.
   Soit C une decomposition cellulaire de X ; pour toute cellule C appartenant a C , soit
fC un isomorphisme de C sur une cellule d[RIGHT SINGLE QUOTATION MARK]un espace affine VC . Pour toute cellule C  in  Cn ,
                                                     ------>
choisissons une orientation epsilonC de l[RIGHT SINGLE QUOTATION MARK]espace vectoriel <fC (C)> et une fonction [MICRO SIGN]C de classe G
                             ------>
de fC (C) dans l[RIGHT SINGLE QUOTATION MARK]espace |lamdan <fC (C)>| des vecteurs-volumes sur cet espace vectoriel ; soit
                                                          ------>
alors [MICRO SIGN]'C l[RIGHT SINGLE QUOTATION MARK]unique fonction de fC (C) a valeurs dans lamdan <fC (C)> tel que [MICRO SIGN]C soit la classe
de ([MICRO SIGN]'C , epsilonC ).
   Il existe alors un unique calibrage [MICRO SIGN] de X tel que pour toute cellule C de dimen-
                                    ------>
sion n et tout element omega  in  An,0 (<fC (C)>), [MICRO SIGN](fC* [omega, epsilonC ]) soit la classe dans G /In G (C)
                                                                                  ------>
de la fonction deduite de (omega, [MICRO SIGN]'C ) par l[RIGHT SINGLE QUOTATION MARK]accouplement de dualite entre An,0 (<fC (C)>) et
   ------>
lamdan <fC (C)> et composition par fC . On resume cette propriete en ecrivant [MICRO SIGN]|C = fC* [MICRO SIGN]C .
   Si les fonctions [MICRO SIGN]C sont constantes (resp. positives, resp. ...), alors le calibrage [MICRO SIGN] est
G-localement constant (resp. positif, resp. ...).
   Pour tout calibrage [MICRO SIGN] sur X, il existe une decomposition cellulaire C , une famille (fC )
d[RIGHT SINGLE QUOTATION MARK]isomorphismes et une famille ([MICRO SIGN]C ) qui definissent [MICRO SIGN] comme ci-dessus.

                     [SECTION SIGN] 4.5. CALIBRAGES ET INTEGRALES DES (n, n)-FORMES                          109

  Si le calibrage [MICRO SIGN] est G-localement constant (resp. positif), on peut en outre exiger que
pour toute cellule C de dimension n, la fonction [MICRO SIGN]C soit constante (resp. positive) sur C.
  Si X est muni d[RIGHT SINGLE QUOTATION MARK]une immersion par morceaux f : X -> V vers un espace affine, on
peut aussi exiger que fC = f |C pour toute cellule C.
  En particulier, si X est un sous-espace paralineaire d[RIGHT SINGLE QUOTATION MARK]un espace affine V , il existe une
decomposition cellulaire forte C de X telle que pour toute cellule C  in  C , le calibrage [MICRO SIGN]|C
\end{lstlisting}

\subsubsection{CLD.txt: cached text lines
6308--6332}\label{cld.txt-cached-text-lines-63086332}

\begin{lstlisting}
Exemple 5.2.8. -- Soit k un corps complet pour une valeur absolue non archimedienne
et soit X un espace k-analytique, topologiquement separe, et purement de dimension n.
Dans la seconde partie de ce livre, nous allons le munir d[RIGHT SINGLE QUOTATION MARK]une structure d[RIGHT SINGLE QUOTATION MARK]espace G-
tropical vertebre. Dans le but d[RIGHT SINGLE QUOTATION MARK]eclairer les constructions et definitions de la presente
partie, indiquons ici quelques aspects de cette structure a venir.
   Les domaines de X, comme espace G-tropical, sont ses domaines analytiques.

126                    CHAPITRE 5. ESPACES G-TROPICAUX CALIBRES

   Pour definir ses fonctions paralineaires et ses fonctions affines, on procedera a rebours
de l[RIGHT SINGLE QUOTATION MARK]axiomatisation generale en commencant par definir les fonctions affines sur l[RIGHT SINGLE QUOTATION MARK]espace
G-tropical X : ce sont celles qui sont localement combinaisons Q-lineaires de fonctions
de la forme log|f |, ou f est une fonction analytique inversible sur un domaine analytique
de X. Les fonctions paralineaires sur l[RIGHT SINGLE QUOTATION MARK]espace G-tropical X sont alors les fonctions qui
sont G-localement affines.
   Parmi toutes les cartes G-tropicales sur X, il suffira de considerer celles de la forme
ftrop : V -> Rn ou V est un domaine analytique de X, f : V -> Gnm est un mor-
phisme d[RIGHT SINGLE QUOTATION MARK]espaces analytiques (nous appellerons moment de tels morphismes), et ftrop =
(log|f1 |, . . . , log|fn |).
   Les parties squelettiques elementaires de X sont les parties de la forme sigmaf = f -1 (Sn ),
ou V est un domaine de X et f : V -> Gnm est un moment et Sn est le squelette standard
                                                                  (n)
de Gnm . Alors, l[RIGHT SINGLE QUOTATION MARK]ensemble (sigmaf )(n) n[RIGHT SINGLE QUOTATION MARK]est autre que l[RIGHT SINGLE QUOTATION MARK]ensemble sigmaftrop .
   Si X est un domaine analytique d[RIGHT SINGLE QUOTATION MARK]un bon espace, en particulier, si X est sans bord,
nous demontrerons que l[RIGHT SINGLE QUOTATION MARK]espace G-tropical associe a X est decent (definition 6.3.1 et
\end{lstlisting}

\subsubsection{CLD.txt: cached text lines
6400--6414}\label{cld.txt-cached-text-lines-64006414}

\begin{lstlisting}
entiers tels que max(p, q) = n et soit sigma une partie squelettique fermee de X qui sup-
porte fortement omega. Alors, sigma(n) est une partie squelettique fermee de X qui supporte
fortement omega.

Demonstration. -- Comme sigma(n) est une partie paralineaire fermee de sigma, c[RIGHT SINGLE QUOTATION MARK]est une partie
squelettique fermee de X.
  Soit T une partie squelettique fermee de X qui contient sigma(n) ; demontrons que omega|T
est le prolongement par zero de omega|sigma(n) . Quitte a remplacer T par T  union  sigma, on suppose
que T contient sigma. Quitte a raisonner G-localement, on peut supposer que T possede
une decomposition cellulaire adaptee a sigma, et il s[RIGHT SINGLE QUOTATION MARK]agit de demontrer la restriction de omega a
toute cellule fermee qui n[RIGHT SINGLE QUOTATION MARK]est pas contenue dans sigma(n) est nulle. C[RIGHT SINGLE QUOTATION MARK]est en effet le cas d[RIGHT SINGLE QUOTATION MARK]une
cellule qui n[RIGHT SINGLE QUOTATION MARK]est pas dans sigma, par definition d[RIGHT SINGLE QUOTATION MARK]un support fort, et d[RIGHT SINGLE QUOTATION MARK]une cellule qui n[RIGHT SINGLE QUOTATION MARK]est
pas dimension n, car la restriction d[RIGHT SINGLE QUOTATION MARK]une (p, q)-forme a une telle cellule est nulle.
  Prouvons maintenant que le support de omega est contenu dans sigma(n) . Soit x un point
de sigma(n) ; prouvons que omega est nulle au voisinage de x. L[RIGHT SINGLE QUOTATION MARK]assertion est G-locale au voisinage
\end{lstlisting}

\subsubsection{CLD.txt: cached text lines
6753--6812}\label{cld.txt-cached-text-lines-67536812}

\begin{lstlisting}
de X ; pour tout i, soit lamdai  in  Vpfp (Xi ). Si lamdai |Xi  intersect Xj = lamdaj |Xi  intersect Xj pour tout (i, j) il existe une
unique forme-volume lamda  in  Vpfp (X) telle que lamda|Xi = lamdai pour tout i.

Demonstration. -- L[RIGHT SINGLE QUOTATION MARK]unicite de lamda provient du lemme 5.6.2 ci-dessus. En vertu de cette
assertion d[RIGHT SINGLE QUOTATION MARK]unicite, l[RIGHT SINGLE QUOTATION MARK]existence peut etre prouvee apres raffinement du recouvrement (Yi ).
Puisque X est paracompact, on peut donc supposer que le recouvrement (Yi ) est loca-
lement fini. Pour tout i, choisissons une partie squelettique fermee sigmai de Yi telle que lamdai
provienne d[RIGHT SINGLE QUOTATION MARK]un element lamdasigmai  in  V p (sigmai ). Comme le G-recouvrement (Yi ) est localement fini,
la reunion sigma des sigmai est une partie squelettique fermee de X (proposition 5.1.5) dont la
famille (sigmai ) est un G-recouvrement. Les lamdasigmai se recollent alors en une section lamdasigma  in  V p (sigma).
Notons lamda la classe de lamdasigma dans V p (X) ; alors lamda|Xi = lamdai quel que soit i.

134                   CHAPITRE 5. ESPACES G-TROPICAUX CALIBRES

Corollaire 5.6.5. -- Soit Y un domaine paracompact de X et soit lamda  in  VXFG,p (Y ) une
forme-volume p-dimensionnelle sur Y . Il existe une partie squelettique fermee sigma de Y
de dimension <= p et une p-forme-volume lamdasigma sur sigma qui definit lamda.

  Il est vraisemblable que le corollaire ne vaille pas sans l[RIGHT SINGLE QUOTATION MARK]hypothese que Y est paracom-
pact.
5.6.6. -- Soit Y un domaine ferme de X ; notons i l[RIGHT SINGLE QUOTATION MARK]immersion canonique de Y dans X.
Comme toute partie squelettique fermee de Y est une partie squelettique fermee de X,
le prolongement par zero definit un homomorphisme de G-prefaisceaux de i* VYpG ,pf
dans VXpG ,pf . On en deduit un morphisme de G-faisceaux de i* VYpG dans VXpG . Si lamda est une
forme-volume p-dimensionnelle sur Y , on notera i* lamda la forme-volume p-dimensionnelle
correspondante sur X.
5.6.7. -- On dit qu[RIGHT SINGLE QUOTATION MARK]une forme-volume p-dimensionnelle lamda sur X est positive s[RIGHT SINGLE QUOTATION MARK]il existe
un G-recouvrement (Xi ) de X tel que, pour tout i, lamda soit representee par une formule-
volume positive sur une partie squelettique fermee de Xi .
  Dans ce cas, toute forme-volume sur une partie squelettie fermee d[RIGHT SINGLE QUOTATION MARK]un domaine Y
de X induisant la forme-volume omega|Y est elle-meme positive.




         [SECTION SIGN] 5.7. MESURES ASSOCIEES A UNE FORME-VOLUME
                          p-DIMENSIONNELLE

   Soit X un espace G-tropical vertebre de niveau n et soit p un entier naturel tel que
p <= n. Nous allons associer a toute forme-volume lamda, p-dimensionnelle, a coefficients
boreliens, sur XG une mesure de Borel reguliere lamdaBor sur X et, lorsque cette mesure est
localement finie, une mesure de Radon lamdaRad sur X telle que lamdaBor = |lamdaRad | qui est positive
lorsque lamda est une forme-volume positive.
   Soit G une classe de coefficients qui est contenue dans B.
5.7.1. -- Soit sigma et sigma' des parties squelettiques fermees de X telles que sigma  subset  sigma' . Soit lamda
appartenant a V G ,p (sigma) et soit lamda' son prolongement par zero a sigma' . Notons [MICRO SIGN] et [MICRO SIGN]' les
mesures de Borel sur sigma et sigma' associees a lamda et lamda' ; lorsqu[RIGHT SINGLE QUOTATION MARK]on les considere comme mesures
de Borel sur X, elles sont egales.
   Les mesure de Borel [MICRO SIGN] et [MICRO SIGN]' sont simultanement localement finies ou non. Lorsqu[RIGHT SINGLE QUOTATION MARK]elles
le sont, les formes-volume lamda et lamda' donnent lieu a des mesures de Radon [MICRO SIGN] et [MICRO SIGN]' sur sigma
et sigma' dont [MICRO SIGN] et [MICRO SIGN]' sont les valeurs absolues. On a de plus [MICRO SIGN] = [MICRO SIGN]' lorsque ces mesures de
Radon sont identifiees a des mesures de Radon sur X.
   Par consequent, il existe un unique morphisme de faisceaux lamda 7-> lamdaBor sur XG qui, pour
tout domaine Y de X, toute partie squelettique fermee sigma de Y et toute forme-volume lamdasigma
sur sigma, p-dimensionnelle et a coefficients dans G , applique la classe de lamdasigma dans V G ,p (Y )
sur la mesure de Borel associee a lamdasigma , consideree comme une mesure de Borel sur Y .
   Notons VXGG,p,pf le sous-faisceau de VXGG,p constitue des formes-volume lamda telles que lamdaBor
soit localement finie. Il existe un unique morphisme lamda 7-> lamdaRad de ce faisceau a valeurs
dans le faisceau des mesures de Radon sur X qui, pour tout domaine Y de X, toute
partie squelettique fermee sigma de Y et toute forme-volume lamdasigma sur sigma dont la mesure de
\end{lstlisting}

\subsubsection{CLD.txt: cached text lines
6934--6977}\label{cld.txt-cached-text-lines-69346977}

\begin{lstlisting}
   Les G-formes localement integrables constituent un sous-espace vectoriel de F n,n (XG )
contenant les G-formes lisses. Les G-formes integrables en sont un sous-espace vectoriel
qui contient les G-formes lisses a support compact.
5.8.6. -- Soit u une fonction localement integrable sur X et soit omega une G-forme lisse
de type (n, n) sur X. Alors, u est localement integrable par rapport a la mesure omegaBor ,
de sorte que uomegaRad est une mesure de Radon sur X.
   L[RIGHT SINGLE QUOTATION MARK]application (u, omega) 7-> uomegaRad est bilineaire ; elle induit une application lineaire de
Lloc
   1
     (XG )  tensor  A n,n (XG ) dans l[RIGHT SINGLE QUOTATION MARK]ensemble des mesures de Radon sur X. Cette application
se factorise par L1loc (XG )  tensor  A n,n (XG ). On la notera bien sur alpha 7-> alphaRad .
   Cette construction s[RIGHT SINGLE QUOTATION MARK]etend en des morphismes de G-faisceaux
                  Lloc,X
                    1
                        G
                           tensor AXG AXn,n
                                  G
                                     -> Lloc,X
                                         1
                                              G
                                                 tensor AXG AXn,n
                                                        G
                                                           -> MRad,XG

sur X, ou MRad,XG est le G-faisceau des mesures de Radon sur X.
  Une section du faisceau Lloc,X
                               1
                                   G
                                      tensor AXG AXp,q
                                             G
                                                sera appelee G-forme de type (p, q) a
coefficients localement integrables.




                             [SECTION SIGN] 5.9. MESURES DE BORD

  Soit X un espace G-tropical calibre de niveau n.
5.9.1. -- Soit omega  in  FXn-1,nG
                            (X), soit sigma une partie squelettique fermee de X qui supporte
fortement omega. En accouplant la restriction a sigma(n) de la forme omega et le calibrage canonique
de sigma(n) par la construction du [SECTION SIGN]4.6.3, on obtient une forme-volume (n-1)-dimensionnelle
sur sigma(n) . D[RIGHT SINGLE QUOTATION MARK]apres la proposition 4.6.5 et la definition d[RIGHT SINGLE QUOTATION MARK]un support fort, la forme-volume
\end{lstlisting}

\subsubsection{CLD.txt: cached text lines
7070--7101}\label{cld.txt-cached-text-lines-70707101}

\begin{lstlisting}



               [SECTION SIGN] 5.10. INTEGRALES ET TROPICALISATIONS

  Soit X un espace G-tropical calibre de niveau n.
Proposition 5.10.1. -- Soit f : X -> E une carte G-tropicale agreable, soit pif le
polyedre tropical de la carte f et soit [MICRO SIGN]pif son calibrage.
  Soit p  in  {n-1, n} et soit alpha une (p, n)-forme a coefficients boreliens sur un sous-espace
paralineaire P de E contenant f (X).
  a) Si p = n, on a l[RIGHT SINGLE QUOTATION MARK]egalite
                                   ftrop,* (f * alpha)Bor >= <alpha, [MICRO SIGN]pif >Bor
                             (n)
de mesures de Borel sur pif .
  b) Si p = n - 1, on a l[RIGHT SINGLE QUOTATION MARK]inegalite
                                   ftrop,* (f * alpha)Bor >= <alpha, [MICRO SIGN]pif >Bor
                             (n)
de mesures de Borel sur pif .
  c) La mesure (f * alpha)Bor est localement finie si et seulement si la forme alpha|pif est locale-
ment integrable sur pif ; dans ce cas, on a l[RIGHT SINGLE QUOTATION MARK]egalite
                                   ftrop,* (f * alpha)Rad = <alpha, [MICRO SIGN]pif >Rad
de mesures de Radon sur pif .
   d) En particulier, si alpha est integrable sur pif , la forme f * alpha est integrable sur X, on a
l[RIGHT SINGLE QUOTATION MARK]egalite
                                   Z         Z
                                        *
(5.10.1.1)                            f alpha=       <alpha, [MICRO SIGN]pif >
                                        X           pif

lorsque p = n et l[RIGHT SINGLE QUOTATION MARK]egalite
                                      Z boundary           Z boundary 
                                            *
\end{lstlisting}

\subsubsection{CLD.txt: cached text lines
8130--8168}\label{cld.txt-cached-text-lines-81308168}

\begin{lstlisting}
7.2.6. -- L[RIGHT SINGLE QUOTATION MARK]image directe d[RIGHT SINGLE QUOTATION MARK]un courant positif (resp. faiblement positif, fortement po-
sitif) par un morphisme topologiquement propre d[RIGHT SINGLE QUOTATION MARK]espaces tropicaux est de meme type.
Cela resulte du fait que la meme propriete vaut pour l[RIGHT SINGLE QUOTATION MARK]image inverse des formes ([SECTION SIGN]3.5.2).
Lemme 7.2.7. -- Soit alpha une forme lisse symetrique sur X et soit T un courant syme-
trique sur X.
   a) Supposons que la forme alpha soit fortement positive. Alors, si T est positif ( resp.
faiblement positif, resp. fortement positif ), il en est de meme de T  wedge  alpha.
   b) Supposons que la forme alpha soit positive. Si T est positif, le courant T  wedge  alpha est positif.
   c) Supposons que la forme alpha soit faiblement positive. Si T est fortement positif, le
courant T  wedge  alpha est faiblement positif.
Demonstration. -- Ces assertions decoulent des proprietes de positivite d[RIGHT SINGLE QUOTATION MARK]un produit de
deux formes (prop. 3.5.5).

Proposition 7.2.8. -- Soit X un espace tropical de niveau n, localement differentiel-
lement separe.
  a) Toute mesure de Radon positive [MICRO SIGN] sur X induit, par restriction a Ac (X), un courant
positif [[MICRO SIGN]] d[RIGHT SINGLE QUOTATION MARK]ordre 0.
  b) L[RIGHT SINGLE QUOTATION MARK]application [MICRO SIGN] 7-> [[MICRO SIGN]] ainsi definie est une bijection continue de l[RIGHT SINGLE QUOTATION MARK]espace des me-
sures de Radon positives sur X sur l[RIGHT SINGLE QUOTATION MARK]espace des courants positifs de bidimension (0, 0)
sur X.
Demonstration. -- a) Soit [MICRO SIGN] une mesure de Radon positive sur X ; demontrons que
la forme lineaire [[MICRO SIGN]] est un courant ; ce sera alors un courant positif. Soit K une partie
compacte de X et soit D un jeu de cartes couvrant sur X. Il suffit de demontrer que
la restriction de [[MICRO SIGN]] a l[RIGHT SINGLE QUOTATION MARK]espace AD,K
                                  0,0
                                      (X) est continue pour sa topologie C 0 . Soit ||[MICRO SIGN]||K

162                  CHAPITRE 7. COURANTS SUR UN ESPACE TROPICAL

la norme de [MICRO SIGN] sur K : pour toute fonction continue a support compact sur K, on a
|[MICRO SIGN](f )| <= ||[MICRO SIGN]||K ||f ||K . Cette inegalite vaut a fortiori pour toute fonction de AD,K (X) et
entraine la continuite voulue.
   b) Cette implication est injective car l[RIGHT SINGLE QUOTATION MARK]espace des fonctions lisses a support compact
est dense dans l[RIGHT SINGLE QUOTATION MARK]espace des fonctions continues a support compact (corollaire 6.2.6).
   Il reste a prouver qu[RIGHT SINGLE QUOTATION MARK]elle est surjective. Soit S un courant positif de bidimension (0, 0)
sur X. Soit K une partie compacte de X. Soit lamda une fonction lisse et positive, a support
compact sur X, qui vaut 1 au voisinage de K. Pour toute fonction lisse f sur X dont
le support est contenu dans K, on a -||f ||K lamda <= f <= ||f ||K lamda. Comme S est positif, il en
resulte l[RIGHT SINGLE QUOTATION MARK]inegalite
\end{lstlisting}

\subsubsection{CLD.txt: cached text lines
16809--16927}\label{cld.txt-cached-text-lines-1680916927}

\begin{lstlisting}

CHAPITRE 12


PARTIES SQUELETTIQUES D[RIGHT SINGLE QUOTATION MARK]UN ESPACE
ANALYTIQUE




          [SECTION SIGN] 12.1. TORES ET LEURS SQUELETTES CANONIQUES

12.1.1. Convention. -- Dans la suite de ce travail, nous travaillerons tres reguliere-
ment avec les analytifies An,an et Gn,an
                                      m   des k-schemas An et Gnm , et seulement excep-
tionnellement avec ces schemas eux-memes. Aussi ecrirons-nous desormais, afin d[RIGHT SINGLE QUOTATION MARK]alleger
l[RIGHT SINGLE QUOTATION MARK]ecriture, An et Gnm au lieu de An,an et Gn,an
                                           m . Il sera donc considere par defaut que A
                                                                                        n

et Gnm sont des espaces analytiques ; si nous avons besoin de travailler par endroit avec
les schemas correspondants, nous le preciserons explicitement.
12.1.2. -- Dans ce travail, le groupe multiplicatif et ses puissances jouent un role
fondamental. Rappelons que l[RIGHT SINGLE QUOTATION MARK]anneau des fonctions algebriques sur Gnm est l[RIGHT SINGLE QUOTATION MARK]anneau
A = k[T1+/-1 , . . . , Tn+/-1 ] des polynomes de Laurent a coefficients dans k en des indetermi-
nees T1 , . . . , Tn ; l[RIGHT SINGLE QUOTATION MARK]espace analytique Gnm est l[RIGHT SINGLE QUOTATION MARK]ensemble des semi-normes multiplicatives
sur cet anneau qui etendent la valeur absolue de k. L[RIGHT SINGLE QUOTATION MARK]espace Gnm est un ouvert de Zariski
dense de An : c[RIGHT SINGLE QUOTATION MARK]est le lieu des points en lesquels aucune fonction coordonnee ne s[RIGHT SINGLE QUOTATION MARK]annule.
                                                                 n
Pour tout polyrayon r = (r1 , . . . , rn ) le point etar de de A     (defini a l[RIGHT SINGLE QUOTATION MARK]exemple 11.1.12)
                     n
                                                                    aI T I 7-> max|aI |rI sur l[RIGHT SINGLE QUOTATION MARK]an-
                                                                P
appartient a Gm . Il correspond a la norme multiplicative
neau des polynomes de Laurent. Si le corps k a besoin d[RIGHT SINGLE QUOTATION MARK]etre rappele nous ecrirons etak,r
au lieu de etar .
   Pour eviter le recours a un choix explicite de coordonnees, nous adopterons la conven-
tion suivante : nous appellerons tore (de dimension n) tout groupe analytique isomorphe
a une puissance Gnm du groupe multiplicatif ; nos tores sont donc implicitement supposes
deployes.
12.1.3. -- Rappelons que toute fonction analytique inversible sur Gnm est un monome,
c[RIGHT SINGLE QUOTATION MARK]est-a-dire de la forme aT1m1 . . . Tnmn , ou a  in  k * et (m1 , . . . , mn )  in  Zn est un vecteur
appele exposant. (En dimension 1, cet enonce est le lemme 9.7.1/1 de [20], le cas general
en decoule par recurrence.) C[RIGHT SINGLE QUOTATION MARK]est en particulier une fonction algebrique.
   Le morphisme Gnm -> Gm induit par une telle fonction est un homomorphisme de
groupes analytiques si et seulement si a = 1.

330          CHAPITRE 12. PARTIES SQUELETTIQUES D[RIGHT SINGLE QUOTATION MARK]UN ESPACE ANALYTIQUE

  Soit T un tore. L[RIGHT SINGLE QUOTATION MARK]ensemble X(T ) des caracteres de T , c[RIGHT SINGLE QUOTATION MARK]est-a-dire des morphismes de
groupes analytiques de T vers Gm , est un groupe abelien.
  On deduit de ce qui precede que tout isomorphisme de T avec Gnm induit un iso-
morphisme de X(T ) avec Zn ; on en deduit aussi que si T et T ' sont deux tores, tout
morphisme d[RIGHT SINGLE QUOTATION MARK]espaces k-analytiques de T vers T ' est affine, i.e. compose d[RIGHT SINGLE QUOTATION MARK]un morphisme
de groupes analytiques et d[RIGHT SINGLE QUOTATION MARK]une translation ; c[RIGHT SINGLE QUOTATION MARK]est donc un morphisme de groupes si et
seulement si il envoie l[RIGHT SINGLE QUOTATION MARK]origine de T sur l[RIGHT SINGLE QUOTATION MARK]origine de T ' . Pour cette raison, lorsque nous
evoquerons un morphisme entre tores qui n[RIGHT SINGLE QUOTATION MARK]est pas necessairement un morphisme de
groupes analytiques, nous le qualifierons de morphisme affine.
12.1.4. Tropicalisations. -- On considere le groupe abelien R x             + comme un R-espace
vectoriel, via l[RIGHT SINGLE QUOTATION MARK]exponentiation. Par ailleurs, le choix d[RIGHT SINGLE QUOTATION MARK]une base de logarithmes permet de
l[RIGHT SINGLE QUOTATION MARK]identifier a l[RIGHT SINGLE QUOTATION MARK]espace R. On fixe une fois pour toutes un tel isomorphisme log : R x        + -> R.
Dans ce texte, il n[RIGHT SINGLE QUOTATION MARK]y a pas d[RIGHT SINGLE QUOTATION MARK]interet particulier a choisir le logarithme neperien ; en
particulier, nous n[RIGHT SINGLE QUOTATION MARK]aurons jamais a deriver cet isomorphisme, vu comme fonction d[RIGHT SINGLE QUOTATION MARK]une
variable reelle. L[RIGHT SINGLE QUOTATION MARK]utilisation de cet isomorphisme log presentera pour nous deux interets :
elle nous permettra d[RIGHT SINGLE QUOTATION MARK]utiliser la geometrie paralineaire sous sa forme standard, plutot
que sa variante multiplicative ; et elle donnera lieu a un formalisme apparente a celui de
la geometrie complexe, dans lequel les logarithmes de fonctions holomorphes inversibles
jouent un role majeur.
   Soit T un tore. Nous noterons Ttrop le R-espace vectoriel Hom(X(T ), R x       + ), que l[RIGHT SINGLE QUOTATION MARK]on peut
identifier a (R x  + )n
                      une fois choisie une trivialisation de T , ou a Rn
                                                                         via l[RIGHT SINGLE QUOTATION MARK]isomorphisme    log
choisi. L[RIGHT SINGLE QUOTATION MARK]espace vectoriel Ttrop est appele tropicalisation du tore T .
   Etant donne un point t de T , l[RIGHT SINGLE QUOTATION MARK]application qui applique un caractere phi  in  X(T ) sur
l[RIGHT SINGLE QUOTATION MARK]element |phi(t)| est un element de Ttrop note trop(t). L[RIGHT SINGLE QUOTATION MARK]application trop : T -> Ttrop ainsi
definie est continue.
   La formation de Ttrop est fonctorielle : tout morphisme affine f : T -> T ' entre tores
                                                   '
induit une application affine ftrop : Ttrop -> Ttrop    telle que le diagramme

                                                   f
                                          T   <-         -> T'
                                          <-




                                                          <-->
                                          ->




                                                  ftrop   '
                                         Ttrop   <-     -> Ttrop

commute.
12.1.5. Squelette de Gnm . -- L[RIGHT SINGLE QUOTATION MARK]application de tropicalisation Gnm -> (Gnm )trop possede
une section naturelle, notee sq et definie comme suit : pour tout point r = (r1 , . . . , rn )  in 
(Gnm )trop = Rn , on note sq(r) le point etar de Gnm . L[RIGHT SINGLE QUOTATION MARK]application sq est continue ; elle induit
un homeomorphisme de (R x         n                                  n        n
                             + ) sur un sous-espace ferme S(Gm ) de Gm que l[RIGHT SINGLE QUOTATION MARK]on appelle
                           n
squelette canonique de Gm .
  On munit S(T ) de l[RIGHT SINGLE QUOTATION MARK]unique structure de groupes pour laquelle sq est un isomomor-
phisme de groupes. Pour tout extension valuee K de k, l[RIGHT SINGLE QUOTATION MARK]application naturelle T (K) ->
S(T ) donnee par t 7-> sq(trop(t)) est un morphisme de groupes.
  Par definition, pour qu[RIGHT SINGLE QUOTATION MARK]un point x P    in  Gnm appartienne a S(Gnm ), il faut et il suffit que
pour tout polynome de Laurent f = am T m , on ait

                         |f (x)| = max|am | * |T1 (x)|m1 . . . |Tn (x)|mn .
                                     m

                      [SECTION SIGN] 12.1. TORES ET LEURS SQUELETTES CANONIQUES                              331

(Si x = sq(r), la formule est vraie par definition ; dans l[RIGHT SINGLE QUOTATION MARK]autre sens, cette formule prouve
que x = etar avec ri = |Ti (x)| pour tout i.)
  Du point de vue de la reduction graduee, cette derniere condition revient a demander
que les elements T] i (x) du corps residuel gradue H (x) soient algebriquement indepen-
\end{lstlisting}

\subsubsection{CLD.txt: cached text lines
18040--18119}\label{cld.txt-cached-text-lines-1804018119}

\begin{lstlisting}
une valeur constante dans Q log|k  x  | hors d[RIGHT SINGLE QUOTATION MARK]une partie compacte de U ; elles s[RIGHT SINGLE QUOTATION MARK]identifient
a des fonctions continues sur U ' . Pour u, v  in  H, les fonctions min(u, v) et max(u, v)
appartiennent encore a H (lemme 12.5.3).
   Soit x, y des points de U ' tels que x /= y, soit a, b  in  R et soit epsilon > 0. Demontrons qu[RIGHT SINGLE QUOTATION MARK]il
existe v  in  H tel que |v(x) - a| < epsilon et |v(y) - b| < epsilon.
   Par symetrie, on peut supposer que x  in  U . Soit V un voisinage compact de x qui
est contenu dans U . D[RIGHT SINGLE QUOTATION MARK]apres le lemme 12.5.6 applique il existe une fonction paralineaire
(resp. paralineaire stricte) u sur X qui est strictement positive en x et nulle hors de V .
En particulier, son support est compact, contenu dans U , et l[RIGHT SINGLE QUOTATION MARK]on a u(y) = 0. Choisissons
s  in  Q log|k  x  | tel que |s - b| < epsilon, puis un nombre rationnel r tel que |ru(x) + s - a| < epsilon.
La fonction v = ru + s est paralienaire (strictement paralineaire stricte) sur X, elle
verifie |v(x) - a| < epsilon et |v(y) - b| < epsilon.

350         CHAPITRE 12. PARTIES SQUELETTIQUES D[RIGHT SINGLE QUOTATION MARK]UN ESPACE ANALYTIQUE

     Selon la proposition 2 de Bourbaki, Topologie generale X, p. 34, ceci entraine que H
est dense dans l[RIGHT SINGLE QUOTATION MARK]espace des fonctions numeriques continues sur U ' .
     Soit f une fonction continue sur X dont le support est compact et contenu dans U .
Soit epsilon > 0. D[RIGHT SINGLE QUOTATION MARK]apres ce qui precede, il existe une fonction g  in  H telle que |f (x) - g(x)| <
epsilon/2 pour tout x  in  U .
     Si U est compact, on prolonge g par 0 sur l[RIGHT SINGLE QUOTATION MARK]ouvert X U . On obtient une fonction
paralineaire (resp. strictement paralineaire) dont le support est compact et contenu
dans U , qui approche f a epsilon pres.
     Supposons que U n[RIGHT SINGLE QUOTATION MARK]est pas compact. Par definition de l[RIGHT SINGLE QUOTATION MARK]ensemble H, il existe a  in 
Q log|k  x  | tel que g ' = g - a soit une fonction paralineaire (resp. paralineaire stricte)
a support compact dans U . Comme f est a support compact dans U , on a |a| <= epsilon/2.
Soit x  in  X. Si x  in  U , on a |f (x) - g ' (x)| <= |f (x) - g(x)| + |a| <= epsilon. Sinon, on a
g ' (x) = 0 = f (x). Cela conclut la demonstration de la proposition.

Corollaire 12.5.8. -- On suppose que la valeur absolue de k n[RIGHT SINGLE QUOTATION MARK]est pas triviale. Soit X
un espace k-analytique ( resp. strictement k-analytique) topologiquement separe et loca-
lement holomorphiquement separe.
  a) L[RIGHT SINGLE QUOTATION MARK]ensemble des fonctions paralineaires ( resp. paralineaires strictes) a support com-
pact sur X est dense dans l[RIGHT SINGLE QUOTATION MARK]espace Cc (X) des fonctions continues a support compact
sur X, muni de la topologie limite inductive des espaces de fonctions a support dans un
compact donne.
  b) L[RIGHT SINGLE QUOTATION MARK]ensemble des fonctions paralineaires ( resp. strictement paralineaires) a support
compact sur X est dense dans l[RIGHT SINGLE QUOTATION MARK]espace C (X) des fonctions continues sur X muni de la
topologie de la convergence uniforme sur tout compact.
   La demonstration se deduit du theoreme 12.5.7 de la meme facon que le corollaire A.2.6
se deduisait du theoreme A.2.5, mutatis mutandis.



        [SECTION SIGN] 12.6. POLYEDRE CARACTERISTIQUE D[RIGHT SINGLE QUOTATION MARK]UN MOMENT

   Si P un espace paralineaire et n un entier, rappelons que l[RIGHT SINGLE QUOTATION MARK]on note P (n) l[RIGHT SINGLE QUOTATION MARK]ensemble des
points x  in  P tels que dimx (P ) >= n ; c[RIGHT SINGLE QUOTATION MARK]est un sous-espace paralineaire ferme de P .
   Les paragraphes qui suivent visent a construire, pour tout espace analytique topologi-
quement separe X, une notion de partie squelettique. D[RIGHT SINGLE QUOTATION MARK]apres la construction generale,
il va suffire de definir la notion de partie squelettique elementaire et le procede que nous
proposons suppose donne un entier naturel n tel que dim(X) <= n. Meme si la theorie
ne voit pas les composantes irreductibles de X de dimension < n, il est commode, pour
des raisons de redaction, de ne pas supposer que X est equidimensionnel, ni que n est
egal a dim(X).
   Si n est un entier naturel, l[RIGHT SINGLE QUOTATION MARK]expression [LEFT-POINTING DOUBLE ANGLE QUOTATION MARK] Soit X un espace analytique de niveau n [RIGHT-POINTING DOUBLE ANGLE QUOTATION MARK]
signifiera ainsi que X est un espace analytique topologiquement separe tel que dim(X) <=
n, et que l[RIGHT SINGLE QUOTATION MARK]on garde trace de cet entier n.
   Lorsque L est une extension complete de k, on considerera implicitement l[RIGHT SINGLE QUOTATION MARK]espace L-
analytique XL obtenu par changement de base comme un espace analytique de niveau n.
(Il est topologiquement separe car le morphisme de changement de base XL -> X est
topologiquement separe.)

                     [SECTION SIGN] 12.6. POLYEDRE CARACTERISTIQUE D[RIGHT SINGLE QUOTATION MARK]UN MOMENT                            351

Definition 12.6.1. -- Soit X un espace k-analytique de niveau n. Soit T un tore et
soit f : X -> T un moment sur X. On appelle polyedre caracteristique du moment f
la reunion des images reciproques par q [WHITE BULLET] f du squelette canonique de Gnm , lorsque q
parcourt l[RIGHT SINGLE QUOTATION MARK]ensemble des morphismes de tores de T dans Gnm . On le note sigmaf .
   Soit x  in  sigmaf . Soit q : T -> Gnm un morphisme de tores tel que q [WHITE BULLET] f (x)  in  S(Gnm ).
On a donc dk (x) >= dk (f (x)) >= dk (q [WHITE BULLET] f (x)) = n = dim(X), donc dk (x) = dim(X) =
dk (f (x)) = n. Par suite, le point x n[RIGHT SINGLE QUOTATION MARK]appartient a aucun ferme de Zariski de dimension <
dim(X) d[RIGHT SINGLE QUOTATION MARK]aucun domaine analytique de X. En particulier, le point x n[RIGHT SINGLE QUOTATION MARK]appartient qu[RIGHT SINGLE QUOTATION MARK]a
une seule composante irreductible de X dans laquelle il est Zariski-dense.
   Cela demontre aussi que sigmaf est vide si dim(T ) < n.
   Notons fred : Xred -> T le moment deduit de f par restriction. On a sigmafred = sigmaf .
Remarque 12.6.2. -- Soit X un espace k-analytique de niveau n et soit f : X -> T
un moment sur X. Le polyedre caracteristique sigmaf est contenu dans la reunion X (n) des
composantes irredutibles de X qui sont de dimension n ; en particulier, cette theorie
ignore les eventuelles composantes de dimension < n. On pourrait travailler sur chacune
\end{lstlisting}

\subsubsection{CLD.txt: cached text lines
19153--19240}\label{cld.txt-cached-text-lines-1915319240}

\begin{lstlisting}

Theoreme 12.11.1. -- Soit X un espace k-analytique sans bord, topologiquement se-
pare et purement de dimension n ; considerons-le comme un espace analytique de ni-
veau n. Alors X est tres riche (definition 5.2.7).
Demonstration. -- D[RIGHT SINGLE QUOTATION MARK]apres le corollaire 12.10.7, l[RIGHT SINGLE QUOTATION MARK]espace X est riche ; il reste a demontrer
que toute partie squelettique compacte sigma de X est contenue dans une partie squelettique
purement de dimension n.
   Soit x un point de sigma. Le germe (X, x) est sans bord, donc presque bon ; d[RIGHT SINGLE QUOTATION MARK]apres la
proposition 12.9.10, il existe un voisinage ouvert Ux de x et un moment fx : Ux -> Gnm
tel que sigma  intersect  Ux  subset  sigmafx . Considerons un voisinage analytique compact Vx de x dans Ux .
Demontrons que x  in  (sigmafx  intersect  Vx )(n) .
   Comme sigma est compact, il existe une partie finie S de sigma telle que les V x , pour x  in  S,
recouvrent sigma. Soit T la reunion des (sigmafs  intersect Vs )(n) , pour s  in  S. c[RIGHT SINGLE QUOTATION MARK]est une partie squelettique
compacte de X, purement de dimension n.
   Demontrons que sigma  subset  T. Soit x  in  sigma ; soit s  in  S tel que x  in  V s . Comme X est sans bord,
l[RIGHT SINGLE QUOTATION MARK]ouvert Vs est sans bord, donc sigmafs  intersect  V s est purement de dimension n (corollaire 12.8.4),
si bien que x appartient a (sigmafs  intersect  Vs )(n) . A fortiori, x  in  T.

Remarque 12.11.2. -- On ne peut pas, en general, enlever l[RIGHT SINGLE QUOTATION MARK]hypothese que X est
sans bord. Supposons par exemple qu[RIGHT SINGLE QUOTATION MARK]il existe r  in       / |k  x  |Q et considerons le domaine
analytique de A1k defini par |T | = r. C[RIGHT SINGLE QUOTATION MARK]est un espace k-analytique de dimension 1, reduit
au singleton {etar }, qui est aussi son bord. Il ne contient aucun squelette de dimension 1,
et n[RIGHT SINGLE QUOTATION MARK]est donc pas tres riche.
Remarque 12.11.3. -- La structure des courbes k-analytiques (pour laquelle nous
renvoyons au chapitre 4 de [8], et a la presentation plus generale de [57]) entraine que
toute courbe k-analytique ne comportant aucun point de type (3) isole est tres riche.
En effet, soit X une telle courbe et soit x  in  X un point tel que dk (x) = 1. Cela entraine
que X est de type (2) ou (3). Si x est de type (2), il existe une branche aboutissant a x ;
il en est de meme si x est de type (3) puisque, par hypothese, il n[RIGHT SINGLE QUOTATION MARK]est pas isole.
   En particulier, pour tout r > 0, le disque de rayon r de A1k est tres riche.
Remarque 12.11.4. -- Nous nous attendons a ce que le polydisque unite de Ank soit
tres riche.
   En revanche, si r  in   / |k  x  |Q, le bidisque X de birayon (r, 1/r) n[RIGHT SINGLE QUOTATION MARK]est pas tres riche.
Considerons en effet l[RIGHT SINGLE QUOTATION MARK]ouvert U de X defini par |T1 T2 - 1| < 1 ; nous allons demontrer
qu[RIGHT SINGLE QUOTATION MARK]il est non vide et contenu dans la fibre de la premiere projection en etar , qui est une
courbe H (etar )-analytique, et ne contient donc aucune partie squelettique de dimension 2.
   Soit x  in  U . Comme |T1 (x)T2 (x) - 1| < 1, on a |T1 (x)||T2 (x)| = 1 ; par definition de X,
on a |T1 (x)| <= r et |T2 (x)| <= 1/r, si bien que |T1 (x) = r|. Alors y = p1 (x) est un point
de A1k tel que |T1 (y)| = r ; comme r  in  / |k  x  |Q , on a y = etar .
   Il reste a verifier que U n[RIGHT SINGLE QUOTATION MARK]est pas vide. Mais p-1                                 1
                                                         1 (etar ) est le domaine de AH (etar ) , de
coordonnee T2 , defini par |T2 | <= 1/r. Le point H (etar )-rationnel de ce domaine defini par
T2 = T1 (etar )-1 appartient a U .


CHAPITRE 13


COMPLEMENTS DE GEOMETRIE ANALYTIQUE




                        [SECTION SIGN] 13.1. ANNEAUX G-LOCAUX

13.1.1. -- Soit X un espace k-analytique et soit x un point de X. Nous avons defini
plus haut la dimension centrale du germe (X, x) (11.5.5). Celle-ci possede une variante
G-topologique : la dimension centrale du G-germe (XG , x), notee dimc (XG , x), et de-
                                               VZar
finie comme le minimum des dimensions de {x}        lorsque V parcourt l[RIGHT SINGLE QUOTATION MARK]ensemble des
domaines analytiques de X contenant x.
   On a de maniere evidente

                                                       XZar
                    dk (x) <= dimc (XG , x) <= dim {x}          <= dimx X ;

en particulier dimc (XG , x) = dimx X des que dk (x) = n.
  Soit V un domaine analytique de X contenant x. L[RIGHT SINGLE QUOTATION MARK]ensemble des domaines analytiques
de V contenant x est cofinal dans l[RIGHT SINGLE QUOTATION MARK]ensemble des domaines analytiques de X contenant x ;
par consequent dimc (VG , x) = dimc (XG , x).
  Soit Z un sous-espace analytique ferme de X contenant x. L[RIGHT SINGLE QUOTATION MARK]ensemble des domaines
analytiques de Z contenant x de la forme V  intersect  Z, ou V est un domaine analytique de
X contenant x, est cofinal dans l[RIGHT SINGLE QUOTATION MARK]ensemble des domaines analytiques de X contenant x.
                                                                          (V  intersect Z)Zar
Par consequent Par consequent dimc (ZG , x) est le minimum des dim {x}              ou V
parcourt l[RIGHT SINGLE QUOTATION MARK]ensemble des domaines analytiques de X contenant x. Mais pour un tel V on
      (V  intersect Z)Zar      VZar
a {x}           = {x} , si bien que dimc (ZG , x) = dimc (XG , x).
13.1.2. -- Nous noterons X gen l[RIGHT SINGLE QUOTATION MARK]ensemble des points x de X tels que dimc (XG , x) =
dimx X. Cet ensemble contient {x  in  X, dk (x) = dimx X}.
   Si V est un domaine analytique de X on a dimc (VG , x) = dimc (XG , x) et dimx V =
dimx X pour tout x  in  V , si bien que V gen = X gen  intersect  V .
   Si Z est un sous-espace analytique ferme de X on a dimc (ZG , x) = dimc (XG , x) pour
tout x  in  Z, si bien que Z gen est l[RIGHT SINGLE QUOTATION MARK]ensemble des points x de Z tels que dimc (ZG , x) =
dimx Z ; cet ensemble ne depend que du ferme de Zariski sous-jacent a Z, et pas de la
structure de sous-espace analytique ferme choisie sur Z.

372              CHAPITRE 13. COMPLEMENTS DE GEOMETRIE ANALYTIQUE
\end{lstlisting}

\subsubsection{CLD.txt: cached text lines
18355--18407}\label{cld.txt-cached-text-lines-1835518407}

\begin{lstlisting}
laire 3.14).

Corollaire 12.7.6. -- On sigmafL  intersect  p-1 (Xu ) = sigmaL (sigmaf  intersect  Xu ). En particulier, si k est
algebriquement clos, on a sigmafL = sigmaL (sigmaf ).

  Lorsqu[RIGHT SINGLE QUOTATION MARK]on ne fait pas d[RIGHT SINGLE QUOTATION MARK]hypothese sur l[RIGHT SINGLE QUOTATION MARK]extension L de k, nous pouvons preciser un
resultat de [50].

Proposition 12.7.7. -- a) Pour toute extension complete L de k, l[RIGHT SINGLE QUOTATION MARK]application p
induit une immersion par morceaux surjective de sigmafL sur sigmaf qui preserve la dimension
locale.
   b) Il existe une extension algebrique separable F de k, de completee Fb, telle que,
pour toute extension complete L de Fb, la projection canonique de sigmafL sur sigmafFb est un
isomorphisme d[RIGHT SINGLE QUOTATION MARK]espaces paralineaires. Si, de plus, X est compact, on peut prendre F
finie sur k (et donc Fb = F ).
   c) En particulier, si k est algebriquement clos, alors la projection canonique de sigmafL
sur sigmaf est un isomorphisme, pour toute extension complete L de k.

356         CHAPITRE 12. PARTIES SQUELETTIQUES D[RIGHT SINGLE QUOTATION MARK]UN ESPACE ANALYTIQUE

Demonstration. -- D[RIGHT SINGLE QUOTATION MARK]apres [50, theoreme 5.1, 3)], l[RIGHT SINGLE QUOTATION MARK]application p induit une immersion
par morceaux surjective de sigmafL sur sigmaF , ce qui demontre la premiere partie de l[RIGHT SINGLE QUOTATION MARK]asser-
tion a). Si X est compact, le theoreme 5.1, 4) de [50] fournit l[RIGHT SINGLE QUOTATION MARK]assertion b), l[RIGHT SINGLE QUOTATION MARK]extension F
de k etant en outre finie. Puisque la formation des polyedres caracteristiques est G-locale,
l[RIGHT SINGLE QUOTATION MARK]assertion b) en resulte en composant toutes les extensions finies de k fournies par un
recouvrement affinoide de X. L[RIGHT SINGLE QUOTATION MARK]assertion c) en decoule immediatement.
   Prouvons maintenant que le morphisme canonique de sigmafL sur sigmaf preserve la dimension
locale. On peut raisonner G-localement sur X, et supposer que X est compact. Soit F
une extension finie separable de k, qu[RIGHT SINGLE QUOTATION MARK]on suppose galoisienne, stabilisant les polyedres
caracteristiques et soit M l[RIGHT SINGLE QUOTATION MARK]extension composee de L et F . Soit y  in  sigmafL , soit x son
image dans sigmaf , soit z un point de sigmafM au-dessus de y et t sa projection dans sigmaf .
Comme la projection sigmafM -> sigmafF est un isomorphisme, on a dimz (sigmafM ) = dimt (sigmaf ).
D[RIGHT SINGLE QUOTATION MARK]apres la proposition 12.7.2, le polyedre caracteristique sigmaf est le quotient de sigmafF par
l[RIGHT SINGLE QUOTATION MARK]action de Gal(F/k) ; on a donc dimt (sigmafF ) = dimx (sigmaf ), donc dimz (sigmafM ) = dimx (sigmaf ).
Comme les projections sigmafM -> sigmafL et sigmafL -> sigmaf sont des immersions par morceaux
d[RIGHT SINGLE QUOTATION MARK]espaces paralineaires, on a dimz (sigmafM ) <= dimy (sigmafL ) <= dimx (sigmaf ). Cela demontre que
dimy (sigmafL ) = dimx (sigmaf ).



      [SECTION SIGN] 12.8. POLYEDRES CARACTERISTIQUES ET DIMENSION
                            TROPICALE

12.8.1. -- Soit X un espace k-analytique de niveau n, soit m un entier et soit f =
X -> Gm   m un moment sur X. On note (f1 , . . . , fm ) les composantes de f . Pour toute
partie I de {1, . . . , m}, on note qI la projection canonique de Gm      I
                                                                   m sur Gm .
  Rappelons (proposition 12.6.8) que le polyedre caracteristique de f est une partie
paralineaire de X et que l[RIGHT SINGLE QUOTATION MARK]application ftrop |sigmaf est une immersion par morceaux d[RIGHT SINGLE QUOTATION MARK]espaces
paralineaires.
  Le but de ce numero est de relier la dimension tropicale de ftrop aux proprietes de sigmaf .
Theoreme 12.8.2. -- Soit x un point de X.
                                                                    (n)
\end{lstlisting}

\subsubsection{CLD.txt: cached text lines
20320--20360}\label{cld.txt-cached-text-lines-2032020360}

\begin{lstlisting}
pour tout corps residuel F de L tensor 
                               b k K.

Demonstration. -- Choisissons une presentation analytique (X, x) de L. D[RIGHT SINGLE QUOTATION MARK]apres le
lemme 13.8.6, on peut supposer qu[RIGHT SINGLE QUOTATION MARK]il existe un morphisme f : X -> Ank tel que
|fi (x) - ti | < |ti | pour tout i. Alors, les extensions K/k(t)
                                                            d et K/k(f \  (x)) sont finies, ont
meme degre et meme degre residuel ; elles ont donc meme defaut.
   On peut ainsi supposer que ti = fi (x) pour tout i. Le point f (x) de Ank est alors
egal a etar , ou r = (|t1 |, . . . , |tn |). Le morphisme f est dimension nulle en x ; quitte a
restreindre X, on peut supposer que f -1 (etar ) = {x}.
   Notons pi : XK -> X le morphisme de changement de base.
   D[RIGHT SINGLE QUOTATION MARK]apres la proposition 13.8.8, f ), les differents corps residuels F de l[RIGHT SINGLE QUOTATION MARK]anneau arti-
nien L tensor   b k K sont permutes par Aut(K/k), de sorte que les differents degres residuels
[F : (L * K)]
 e      e     e sont egaux a un meme entier d.        e D[RIGHT SINGLE QUOTATION MARK]apres la proposition 13.8.8, c), ces

                       [SECTION SIGN] 13.8. EXTENSIONS D[RIGHT SINGLE QUOTATION MARK]ABHYANKAR ET LEUR DEFAUT                         393

corps F sont de la forme H (y), pour un unique point y  in  pi -1 (x), et l[RIGHT SINGLE QUOTATION MARK]assertion d )
affirme que les degres residuels consideres sont finis.
   Notons delta le defaut de H (x) sur H (etar ).
   Il reste a demontrer que l[RIGHT SINGLE QUOTATION MARK]on a delta = elld/
                                        e ell.
                                           e On a :
                    [H (x) : H (etar )]   [(H (x)  tensor H (etar ) H (etaK,r )) : H (etaK,r )]
               delta=                     =                                          .
                    [H
                     ^ (x) : H
                             ^ (etar )]   [(H
                                          ^ (x)  tensor         H^(etaK,r )) : H^
                                                              H
                                                              ^ (etar )
                                                                        (etaK,r )]

Fixons un point y  in  f -1 (etaK,r ) et posons d = [H (y) : H (etaK,r )].
  Puisque les points de f -1 (etaK,r ) sont permutes par Aut(K/k), on a
                                                      X
          [(H (x)  tensor H (etar ) H (etaK,r )) : H (etaK,r )] =        [H (z) : H (etaK,r )]
                                                                 z in f -1 (etaK,r )
\end{lstlisting}

\subsubsection{CLD.txt: cached text lines
20700--20870}\label{cld.txt-cached-text-lines-2070020870}

\begin{lstlisting}

                                  [SECTION SIGN] 14.2. MOMENTS ET DEGRES                                   399

                                        (n)
d[RIGHT SINGLE QUOTATION MARK]un moment agreable. Par suite, sigmaf est reunion d[RIGHT SINGLE QUOTATION MARK]une famille denombrable de sous-
espaces compacts. Cela entraine qu[RIGHT SINGLE QUOTATION MARK]il est paracompact.
                                                                         (n)
  D[RIGHT SINGLE QUOTATION MARK]apres le lemme 14.1.5, il existe une decomposition cellulaire C1 de sigmaf en f -cellules
                                                                                               (n)
commodes. D[RIGHT SINGLE QUOTATION MARK]apres la proposition 2.3.6, il existe une decomposition cellulaire C de sigmaf
                                                                                              (n)
qui raffine C1 et une decomposition cellulaire D adaptee a l[RIGHT SINGLE QUOTATION MARK]inclusion de ftrop (sigmaf )
dans Ttrop tel que le couple (C , D) soit adapte a l[RIGHT SINGLE QUOTATION MARK]immersion par morceaux compacte ftrop .
Ses cellules sont des f -cellules commodes car la decomposition cellulaire C raffine C1 .
  La derniere assertion decoule de la definition d[RIGHT SINGLE QUOTATION MARK]une cellule commode.



                           [SECTION SIGN] 14.2. MOMENTS ET DEGRES

Lemme 14.2.1. -- Soit X un espace analytique de niveau n, soit f : X -> Gnm un
moment. Alors, f est quasi-fini et plat en tout point de sigmaf .
Demonstration. -- Soit x  in  sigmaf . On a f (x)  in  S(Gnm ), donc dk (f (x)) = n. D[RIGHT SINGLE QUOTATION MARK]apres
le lemme 11.5.1, le morphisme f est quasi-fini en x. Comme Gnm est reduit, le theo-
reme 10.3.7 de [52] entraine que f est plat en x.

Proposition 14.2.2. -- Soit X un espace analytique de niveau n, soit f : X -> T un
moment sur X et soit C une f -cellule commode de dimension n de X. Soit q : T -> Gnm
un morphisme de tores tel que qtrop |ftrop (C) est injectif. Alors le morphisme q [WHITE BULLET] f est de
dimension relative nulle en tout point de C et il existe un unique entier d > 0 tel que ce
morphisme soit fini et plat de degre d en tout point de C.
Demonstration. -- Soit q : T -> Gnm un morphisme de tores tel que qtrop soit injectif sur
ftrop (C). La dimension tropicale de q [WHITE BULLET] f est egale a n en tout point de C, donc C  subset 
(q [WHITE BULLET]f )-1 (S(Gnm )) d[RIGHT SINGLE QUOTATION MARK]apres la prop. 12.2.8, (iv). D[RIGHT SINGLE QUOTATION MARK]apres le lemme 11.5.1, le morphisme q [WHITE BULLET]f
est de dimension relative nulle en tout point de C.
   Comme C ne rencontre pas boundary (X), il resulte du lemme 11.6.5 que le morphisme q [WHITE BULLET]f est
fini et plat en tout point de C ; on a note degx (q [WHITE BULLET] f ) son degre et il s[RIGHT SINGLE QUOTATION MARK]agit de demontrer
que l[RIGHT SINGLE QUOTATION MARK]entier degx (q [WHITE BULLET] f ) est independant du point x  in  C.
   Soit D la cellule sq(qtrop (ftrop (C))) de Gnm . Si un point x  in  Int(X) verifie q [WHITE BULLET] f (x)  in  D,
                                                                                 (n)
alors q[WHITE BULLET]f est de dimension tropicale n en x, donc f aussi, si bien que x  in  sigmaf . Autrement
dit, on a les inclusions
                                                                (n)
                             C  subset  (q [WHITE BULLET] f )-1 (D)  intersect  Int(X)  subset  sigmaf .
                                                                        (n)
   Soit maintenant x  in  C. Par hypothese, C est ouvert dans sigmaf ; il est donc ouvert
dans (q [WHITE BULLET] f )-1 (D)  intersect  Int(X). Comme q [WHITE BULLET] f est fini et plat en x, il existe un voisinage
ouvert U de x dans Int(X) et un voisinage ouvert V de q [WHITE BULLET] f (x) dans Gnm tels que
q [WHITE BULLET] f induise un morphisme fini et plat de degre degx (q [WHITE BULLET] f ) de U sur V , et tel que
(q [WHITE BULLET] f )-1 (D)  intersect  U  subset  C. Un point y  in  C  intersect  U est alors le seul antecedent dans U de
q [WHITE BULLET] f (y) par q [WHITE BULLET] f ; on a donc degy (q [WHITE BULLET] f ) = degx (q [WHITE BULLET] f ) pour tout y  in  C  intersect  U . La
fonction x 7-> degx (q [WHITE BULLET] f ) est localement constante sur l[RIGHT SINGLE QUOTATION MARK]ensemble C ; comme celui-ci est
connexe, elle est constante.

400           CHAPITRE 14. CALIBRAGE CANONIQUE D[RIGHT SINGLE QUOTATION MARK]UN ESPACE ANALYTIQUE

Proposition 14.2.3. -- Soit X un espace k-analytique paracompact de niveau n et
soit f : X -> Gnm un moment. Soit sigma un sous-espace paralineaire paracompact de sigmaf .
Il existe une decomposition cellulaire C de sigma en f -cellules telle que pour toute cellule
C  in  C , l[RIGHT SINGLE QUOTATION MARK]application x 7-> degx (f ) soit finie et constante sur C.

    Rappelons que le degre degx (f ) d[RIGHT SINGLE QUOTATION MARK]un morphisme f en un point x a ete defini au [SECTION SIGN]13.3.1
sans hypothese de finitude ou de platitude. Il peut eventuellement etre infini.
Demonstration. -- Pour tout x  in  sigmaf , on a dk (x) = n, donc f est quasi-fini en x
(lemme 11.5.1) ; cela prouve que degx (f ) est fini.
    Nous procedons alors par recurrence sur le niveau n, et supposons donc la proposi-
tion vraie en niveau < n. Puisque sigma est paracompact, il possede une decomposition
cellulaire D constituee de f -cellules commodes (lemme 14.1.5).
    Soit D une cellule de D.
    Si D est de dimension n, la proposition 14.2.2 assure que f est fini et plat de degre
constant sur D ; en particulier, la fonction x 7-> degx (f ) est constante sur D.
    Supposons maintenant que dim(D) < n. Il existe un domaine analytique compact V
de X dont l[RIGHT SINGLE QUOTATION MARK]intersection avec sigma est egale a V (lemme 3.2.8 (e)) ; on a donc D = sigmaf |V .
Comme D est une f -cellule, f (D) est une partie du squelette canonique S(Gnm ), et
f (D)trop est une d-cellule forte [INCREMENT] de Rn . Il existe un automorphisme phi de Gnm tel que
l[RIGHT SINGLE QUOTATION MARK]espace affine engendre par phitrop ([INCREMENT]) soit defini par les n-d egalites xd+1 = * * * = xn = 0.
Quitte a changer f en phi [WHITE BULLET] f , on se ramene au cas ou phi est l[RIGHT SINGLE QUOTATION MARK]identite. La cellule D est
ainsi contenue dans f -1 (S(Gnm )  intersect  p-1 (eta1 )), ou p designe la projection de Gnm sur les n - d
                                                                                   n-d
dernieres composantes et ou eta1 designe, par abus, le point eta1,...,1 de Gm              .
    Notons V la fibre (p [WHITE BULLET] f ) (eta1 )  intersect  V . Puisque dk (eta1 ) = n - d, la dimension de V '
                   '                      -1

est au plus n - (n - d) = d ; on le considere comme un espace H (eta1 )-analytique de
niveau d. Notons f ' : V ' -> Gdm,H (eta1 ) le morphisme d[RIGHT SINGLE QUOTATION MARK]espaces H (eta1 )-analytiques donne
par (f1 , . . . , fd ). On a D = sigmaf ' . En effet, pour x  in  V ' , on a x  in  sigmaf ' si et seule-
ment si f]     1 (x), . . . , fd (x) sont algebriquement independants sur H (eta1 ) ; comme ce corps
                              ]                                             ^
est transcendant pur engendre par les f]             i (x) pour i > d + 1, cela revient a exiger que
f]                ]
  1 (x), . . . , fn (x) soit algebriquement independant sur k, c[RIGHT SINGLE QUOTATION MARK]est-a-dire x  in  sigmaf , soit encore
x  in  D puisque D = sigmaf  intersect  V .
    Puisque d < n, l[RIGHT SINGLE QUOTATION MARK]hypothese de recurrence assure l[RIGHT SINGLE QUOTATION MARK]existence d[RIGHT SINGLE QUOTATION MARK]une decomposition cel-
lulaire DD de D telle que x 7-> degx (f ' ) soit constant sur toute cellule ouverte de DD .
Comme les fibres de f ' sont celles de f , on a degx (f ' ) = degx (f ) pour tout x  in  D. Par
consequent, x 7-> degx (f ) est constant sur toute cellule ouverte de DD .
    Choisissons maintenant une decomposition cellulaire D ' de la reunion des cellules D
de D de dimension < n et qui raffine chacune des DD .
    La famille C formee des cellules de D ' et des cellules de D qui sont de dimension n
est une decomposition cellulaire de sigma. Elle repond aux exigences de l[RIGHT SINGLE QUOTATION MARK]enonce.

14.2.4. -- Conservons les notations de la proposition precedente et notons dC la valeur
constante de degx (f ) sur C. Comme x 7-> degx (f ) est semi-continue superieurement, on
a dD <= dC si C  subset  D. Lorsque X est sans bord, on a des relations plus precises donnees
par le theoreme suivant qui montre, en outre, que toute decomposition cellulaire de sigmaf
convient.

                               [SECTION SIGN] 14.2. MOMENTS ET DEGRES                                401

Theoreme 14.2.5. -- Soit X un espace k-analytique de niveau n, sans bord, soit
f : X -> Gnm un moment sur X et soit C une decomposition cellulaire de sigmaf en f -
cellules.
   a) Pour toute cellule C de C , il existe un unique entier dC > 0 tel que f soit fini et
plat de degre dC en tout point de C.
   b) Soit C  in  C et soit d un entier tel que dim(C)T<= d <= n. Soit [INCREMENT] un ensemble
non vide de cellules de C qui contiennent C, tel que D in [INCREMENT] ftrop (D) soit une cellule de
dimension d, et maximal pour cette propriete ; notons [INCREMENT]* l[RIGHT SINGLE QUOTATION MARK]ensemble de ses elements
minimaux. On a                               X
                                      dC =       dD .
                                            D in [INCREMENT]*

Demonstration. -- Sous les hypotheses du theoreme, toute cellule de C est une f -cellule
commode (remarque 14.1.4, b)). D[RIGHT SINGLE QUOTATION MARK]apres le lemme 11.6.5, le morphisme f est fini et plat
en tout point de sigmaf ; le contenu du theoreme est la constance du degre sur chaque cellule
ouverte et une formule de recurrence pour ces degres.
   On demontre ces resultats par recurrence descendante sur c = dim(C). Lorsque c = n,
la premiere assertion decoule de la proposition 14.2.2, et la seconde revient a l[RIGHT SINGLE QUOTATION MARK]egalite
triviale dC = dC , car on a alors necessairement d = n et [INCREMENT] = [INCREMENT]* = {C}.
   Soit c un entier < n ; supposons les assertions (1) et (2) verifiees pour toute cellule C
de dimension > c. Pour une telle cellule C, on a donc un degre dC bien defini.
   Soit C  in  C une cellule de dimension c. Soit d un entier tel que c <= d T    <= n. Soit EC,d
l[RIGHT SINGLE QUOTATION MARK]ensemble des parties [INCREMENT] de C formees de cellules contenant C et telles que D in [INCREMENT] ftrop (D)
soit de dimension d, et maximale pour cette propriete. Pour [INCREMENT]  in  EC,d , notons [INCREMENT]* l[RIGHT SINGLE QUOTATION MARK]en-
semble des elements minimaux de [INCREMENT], lorsque ce dernier est ordonne par l[RIGHT SINGLE QUOTATION MARK]inclusion.
   Comme sigmaf est purement de dimension n, il existe une n-cellule de C qui contient C.
Il existe alors une cellule D  in  C de dimension d qui contient C. Comme ftrop (D) est de
dimension d, il existe une partie [INCREMENT] appartenant a EC,d qui contient D. En particulier,
EC,d n[RIGHT SINGLE QUOTATION MARK]est pas vide.
Lemme 14.2.6. -- (Sous l[RIGHT SINGLE QUOTATION MARK]hypothese de recurrence.) Supposons d > c. Pour tout [INCREMENT]  in 
EC,d et tout x  in  C, on a degx (f ) = D in [INCREMENT]* dD .
                                     P

   Admettons provisoirement ce lemme.
   Rappelons que c < n ; applique a d = n et a un element de l[RIGHT SINGLE QUOTATION MARK]ensemble non vide EC,n ,
le lemme 14.2.6 entraine que la fonction x 7-> degx (f ) est constante sur l[RIGHT SINGLE QUOTATION MARK]ensemble non
vide C ; notons dC cette valeur. Cela etablit deja l[RIGHT SINGLE QUOTATION MARK]assertion (1) pour C.
   Soit d un entier tel que d <= c <= n. Le lemme 14.2.6 etablit alors egalement l[RIGHT SINGLE QUOTATION MARK]asser-
tion (2) des que d > c. Dans le cas restant ou d = c, on a [INCREMENT] = [INCREMENT]* = {C} et l[RIGHT SINGLE QUOTATION MARK]assertion (2)
se reduit alors a l[RIGHT SINGLE QUOTATION MARK]egalite tautologique dC = dC .
   Cela termine la demonstration du theoreme et il reste a demontrer le lemme 14.2.6
que nous avions admis.
   Notons WC la reunion des cellules ouvertes D, ou D parcourt l[RIGHT SINGLE QUOTATION MARK]ensemble des cellules
de C contenant C ; c[RIGHT SINGLE QUOTATION MARK]est un voisinage ouvert de C dans sigmaf . Soit V un voisinage ouvert
de x dans X tel que V  intersect  sigmaf  subset  WC et tel que f induise un morphisme fini et plat de
degre degx (f ) de V sur un voisinage ouvert U de f (x) dans Gnm .
   Notons P l[RIGHT SINGLE QUOTATION MARK]intersection des ftrop (D), pour D  in  [INCREMENT] ; c[RIGHT SINGLE QUOTATION MARK]est une cellule de dimension d
contenant ftrop (C). Soit E  in  C contenant C telle que E / in  [INCREMENT] ; alors [INCREMENT] union {E} n[RIGHT SINGLE QUOTATION MARK]appartient

402           CHAPITRE 14. CALIBRAGE CANONIQUE D[RIGHT SINGLE QUOTATION MARK]UN ESPACE ANALYTIQUE

pas a EC,d , par maximalite de [INCREMENT], de sorte que P  intersect  ftrop (E) est une cellule de dimension <
d. Notons Q1 la reunion de ces P  intersect  ftrop (E), lorsque E parcourt l[RIGHT SINGLE QUOTATION MARK]ensemble des cellules
de C contenant C et n[RIGHT SINGLE QUOTATION MARK]appartenant pas a [INCREMENT] ; c[RIGHT SINGLE QUOTATION MARK]est un sous-espace paralineaire ferme
de P , de dimension < d.
   D[RIGHT SINGLE QUOTATION MARK]autre part, soit D  in  [INCREMENT]* et soit E une cellule contenue dans boundary (D) disjointe de C ;
alors P  intersect  ftrop (E) est une cellule, disjointe de ftrop (C), car ftrop est injectif sur D. La
reunion Q2 de ces cellules P  intersect  ftrop (E) est un sous-espace paralineaire de P , ferme,
disjoint de ftrop (C), de sorte que P Q2 est un voisinage de ftrop (C) dans P .
   Posons Q = Q1  union Q2 ; c[RIGHT SINGLE QUOTATION MARK]est un sous-espace paralineaire de P , ferme. L[RIGHT SINGLE QUOTATION MARK]ensemble P Q2
est un voisinage de ftrop (C) dans P ; choisissons un ouvert [OHM SIGN] de P tel que ftrop (C)  subset 
[OHM SIGN]  subset  P Q2 . Comme Q1 est d[RIGHT SINGLE QUOTATION MARK]interieur vide dans P , sa trace sur [OHM SIGN] est d[RIGHT SINGLE QUOTATION MARK]interieur vide
\end{lstlisting}

\subsubsection{CLD.txt: cached text lines
20924--20976}\label{cld.txt-cached-text-lines-2092420976}

\begin{lstlisting}
                                                                                         '
groupe, le point yK est le seul antecedent de y dans (Gnm )K . De meme, le point yK        est
                         '         n
le seul antecedent de y dans (Gm )K .
   Notons p : XK -> X le morphisme de changement de base canonique. D[RIGHT SINGLE QUOTATION MARK]apres ce qui
precede, p-1 (x)  subset  fK-1 (yK ) et p-1 (x)  subset  (f ' )-1    '
                                                   K (yK ). Appliquons le lemme 13.3.3. On
obtient
                            X                                    X
              degx (f ) =         degz (fK ) et degx (f ' ) =         degz (fK' ).
                             z in p-1 (x)                                  z in p-1 (x)

  Soit z  in  p-1 (x). Le point z est interieur a XK . De plus, il appartient a l[RIGHT SINGLE QUOTATION MARK]image
reciproque du squelette canonique par fK et par fK' . Quitte a etendre les scalaires a K,
on peut ainsi supposer que k est algebriquement clos.
  Notons (f1 , . . . , fn ) et (f1' , . . . , fn' ) les composantes de f et de f ' . Soit i  in  {1, . . . , n}.
La fonction numerique |fi' /fi | est constante sur X ; comme x  in  Int(X) et comme le
corps k est maintenant suppose algebriquement clos, la proposition 12.2.12 affirme qu[RIGHT SINGLE QUOTATION MARK]il
existe un element lamdai  in  k  x  et une fonction gi  in  OX,x             x 
                                                                    tels que |gi | < |lamdai | et fi = lamdai + gi
au voisinage de x. Quitte a composer f avec la translation par (lamda1 , . . . , lamdan ) dans Gnm , on
suppose que lamda1 = * * * = lamdan = 1. Alors, trop(y) = trop(y ' ), donc y = y ' puisque ces deux
points appartiennent a S(Gnm ).
  On a dk (x) = n = dim(X) ; l[RIGHT SINGLE QUOTATION MARK]anneau local OX,x est donc artinien ([SECTION SIGN]11.5.5) ; c[RIGHT SINGLE QUOTATION MARK]est
en particulier un anneau de Cohen-Macaulay. Choisissons alors un voisinage affinoide
connexe V de y dans Gnm et et un voisinage affinoide W de x dans Int(X) tels que les
conditions suivantes soient satisfaites :
  - Les fonctions g sont definies sur W et l[RIGHT SINGLE QUOTATION MARK]on a |gi | < 1 et fi' = (1 + gi )fi sur W ;
  - L[RIGHT SINGLE QUOTATION MARK]espace W est de Cohen-Macaulay ;
  - Le point x est le seul antecedent de f (x) par f dans W ;
  - Le morphisme f definit par restriction un morphisme fini et plat de W vers V .
Comme x est le seul antecedent de f (x) par f dans W , on a degf (x) (f ) = degx (f ).
Comme V est connexe, il s[RIGHT SINGLE QUOTATION MARK]ensuit le morphisme fini et plat f |W : W -> V est de degre
constant d.
   Soit R un nombre reel > 1 tel que R|g| < 1 sur W et soit B = D(0, R) le disque
ouvert de rayon R dans la droite affine ; on note theta la fonction coordonnee sur B. Soit h
le morphisme de W  x  B vers Gnm  x  B defini par la famille

                                   ((1 + thetag1 )f1 , . . . , (1 + thetagn )fn ).

Pour tout t  in  B, on note ht le morphisme de WH (t) dans Gnm,H (t) deduit de h par
restriction a la fibre de la derniere projection en t. En particulier, on a h0 = f |W et
h1 = f ' |W .
  Rappelons ([SECTION SIGN]12.1.10) que pour r  in  Rn et t  in  B, on a note sqt (r) l[RIGHT SINGLE QUOTATION MARK]unique point
de Gnm  x k B caracterise par les proprietes suivantes :

404          CHAPITRE 14. CALIBRAGE CANONIQUE D[RIGHT SINGLE QUOTATION MARK]UN ESPACE ANALYTIQUE

  - son image par la composition de la premiere projection et de l[RIGHT SINGLE QUOTATION MARK]application de tropi-
calisation est r ;
  - son image par la seconde projection est t ;
  - il appartient au squelette canonique du H (t)-tore Gnm,H (t) , immerge dans Gnm  x  B.
\end{lstlisting}

\subsubsection{CLD.txt: cached text lines
21545--21664}\label{cld.txt-cached-text-lines-2154521664}

\begin{lstlisting}
de x ; en particulier, l[RIGHT SINGLE QUOTATION MARK]anneau local OX,x est un corps. Le theoreme 6.3.7 de [52] entraine
alors que le morphisme d[RIGHT SINGLE QUOTATION MARK]anneaux locaux OX,x -> OXy ,x est regulier. En consequence, la

414           CHAPITRE 14. CALIBRAGE CANONIQUE D[RIGHT SINGLE QUOTATION MARK]UN ESPACE ANALYTIQUE

fibre h-1 (y) est reduite en x. Le degre du morphisme h en x est alors egal au degre de
l[RIGHT SINGLE QUOTATION MARK]extension H (x)/H (y), c[RIGHT SINGLE QUOTATION MARK]est-a-dire P .
   Soit [OHM SIGN] un voisinage affinoide de x dans X tel que le polyedre calibre pif |[OHM SIGN] soit un
voisinage de 0, etoile en 0, dans le cone calibre pif,x ; on exige en outre que x soit le seul
antecedent de eta1 dans [OHM SIGN]. Le morphisme g|[OHM SIGN] est alors fini et plat au-dessus de eta1 . Soit
W un voisinage affinoide de eta1 dans Gm tel que g induise un morphisme fini et plat
surjectif de [OHM SIGN]  intersect  g -1 (W ) dans W . Il existe une cellule C de Rn , voisinage de l[RIGHT SINGLE QUOTATION MARK]origine,
telle que Csq  subset  W . Soit W ' = trop-1 (C) ; c[RIGHT SINGLE QUOTATION MARK]est un domaine affinoide de Gnm . Alors
gtrop ([OHM SIGN]  intersect  g -1 (W  intersect  W ' )) = C et g([OHM SIGN]  intersect  g -1 (W  intersect  W ' )) contient Csq . Quitte a remplacer [OHM SIGN]
par [OHM SIGN]  intersect  g -1 (W  intersect  W ' ), on peut donc supposer que gtrop ([OHM SIGN]) est une cellule, voisinage de
l[RIGHT SINGLE QUOTATION MARK]origine dans Rn , et que g|[OHM SIGN] est fini et plat au-dessus de gtrop ([OHM SIGN])sq .
   Soit C une decomposition cellulaire du cone tropical pif,x , dont les cellules sont des
cones et qui est adaptee a son calibrage ; on suppose en outre que son image par qtrop se
prolonge en une decomposition cellulaire de Rn .
   Soit C une n-cellule ouverte de pif,x ; posons C ' = gtrop (C). Notons aussi D = C  intersect 
                                                                             -1
ftrop ([OHM SIGN]) et D' = qtrop (D). L[RIGHT SINGLE QUOTATION MARK]ensemble D est ouvert et ferme dans qtrop         (D' )  intersect  ftrop ([OHM SIGN]),
         -1                                         -1
donc ftrop  (D)  intersect  [OHM SIGN] est ouvert et ferme dans gtrop      (D' )  intersect  [OHM SIGN]. Par suite, la restriction de g
    -1                                                '             '
a ftrop (D)  intersect  [OHM SIGN] est finie et plate au-dessus de Dsq     . Comme Dsq    est connexe et non vide,
il existe un unique entier NC,[OHM SIGN] tel que cette restriction soit de degre NC,[OHM SIGN] . Le vecteur-
volume [MICRO SIGN]C correspondant au calibrage de la cellule C de pif,x est egal a NC,[OHM SIGN] sigmaC (e), ou
e est le vecteur-volume canonique de Rn et sigmaC : Rn -> Rm est l[RIGHT SINGLE QUOTATION MARK]unique section de qtrop
d[RIGHT SINGLE QUOTATION MARK]image <C>. Cela prouve que l[RIGHT SINGLE QUOTATION MARK]entier NC,[OHM SIGN] ne depend pas du choix du voisinage [OHM SIGN] choisi
comme ci-dessus ; on le notera NC .
   Le cone tropical pif ,xi jouit d[RIGHT SINGLE QUOTATION MARK]une description similaire, a l[RIGHT SINGLE QUOTATION MARK]aide d[RIGHT SINGLE QUOTATION MARK]un voisinage affi-
noide upsilon de xi dans V an et de decompositions cellulaires de f (upsilon), de pif ,xi et de Rn adap-
tees a qtrop . Par raffinement, on suppose que ces decompositions cellulaires coincident
avec celles introduites precedemment pour pif,x .
   Pour toute n-cellule C de pif ,xi qui rencontre la n-cellule ouverte D de C , d[RIGHT SINGLE QUOTATION MARK]image D'
                                  -1                                              '
dans Rn , la restriction de g a ftrop (C)  intersect  upsilon est finie et plate au-dessus de Dsq  ; on note NC
le degre correspondant, de sorte que le vecteur-volume [MICRO SIGN]C donne par le calibrage de pif ,xi
est egal a NC sigmaC (e).
   D[RIGHT SINGLE QUOTATION MARK]apres la proposition suivante, on a NC = NC pour toute n-cellule C de pif,x . Le
theoreme s[RIGHT SINGLE QUOTATION MARK]ensuit.

Proposition 14.6.4. -- Pour toute n-cellule C de pif,x , on a NC = NC .

Demonstration. -- On peut definir la n-cellule       conique C comme le lieu de validite
                                                                                               '
dans Rm d[RIGHT SINGLE QUOTATION MARK]une conjonction d[RIGHT SINGLE QUOTATION MARK]inegalites phi = { P aV
                                                V P
                                                          r,i xP
                                                               i < 0}. De meme, le cone C
de Rn est defini par une conjonction d[RIGHT SINGLE QUOTATION MARK]inegalites phi' = { a's,j yj < 0}.
  On note psi la conjonction des inegalites monomiales |fi |P ar,i < 1 ; elle definit un
                                                              Q
                             e la conjonction des inegalites monomiales Q|f]       P ar,i
ouvert de X. On note aussi psi                                                 i (x)|       < 1;
                 ]
en interpretant fi (x) comme une fonction rationnelle sur V , cette conjection definit un
ouvert de V an .
                                                                '
  De meme, notons psi' la conjonction des inegalites |Tj |as,j < 1, definissant un ouvert
                                                       Q

de Gnm,k . On notera psie' la meme conjonction d[RIGHT SINGLE QUOTATION MARK]inegalites definissant un ouvert de Gn .
                                                                                          m,k

                                    [SECTION SIGN] 14.6. CONES TROPICAUX                                  415

  Le degre NC peut etre caracterise ainsi : dans l[RIGHT SINGLE QUOTATION MARK]ouvert de [OHM SIGN] defini par la condition psi,
tout point de S(Gnm ) satisfaisant psi' et suffisamment proche de eta1 a NC antecedents
(comptes avec degres de l[RIGHT SINGLE QUOTATION MARK]extension de corps residuels et multiplicite) par le morphisme g.
                                                                                            P
  Pour tout i, choisissons une fonction hi  in  OY,y telle que |hi (y)| = 1 et h]        ]
                                                                              i (y) = fi (x) ;
                                P
c[RIGHT SINGLE QUOTATION MARK]est possible puisque H   ^  (x)  subset  F = H  ^ (y) et puisque, Y = Y an etant sans bord
(donc presque bon en tout point), kappa(y) est dense dans H (y). Au voisinage de x, les
              P
fonctions fQ  i et hi [WHITE BULLET] h ont donc meme valeur absolue. Notons theta la conjonction des
inegalites |hi |ar,i < 1. Elle definit un ouvert de Y dont l[RIGHT SINGLE QUOTATION MARK]image reciproque            par h se
decrit au voisinage, de x par la conjonction des inegalites monomiales |fi |P ar,i < 1,
                                                                                  Q
c[RIGHT SINGLE QUOTATION MARK]est-a-dire par psi.
   L[RIGHT SINGLE QUOTATION MARK]extension F definit un revetement fini etale Y f de (Gn ) au voisinage de son point
                                                                m k
generique et le morphisme g : (V , xi) -> (Gm )k se factorise par un morphisme p : Y
                                    an        n                                           g an ->

(Gnm )k . Soit zeta = p(xi) ; alors le morphisme h : (V an , xi) -> (Y  fan , zeta) est fini et plat de
degre P . De plus, le morphisme p est fini et etale en zeta. La condition psi          e ne portant,
par construction, que sur des normes de fonctions sur Y        f, elle definit une condition theta  e
sur Yf .an



Lemme 14.6.5. -- a) Par le morphisme g, tout point de S(Gnm ) satisfaisant psi' et
suffisamment proche de eta1 a exactement NC /P antecedents (comptes avec degres de
l[RIGHT SINGLE QUOTATION MARK]extension de corps residuels) dans Y satisfaisant theta.
   b) Par le morphisme g, tout point de S(Gnm,k ) satisfaisant psi
                                                                 e' et suffisamment proche
de eta1,k a NC /P antecedents (comptes avec degres de l[RIGHT SINGLE QUOTATION MARK]extension de corps residuels)
dans Yfan satisfaisant theta.
                       e

Demonstration. -- Nous ne demontrons que la premiere assertion, la seconde en est le
cas particulier d[RIGHT SINGLE QUOTATION MARK]un corps trivialement value.
  On commence par choisir des voisinages affinoides X1 de x dans [OHM SIGN], Y1 de y dans Y et Z1
de eta1 dans Gnm , tels que h induise un morphisme fini et plat de degre P , h1 : X1 -> Y1 ,
et p : Y -> Gmn induise un morphisme fini etale p1 : Y1 -> Gnm . C[RIGHT SINGLE QUOTATION MARK]est possible car le
morphisme h est fini et plat en x et le morphisme p est fini et etale en y.
  Comme x est le seul antecedent de eta1 par g dans [OHM SIGN], le voisinage X1 de x est un voisinage
de la fibre de g ; il contient donc l[RIGHT SINGLE QUOTATION MARK]image inverse d[RIGHT SINGLE QUOTATION MARK]un voisinage de eta1 . Autrement dit,
tous les antecedents dans [OHM SIGN] d[RIGHT SINGLE QUOTATION MARK]un point assez proche de eta1 appartiennent a X1 . De meme,
tous les antecedents dans Y d[RIGHT SINGLE QUOTATION MARK]un point assez proche de eta1 appartiennent a Y1 .
  Soit z un point de Gnm assez proche de eta1 satisfaisant la condition psi' . Notons y1 , . . . , ya
ses antecedents dans Y1 par p. Pour i  in  {1, . . . , a}, notons xij les antecedents de yi
dans X1 par h ; on remarque que xij verifie la condition psi si et seulement si yi verifie la
condition theta. Par suite, la definition de NC entraine que
                      X                        X
              NC =         rgOZ,eta1 (OX,xij ) =   rgOZ,eta1 (OY,yi ) rgOY,y (OX,xij ).
\end{lstlisting}

\subsubsection{CLD.txt: cached text lines
21712--21776}\label{cld.txt-cached-text-lines-2171221776}

\begin{lstlisting}
   Retraduisons ceci en termes de valuations. La valuation etag du corps k(T1 , . . . , Tn )
possede ainsi NC /P extensions a k(Y ) et satisfaisant theta, comptees avec multiplicites. A
priori, la multiplicite associee a une telle extension est le degre de son henselise sur celui
du henselise de (k(T1 , . . . , Tn ), eta1 ) ; par stabilite, cette multiplicite est simplement egale
au produit des indices d[RIGHT SINGLE QUOTATION MARK]inertie et de ramification de l[RIGHT SINGLE QUOTATION MARK]extension.
   On prouve de meme que la valuation etag,k possede N C /P extensions a e          k(Yf) satisfaisant
Teheta, comptes avec multiplicites ; et la multiplicite associee a une telle extension est
simplement egale au produit de ses indices d[RIGHT SINGLE QUOTATION MARK]inertie et de ramification.
   Il resulte alors du lemme suivant que NC /P = N            fC /P , si bien que NC = N  eC , ce qui
conclut la demonstration de la proposition 14.6.4.

Lemme 14.6.6. -- Par composition avec la valuation associee a y, toute valuation v
sur k(Y  ) triviale sur k induit une valuation v sur k(Y ).

                         [SECTION SIGN] 14.7. VARIATION DES TROPICALISATIONS                           417

   a) Cette application v 7-> v induit une bijection de l[RIGHT SINGLE QUOTATION MARK]ensemble des prolongements
a k(Y  ) de etag,k sur l[RIGHT SINGLE QUOTATION MARK]ensemble des prolongements a k(Y ) de etag .
   Soit v et v deux telles valuations.
   b) Les indices d[RIGHT SINGLE QUOTATION MARK]inertie de v et de v sont egaux.
   c) Les indices de ramification de v et de v sont egaux.
   d) La valuation v verifie la condition theta si et seulement si la valuation v verifie la
condition theta.
Demonstration. -- La valuation etag raffine eta1 qui n[RIGHT SINGLE QUOTATION MARK]a qu[RIGHT SINGLE QUOTATION MARK]un seul prolongement a k(Y ).
En effet, comme k(Y ) est finie sur k(T1 , . . . , Tn ), le groupe des valeurs d[RIGHT SINGLE QUOTATION MARK]une valuation
de k(Y ) qui prolonge eta1 contient celui de H (eta1 ) comme sous-groupe d[RIGHT SINGLE QUOTATION MARK]indice fini, donc
coincide avec celui-ci puisqu[RIGHT SINGLE QUOTATION MARK]il est egal a |k  x  | et que k est algebriquement clos. En
particulier, tout prolongement de eta1 a k(Y ) appartient a Y an , donc est egal a y qui
est le seul antecedent de eta1 dans Y an . La premiere assertion resulte alors de ce que la
valuation etag est la composee de la valuation de Gau[LATIN SMALL LETTER SHARP S] usuelle eta1 sur k(T1 , . . . , Tn ) et de
la valuation etag,k sur k(T1 , . . . , Tn ),
   Puisque les extensions residuelles correspondantes coincident, les indices d[RIGHT SINGLE QUOTATION MARK]inertie de v
et de v sont egaux. Puisque |H (y) x  | = |H (eta1 ) x | = |k  x  |, leurs indices de ramification
sont aussi egaux.
   Demontrons maintenant que v verifie la condition theta si et seulement si v la condition theta.
Par construction, theta est une conjonction d[RIGHT SINGLE QUOTATION MARK]inegalites de la forme |phi| < 1, ou phi  in  OY,y
verifie |phi(y)| = 1, et theta est la conjonction des inegalites |phi(y)|    < 1 (pour les memes
fonctions phi) ou phi(y)   in H   ^    (y) est considere comme un element de O   . Rappelons que
                                                                         Y ,zeta
la valuation v est la composee de la valuation associee a y et de v ; comme |phi(y)| = 1,
les conditions |phi(v)| < 1 et |phi(v)| < 1 sont equivalentes. Le lemme en resulte
   Terminons ce paragraphe par un complement sur la structure des cones tropicaux
lorsqu[RIGHT SINGLE QUOTATION MARK]ils ne sont pas de dimension maximale qui indique neanmoins qu[RIGHT SINGLE QUOTATION MARK]elles ne peuvent
etre arbitraire. En particulier, il resultera du theoreme 15.6.1 qu[RIGHT SINGLE QUOTATION MARK]ils peuvent etre munies
d[RIGHT SINGLE QUOTATION MARK]un calibrage partout strictement positif et harmonieux.



                [SECTION SIGN] 14.7. VARIATION DES TROPICALISATIONS

Lemme 14.7.1. -- Soit X un espace analytique compact, soit f : X -> Y un mor-
phisme d[RIGHT SINGLE QUOTATION MARK]espaces analytiques, soit g : X -> Gnm un moment. Pour tout y  in  Y , on pose
Xy = f -1 (y). Soit y un point rigide de Y . Il existe un voisinage U de y dans Y et un
polyedre [INCREMENT]  subset  gtrop (Xy ) de dimension <= dim(Xy ) - 1 tels que gtrop (Xu )  subset  gtrop (Xy ) et
gtrop (Xu  intersect  boundary (X))  subset  [INCREMENT] pour tout u  in  U .
                                               -1
Demonstration. -- Le domaine analytique gtrop      (gtrop (Xy )) de X contient le ferme de
Zariski Xy , donc en est un voisinage (lemme 11.4.10). Comme X est compact, il existe
                                           -1
un voisinage U de y tel que f -1 (U )  subset  gtrop (gtrop (Xy )). Ainsi, si u  in  U , gtrop (Xu ) est
contenu dans gtrop (Xy ).
   Le bord boundary (X) de X est contenu dans une reunion finie de compacts qui sont chacun de
la forme phi-1 (etar ), ou phi est un morphisme d[RIGHT SINGLE QUOTATION MARK]un domaine affinoide V de X vers A1 et r est
un nombre reel > 0 ([50], lemme 3.1). Soit phi : V -> A1 un tel morphisme et soit r > 0.
\end{lstlisting}

\subsubsection{CLD.txt: cached text lines
22100--22151}\label{cld.txt-cached-text-lines-2210022151}

\begin{lstlisting}
pace C (X) des fonctions continues sur X muni de la topologie de la convergence uniforme
sur tout compact.
Demonstration. -- Comme X est localement differentiellement separe, c[RIGHT SINGLE QUOTATION MARK]est une conse-
quence du corollaire 6.2.6.

Proposition 15.2.8 (Partitions de l[RIGHT SINGLE QUOTATION MARK]unite lisses). -- Soit X un espace analytique
localement holomorphiquement separe, paracompact. Soit (Ui )i in I un recouvrement ouvert
de X. Il existe une partition de l[RIGHT SINGLE QUOTATION MARK]unite subordonnee a (Ui ) formee de fonctions lisses
sur X.
Demonstration. -- Comme X est localement differentiellement separe, cela resulte de la
proposition 6.2.7.

Corollaire 15.2.9. -- Soit X un espace analytique localement holomorphiquement se-
pare et soit F une partie fermee de X qui est denombrable a l[RIGHT SINGLE QUOTATION MARK]infini (par exemple, com-
pacte). Soit (Ui )i in I un recouvrement de K par des ouverts de X. Il existe une famille
localement finie (thetai ) de fonctions lisses sur X, a support
                                                          P compact, positives, telle que
pour tout i, le support de thetai soit contenu dans Ui et que   thetai vaille 1 au voisinage de F .
Demonstration. -- Comme X est localement differentiellement separe, cela resulte du
corollaire 6.2.8.

Corollaire 15.2.10. -- Soit X un espace analytique localement holomorphiquement
separe, paracompact. Les faisceaux AXp,q de formes lisses sur X sont fins.

                         [SECTION SIGN] 15.4. INTEGRALES ET MORPHISMES FINIS                        425

Demonstration. -- Comme X est localement differentiellement separe, cela resulte du
corollaire 6.2.9.




                     [SECTION SIGN] 15.3. EXTENSION DES SCALAIRES

  Soit X un espace k-analytique de niveau n. On le considere comme un espace G-
tropical de niveau n. On a montre au [SECTION SIGN]5.8.5 comment toute (n, n)-forme omega sur XG a
coefficients tropicalement boreliens definit une mesure de Borel omegaBor sur X, et lorsque
cette G-forme est localement integrable, une mesure de Radon omegaRad .

Proposition 15.3.1. -- Soit L une extension complete de k, soit p : XL -> X le mor-
phisme de changement de base. Pour toute forme omega de type (n, n) sur XG a coefficients
tropicalement boreliens, on a p* (omegaL,Bor ) = omegaBor . Lorsque omega est localement integrable, on
a p* (omegaL,Rad ) = omegaRad .

Demonstration. -- En revenant aux definitions, on deduit de la compatibilite du cali-
brage canonique a l[RIGHT SINGLE QUOTATION MARK]extension des scalaires (proposition 14.3.6) l[RIGHT SINGLE QUOTATION MARK]egalite p* omega
                                                                              fL = omegaL de
formes-volume n-dimensionnelles. La proposition en resulte.

Proposition 15.3.2. -- Soit L une extension complete de k, soit p : XL -> X le mor-
phisme de changement de base. Pour toute forme omega de type (n-1, n) sur XG a coefficients
tropicalement boreliens, on a p* (omegaL,Bor ) = omegaBor . Lorsque omega est localement integrable, on
\end{lstlisting}

\subsubsection{CLD.txt: cached text lines
29080--29148}\label{cld.txt-cached-text-lines-2908029148}

\begin{lstlisting}

Remarque 19.1.9. -- Placons-nous encore sous les hypotheses de la proposition. No-
tons L le fibre en droites L muni de la metrique formelle deduite du prolongement L.
Soit g = dim(A) et soit gamma  in  A l[RIGHT SINGLE QUOTATION MARK]unique point dont la reduction est le point generique
de As .
  a) La mesure c1 (L)g est supportee par gamma. Lorsque k est algebriquement clos, le theo-
reme 18.4.7 assure que sa masse est (c1 (Ls )g  intersect  [As ]) ; c[RIGHT SINGLE QUOTATION MARK]est encore le cas dans le cas
general, comme on le voit en appliquant le resultat a l[RIGHT SINGLE QUOTATION MARK]extension du schema formel A
   [WHITE BULLET]
ak .
  Lorsque A est algebrique, cette masse est egalement egale a (c1 (L)g  intersect  [A]) (theo-
reme 18.3.11).
  b) Il resulte de l[RIGHT SINGLE QUOTATION MARK]exemple 18.8.10, b), que les trois proprietes de positivite suivantes
sont equivalentes :
  (i) Le courant c1 (L) est positif ;
  (ii) Le courant c1 (L) est positif au voisinage de gamma ;
  (iii) Le fibre en droites Ls est numeriquement effectif.

                     [SECTION SIGN] 19.2. METRIQUES CUBISTES ET UNIFORMISATION                         551

L[RIGHT SINGLE QUOTATION MARK]implication (i)=>(ii) est evidente ; l[RIGHT SINGLE QUOTATION MARK]implication (ii)=>(iii) est expliquee dans l[RIGHT SINGLE QUOTATION MARK]exemple
indique, de meme que l[RIGHT SINGLE QUOTATION MARK]implication (iii)=>(i).

   Dans le paragraphe suivant, nous allons exploiter l[RIGHT SINGLE QUOTATION MARK]uniformisation des varietes abe-
loides pour obtenir une description similaire des metriques cubistes sur les fibres en
droites et des puissances de leurs courants de courbure.



          [SECTION SIGN] 19.2. METRIQUES CUBISTES ET UNIFORMISATION

19.2.1. -- Rappelons l[RIGHT SINGLE QUOTATION MARK]uniformisation des varietes abeloides, d[RIGHT SINGLE QUOTATION MARK]apres [8, 23, 78, 98] ;
nous renvoyons aussi le lecteur a l[RIGHT SINGLE QUOTATION MARK]appendice B.5.
   Soit A une variete abeloide sur k ; notons g sa dimension. Il existe une variete abe-
loide B a bonne reduction potentielle, un tore algebrique T de dimension r, une exten-
sion E de B par T , une suite exacte de groupes analytiques
                                                  p
                                  0 -> lamdaan -> E ->
                                              - A->0
ou lamda est un k-groupe localement isomorphe a Zr pour la topologie etale de Spec(k).
   Nous dirons que la situation est deployee si B a bonne reduction et si le tore T est
deploye. Dans ce cas, lamda est isomorphe a Zr , la suite exacte precedente s[RIGHT SINGLE QUOTATION MARK]interprete comme
le revetement universel de l[RIGHT SINGLE QUOTATION MARK]espace topologique Aan et lamda comme son groupe fondamental.
   Nous nous placons desormais dans ce cas.
19.2.2. -- Notons q : E -> B la projection de l[RIGHT SINGLE QUOTATION MARK]extension E. Son noyau est le tore
deploye T , de sorte que cette extension est localement triviale pour la topologie de
Berkovich.
   L[RIGHT SINGLE QUOTATION MARK]extension E est construite a partir d[RIGHT SINGLE QUOTATION MARK]une extension analytique N de B par le tore af-
finoide T 1 ; cet espace N est un domaine analytique compact de A, muni d[RIGHT SINGLE QUOTATION MARK]une structure
de sous-groupe formellement analytique compact.
   Soit V un recouvrement formellement affinoide de B trivialisant l[RIGHT SINGLE QUOTATION MARK]extension N ; pour
tout V  in  V, soit fV : V -> N une section de la projection q|N : N -> B an . Sur V  intersect  W , les
sections fV et fW different de la multiplication par un morphisme fV W : V  intersect W -> T 1 . Ces
morphismes (fV W ) forment un cocycle a valeurs dans T 1 qui represente l[RIGHT SINGLE QUOTATION MARK]extension N ,
mais egalement, si on les considere comme un cocycle a valeurs dans T , l[RIGHT SINGLE QUOTATION MARK]extension E.
   Pour V  in  V, la section fV permet d[RIGHT SINGLE QUOTATION MARK]identifier q -1 (V ) au produit V  x  T ; composons la
projection sur le second facteur avec l[RIGHT SINGLE QUOTATION MARK]application de tropicalisation de T vers Ttrop ; on
obtient une application phiV de q -1 (V ) dans Ttrop qui est lisse et pl : lorsqu[RIGHT SINGLE QUOTATION MARK]on identifie T
a Grm et Ttrop a Rr , chaque coordonnee est le logarithme de la valeur absolue d[RIGHT SINGLE QUOTATION MARK]une
fonction holomorphe inversible (un caractere du tore).
   Comme le cocycle (fV W ) est a valeurs dans T 1 , ces applications phiV se recollent en une
application G-lisse, pl, de E dans Ttrop . Notons-la tau .
   La restriction de tau au tore T est l[RIGHT SINGLE QUOTATION MARK]application de tropicalisation usuelle de T sur Ttrop .
Comme tau |T est propre et que B est compacte, l[RIGHT SINGLE QUOTATION MARK]application tau : E -> Ttrop est surjective
et propre.
   Si E = B, le tore T est de dimension 0 et tau est l[RIGHT SINGLE QUOTATION MARK]application constante.
   L[RIGHT SINGLE QUOTATION MARK]image du reseau lamda par tau est un reseau cocompact de Ttrop (theoreme 1.2 de [24]).
Par suite, tau induit une application surjective, pl, de A dans le tore reel Ttrop /tau (lamda).
\end{lstlisting}

\subsubsection{GR.txt: cached text lines
302--368}\label{gr.txt-cached-text-lines-302368}

\begin{lstlisting}
pi-1
  k r /k
         (x) = M (H (x) tensor b k r ) = M (H (x) r ) (which is the weighted Gauss point eta r ). We
call sigma the Shilov section. The Shilov section is continuous: see [Ber90, Lemma 3.2.2(i)],
[Duc18, 1.2.16].

2.11. Relative boundary. An important notion in this paper is that of the relative bound-
ary of a morphism greek phi symbol : X -> Y of k-analytic spaces. This is a closed subset of X which we
denote by boundary  (X /Y ). We will not give the details of this definition, but let us try to explain
the literature on the subject. In [Ber90, [SECTION SIGN]2.5], Berkovich defines the relative boundary
boundary  (X /Y ) for k-affinoid spaces X and Y . Then in [Ber90, [SECTION SIGN]3.1], he generalizes boundary  (X /Y ) for
good k-analytic spaces using that all points of X and Y have k-affinoid neighborhoods.
Berkovich calls greek phi symbol closed if boundary  (X /Y ) = ;. Berkovich defines boundary  (X /Y ) for arbitrary k-analytic
spaces in [Ber93, 1.5.3(ii)], as follows. A morphism greek phi symbol is closed if for every morphism
Y ' -> Y from a good k-analytic space Y ' , the fiber product X ' := X  x Y Y ' is good and the
canonical morphism X ' -> Y ' is closed in the previous sense. Finally in [Ber93, Defini-
tion 1.5.4], a point x  in  X is defined to be in Int(X /Y ) if x has an open neighborhood U
such that greek phi symbol induces a closed morphism U -> Y in the above sense. We refer to [Tem15,
Definition 4.2.4.1] for a more explicit equivalent definition of boundary  (X /Y ). Note that in this
paper, we will use the notion closed map in the usual topological sense meaning that it
maps closed sets to closed sets. This is different from Berkovich[RIGHT SINGLE QUOTATION MARK]s terminology.
   The complement Int(X /Y ) := X \ boundary  (X /Y ) is called the relative interior. We will usually
apply this in the case Y = M (k), in which case we call boundary  X := boundary  (X /M (k)) the boundary
of X and Int(X ) := X \ boundary  X the interior. If boundary  X = ;, we say that X boundaryless. A bound-
aryless space is good by definition. Note that Berkovich calls boundaryless k-analytic
spaces closed.
   We list here some properties of the boundary and interior which we will use throughout
the paper:
   (1) A point x  in  X is contained in Int(X /Y ) if and only if there is an open neighbor-
       hood U of x such that Int(U/Y ) = U: this is part of the above definition.
   (2) If greek phi symbol : X -> Y is finite, then Int(X /Y ) = X [Ber93, Example 1.5.3(iii)].
   (3) If X is an analytic domain in Y , then Int(X /Y ) is the topological interior of X in
       Y [Ber93, Proposition 1.5.5(i)].
             greek phi symbol     psi
   (4) If X -> Y -> S are morphisms of k-analytic spaces, then

(2.11.1)                   Int(X /Y )  intersect  greek phi symbol -1 Int(Y /S)  subset  Int(X /S)

        [Ber93, Proposition 1.5.5(ii)], with equality if psi is locally separated [Tem04,
        Corollary 5.7]. Taking S = M (k), this becomes

(2.11.2)                      Int(X /Y )  intersect  greek phi symbol -1 Int(Y )  subset  Int(X ).

   (5) If k' /k is an analytic extension field, then

                                pi-1
                                 k' /k
                                       Int(X /Y )  subset  Int(X k' /Yk' )

8                                      W. GUBLER AND J. RABINOFF

         [Ber93, Proposition 1.5.5(iii)]. Taking Y = M (k), this becomes
                                       pi-1
                                        k' /k
                                              Int(X )  subset  Int(X k' ).
For convenience, we restate these properties in terms of the relative boundary:
   (1' ) A point x  in  X is contained in boundary  (X /Y ) if and only if boundary  (U/Y ) /= ; for every open
         neighborhood U of x.
   (2' ) If greek phi symbol is finite, then boundary  (X /Y ) = ;.
   (3' ) If X is an analytic domain in Y , then boundary  (X /Y ) is the topological boundary of X
         in Y .
               greek phi symbol     psi
     '
   (4 ) If X -> Y -> S are morphisms of k-analytic spaces, then
(2.11.3)                        boundary  (X /S)  subset  boundary  (X /Y )  union  greek phi symbol -1 (boundary  (Y /S)),
         with equality if psi is locally separated. Taking S = M (k), this becomes
(2.11.4)                             boundary  X  subset  boundary  (X /Y )  union  greek phi symbol -1 (boundary  Y ).
\end{lstlisting}


\end{document}
